\chapter{The tautological quotient and the universal property of $\mathrm{Gr}(r,d)$}
\label{chap:Picard_GrassmannianQuot}

\providecommand{\ShMod}{\mathrm{SheafOfModules}}

The previous chapter constructed the Grassmannian scheme \(\mathrm{Gr}(r,d)\) over
\(\mathbb{Z}\) by gluing the affine charts \(U^I\)
(\cref{def:gr_affine_chart}) along the transition isomorphisms
(\cref{def:gr_transition}) satisfying the cocycle condition
(\cref{lem:gr_cocycle}), and proved that the structure morphism
\(\mathrm{Gr}(r,d) \to \Spec \mathbb{Z}\) is separated (\cref{lem:gr_separated})
and proper (\cref{lem:gr_proper}). This chapter adds the next layer: the
\emph{tautological rank-\(d\) quotient}
\(u : \mathcal{O}_{\mathrm{Gr}(r,d)}^{\,r} \twoheadrightarrow \mathcal{U}\) of the
trivial rank-\(r\) sheaf, and the \emph{universal property} that makes
\(\mathrm{Gr}(r,d)\) the fine moduli space of rank-\(d\) locally free quotients of
\(\mathcal{O}^{r}\) -- equivalently, the representing scheme of the Quot functor
of the trivial sheaf.

Throughout, \(r \ge d \ge 1\) are the fixed integers of the previous chapter, and
for a size-\(d\) subset \(I \subseteq \{1,\dots,r\}\) we write \(R^I\) for the
chart ring \(\mathbb{Z}[X^I]\) of \(U^I = \Spec R^I\) and \(R^I_J\) for the
localisation along \(P^I_J = \det(X^I_J)\) defining the overlap
\(U^I_J \subseteq U^I\).

\section{Infrastructure: locally free sheaves and module gluing}
\label{sec:grquot_infra}

The constructions of this chapter -- the universal quotient sheaf
\(\mathcal{U}\), the tautological quotient \(u\), and the functor of points -- rest
on two pieces of foundational machinery for sheaves of modules. The first is the
predicate ``locally free of rank \(d\)'' (\cref{def:is_locally_free_of_rank}),
which records that a sheaf of modules is, on the members of some open cover,
isomorphic to the trivial bundle \(\mathcal{O}^{d}\); it is the condition cut out
in the definition of the Grassmannian functor. That predicate is the project-local
\(\mathrm{SheafOfModules.IsLocallyFreeOfRank}\) already set up in the Quot-scheme
chapter (\cref{def:is_locally_free_of_rank}), which this chapter reuses verbatim;
the only foundational machinery genuinely new to this chapter is the second piece:
a gluing primitive that descends a sheaf of modules (and a morphism of such
sheaves) from per-chart data along an open cover, given a transition cocycle
compatible with the cover's own gluing; it is the vehicle that assembles the
per-chart free sheaves \(\mathcal{O}^{d}_{U^I}\) and the per-chart quotients
\(u^I\) into the global \(\mathcal{U}\) and \(u\), and is set up below.

Stating the triple-overlap multiplicativity of the descent datum, however, first
requires a small piece of base-change machinery for the pullback functor on sheaves
of modules: the three transition isomorphisms attached to a triple of charts live a
priori on three different overlaps, and to compare their composite they must be
transported to a single common overlap. The next two blocks supply exactly this
transport -- the pseudofunctorial comparison of iterated pullbacks
(\cref{def:modules_pullbackComp}) and its packaging along the triple-overlap data of a
scheme glue datum (\cref{lem:modules_pullback_basechange_transport}) -- after which the
cocycle condition can be expressed cleanly.

\begin{definition}[Comparison of iterated pullbacks of sheaves of modules]
  \label{def:modules_pullbackComp}
  \lean{AlgebraicGeometry.Scheme.Modules.pullbackComp}
  \mathlibok
  \textit{Provided by Mathlib.}
  Let \(a : X \to Y\) and \(b : Y \to Z\) be morphisms of schemes. For a sheaf of
  \(\mathcal{O}_Z\)-modules \(\mathcal{M}\), pulling back first along \(b\) and then
  along \(a\) agrees with pulling back along the composite \(b \circ a\): there is a
  natural isomorphism of sheaves of \(\mathcal{O}_X\)-modules
  \[
    a^{*}\bigl(b^{*}\mathcal{M}\bigr)
      \;\xrightarrow{\ \sim\ }\;
    (b \circ a)^{*}\mathcal{M},
  \]
  natural in \(\mathcal{M}\), expressing that \(\mathcal{M} \mapsto \mathcal{M}\)'s
  pullback is a pseudofunctor on schemes. These comparison isomorphisms satisfy the
  pentagon (associativity) coherence for a triple of composable morphisms. This is the
  pullback half of the pseudofunctoriality of sheaves of modules and is supplied by
  Mathlib.
\end{definition}

\begin{lemma}\leanok
[Base-change transport of a transition isomorphism to a triple overlap]
  \label{lem:modules_pullback_basechange_transport}
  \lean{AlgebraicGeometry.Scheme.Modules.pullbackBaseChangeTransport}
  \uses{def:modules_pullbackComp, def:gr_the_glue_data}
  Let \(D\) be a glue datum of schemes (\cref{def:gr_the_glue_data}) with charts
  \(U_i\), overlaps \(V_{ij}\), overlap immersions \(f_{ij} : V_{ij} \hookrightarrow
  U_i\), transitions \(t_{ij} : V_{ij} \xrightarrow{\sim} V_{ji}\), and triple-overlap
  data \(t'_{ijk}\) with its compatibility identity \(t_{\mathrm{fac}}\). For an index
  triple \((i,j,k)\) write
  \[
    V_{ijk} \;=\; V_{ij} \times_{U_i} V_{ik}
  \]
  for the common triple overlap, with the two projections
  \(p^{ij}_{ijk} : V_{ijk} \to V_{ij}\) and \(p^{ik}_{ijk} : V_{ijk} \to V_{ik}\), and
  let \(t'_{ijk} : V_{ijk} \xrightarrow{\sim} V_{jki}\) be the glue datum's triple
  transition, satisfying
  \[
    t'_{ijk} \,\circ\, p^{jk}_{jki} \;=\; p^{ij}_{ijk} \,\circ\, t_{ij}
    \qquad (t_{\mathrm{fac}}).
  \]
  Given per-chart sheaves of modules \(\mathcal{M}_i\) and a transition isomorphism
  \[
    g_{ij} \;:\; f_{ij}^{*}\,\mathcal{M}_i
      \;\xrightarrow{\ \sim\ }\;
    (t_{ij} \circ f_{ji})^{*}\,\mathcal{M}_j
  \]
  on the overlap \(V_{ij}\) (as in \cref{def:scheme_modules_glue}), the
  \emph{base-change transport} of \(g_{ij}\) to the triple overlap \(V_{ijk}\) is the
  isomorphism of sheaves of \(\mathcal{O}_{V_{ijk}}\)-modules
  \[
    \widehat{g}_{ij}^{\,k}
      \;:\;
    (f_{ij} \circ p^{ij}_{ijk})^{*}\,\mathcal{M}_i
      \;\xrightarrow{\ \sim\ }\;
    (t_{ij} \circ f_{ji} \circ p^{ij}_{ijk})^{*}\,\mathcal{M}_j
  \]
  obtained by pulling \(g_{ij}\) back along \(p^{ij}_{ijk}\) and reassociating the
  iterated pullbacks via \cref{def:modules_pullbackComp}. With the analogous transports
  \(\widehat{g}_{jk}^{\,i}\) and \(\widehat{g}_{ik}^{\,j}\) (the first taken along
  \(t'_{ijk}\), so that \(t_{\mathrm{fac}}\) identifies its source and target sheaves
  with those of \(\widehat{g}_{ij}^{\,k}\) and \(\widehat{g}_{ik}^{\,j}\)), all three
  transported isomorphisms become endomorphism-composable: \(\widehat{g}_{ij}^{\,k}\)
  and \(\widehat{g}_{jk}^{\,i}\) land on a common sheaf, and their composite is
  comparable to \(\widehat{g}_{ik}^{\,j}\) on \(V_{ijk}\).
\end{lemma}

\begin{proof}
  \uses{def:modules_pullbackComp, def:gr_the_glue_data}
  \leanok
  The pullback functor on sheaves of modules is a pseudofunctor: for composable
  scheme morphisms it carries the comparison isomorphisms of
  \cref{def:modules_pullbackComp}, natural in the sheaf and coherent for triple
  composites. Pulling the transition isomorphism \(g_{ij}\) back along the projection
  \(p^{ij}_{ijk} : V_{ijk} \to V_{ij}\) yields an isomorphism between
  \((p^{ij}_{ijk})^{*} f_{ij}^{*}\mathcal{M}_i\) and \((p^{ij}_{ijk})^{*}(t_{ij} \circ
  f_{ji})^{*}\mathcal{M}_j\); applying \cref{def:modules_pullbackComp} to each side
  reassociates these into the single-pullback form
  \((f_{ij} \circ p^{ij}_{ijk})^{*}\mathcal{M}_i\) and
  \((t_{ij} \circ f_{ji} \circ p^{ij}_{ijk})^{*}\mathcal{M}_j\), giving
  \(\widehat{g}_{ij}^{\,k}\). For the \((j,k)\)-transition the same construction is
  applied along the triple transition \(t'_{ijk}\) rather than a bare projection; the
  glue datum's compatibility identity \(t_{\mathrm{fac}}\),
  \(t'_{ijk} \circ p^{jk}_{jki} = p^{ij}_{ijk} \circ t_{ij}\), guarantees that the
  source of \(\widehat{g}_{jk}^{\,i}\) coincides with the target of
  \(\widehat{g}_{ij}^{\,k}\) (both are the pullback of \(\mathcal{M}_j\) along the same
  morphism \(V_{ijk} \to U_j\)), so the two are composable as endomorphism-style
  isomorphisms over \(V_{ijk}\). The reassociations introduced by
  \cref{def:modules_pullbackComp} are coherent (they satisfy the pentagon identity), so
  the composite \(\widehat{g}_{jk}^{\,i} \circ \widehat{g}_{ij}^{\,k}\) and the third
  transport \(\widehat{g}_{ik}^{\,j}\) are isomorphisms between the same pair of
  sheaves on \(V_{ijk}\), and may be compared directly. This is the module-level
  analogue of the scheme-level triangle \(t'_{ijk}\)/\(t_{\mathrm{fac}}\) and is the
  shape in which the multiplicative cocycle condition is stated below.
\end{proof}

\begin{lemma}[Triple-overlap endpoint bridges]
  \label{lem:gr_glueData_bridges}
  \lean{AlgebraicGeometry.Scheme.Modules.glueData_bridge_src,
        AlgebraicGeometry.Scheme.Modules.glueData_bridge_mid,
        AlgebraicGeometry.Scheme.Modules.glueData_bridge_tgt}
  \uses{def:gr_the_glue_data}
  Let \(D\) be a glue datum of schemes (\cref{def:gr_the_glue_data}) and fix
  indices \(i,j,k\) with triple overlap \(V_{ijk} = V_{ij} \times_{U_i} V_{ik}\).
  The three transports of \cref{lem:modules_pullback_basechange_transport} have
  endpoints pinned down by three scheme-morphism equalities on \(V_{ijk}\), each a
  pure consequence of the glue datum's structural identities:
  \begin{itemize}
    \item (\emph{src}) \(\mathrm{fst} \circ f_{ij} = \mathrm{snd} \circ f_{ik}\)
      -- the defining pullback condition of \(V_{ijk}\);
    \item (\emph{mid}) \(\mathrm{fst} \circ (t_{ij} \circ f_{ji})
      = (t'_{ijk} \circ \mathrm{fst}) \circ f_{jk}\) -- the pullback condition
      composed with the glue datum's \(t_{\mathrm{fac}}\) identity;
    \item (\emph{tgt}) \((t'_{ijk} \circ \mathrm{fst}) \circ (t_{jk} \circ f_{kj})
      = \mathrm{snd} \circ (t_{ik} \circ f_{ki})\) -- the cocycle core, from the
      \(t_{\mathrm{fac}}\) associativity, the reversed pullback condition, the
      \(t_{\mathrm{inv}}\) collapse and the glue datum's cocycle identity.
  \end{itemize}
  These three equalities align the source, middle, and target sheaves of the three
  transports so that the cocycle equation (C2) of \cref{def:scheme_modules_glue}
  typechecks.
\end{lemma}
\begin{proof}  \uses{def:gr_the_glue_data}
  Each equality is an identity of morphisms in the category of schemes, proved by
  rewriting with the glue datum's axioms: \emph{src} is literally
  \(\mathrm{pullback.condition}\); \emph{mid} adds the \(t_{\mathrm{fac}}\)
  factorisation \(D.t_{\mathrm{fac}}\,i\,j\,k\); \emph{tgt} chains
  \(t_{\mathrm{fac}}\)-associativity at \((j,k,i)\), the reversed pullback
  condition, an auxiliary collapse combining \(t_{\mathrm{fac}}\)-associativity at
  \((k,i,j)\) with \(t_{\mathrm{inv}}\), and finally the cocycle identity
  \(D.\mathrm{cocycle}\,i\,j\,k\). No sheaf-of-modules data enters; these are
  bookkeeping equalities of the underlying scheme cover.
\end{proof}

% HISTORICAL (abandoned route): an earlier draft built the glued sheaf as a hand-rolled
% presheaf of compatible families (decls gluePresheaf / gluePresheafModule / gluePresheafIsSheaf).
% That route is abandoned; the gluing is now the equalizer of pushforwards described in the
% Construction paragraph of \cref{def:scheme_modules_glue}, which subsumes it.

\begin{definition}

\leanok
[Gluing a sheaf of modules along an open cover]
  \label{def:scheme_modules_glue}
  \lean{AlgebraicGeometry.Scheme.Modules.glue}
  \uses{def:gr_the_glue_data, def:modules_pullbackComp,
    lem:modules_pullback_basechange_transport, lem:gr_glueData_bridges}
  Let \(D\) be a glue datum of schemes (\cref{def:gr_the_glue_data}), with index set \(J\), charts \(U_i\),
  pairwise overlaps \(V_{ij}\), open immersions
  \(f_{ij} : V_{ij} \hookrightarrow U_i\), and transition isomorphisms
  \(t_{ij} : V_{ij} \xrightarrow{\sim} V_{ji}\) satisfying the glue datum's
  cocycle conditions; write \(X\) for the glued scheme and
  \(\iota_i : U_i \hookrightarrow X\) for the canonical open immersions of the
  charts. Suppose given:
  \begin{itemize}
    \item for each \(i \in J\), a sheaf of \(\mathcal{O}_{U_i}\)-modules \(\mathcal{M}_i\) on the chart \(U_i\);
    \item for each pair \((i,j)\), a \emph{transition isomorphism} of sheaves of
      modules on the overlap \(V_{ij}\),
      \[
        g_{ij} \;:\; f_{ij}^{*}\,\mathcal{M}_i
          \;\xrightarrow{\ \sim\ }\;
        (t_{ij} \circ f_{ji})^{*}\,\mathcal{M}_j,
      \]
      comparing the two charts' sheaves after restriction (pullback) to the
      common overlap \(V_{ij}\);
  \end{itemize}
  \emph{subject to the following module cocycle conditions}, compatible with the
  scheme-level cocycle of \(D\) (its triple transition data \(t_{ij}\) and the
  overlap immersions \(f_{ij}\)), which are required as hypotheses on the family
  \((g_{ij})\):
  \begin{enumerate}
    \item[(C1)] \emph{self-identity:} the self-transition is the identity,
      \(g_{ii} = \mathrm{id}\) on \(f_{ii}^{*}\mathcal{M}_i\) (with
      \(V_{ii}\) and \(t_{ii}\) the diagonal data of \(D\));
    \item[(C2)] \emph{triple-overlap multiplicativity:} over each common triple
      overlap \(V_{ijk} = V_{ij} \times_{U_i} V_{ik}\) the base-change transports
      (\cref{lem:modules_pullback_basechange_transport}) of the three transition
      isomorphisms satisfy
      \[
        \widehat{g}_{jk}^{\,i} \circ \widehat{g}_{ij}^{\,k}
          \;=\; \widehat{g}_{ik}^{\,j},
      \]
      where \(\widehat{g}_{ij}^{\,k}\), \(\widehat{g}_{jk}^{\,i}\),
      \(\widehat{g}_{ik}^{\,j}\) are the transports of \(g_{ij}\), \(g_{jk}\),
      \(g_{ik}\) to \(V_{ijk}\) through the scheme-level triple transition
      \(t'_{ijk}\) and its compatibility \(t_{\mathrm{fac}}\), reassociated by the
      pullback comparison \cref{def:modules_pullbackComp} so that all three land on
      the same sheaf -- the multiplicative, \(\mathrm{GL}_d\)-valued cocycle
      condition.
  \end{enumerate}
  Conditions (C1) and (C2) are exactly the descent datum for the Zariski cover
  \(\{\iota_i\}\); it is subject to them that the gluing is well-defined. Out of this
  data the construction produces a glued sheaf of \(\mathcal{O}_X\)-modules
  \(\mathcal{M}\) on the glued scheme \(X\).

  \emph{Construction.} The category \(\ShMod(\mathcal{O}_X)\) of sheaves of
  \(\mathcal{O}_X\)-modules has all limits, and both pushforward along a morphism and
  these limits preserve the sheaf condition; the glued sheaf \(\mathcal{M}\) is therefore
  built directly as a limit -- an \emph{equalizer of pushforwards} -- rather than through a
  hand-built presheaf of compatible families. Concretely, form the two parallel maps of
  sheaves of \(\mathcal{O}_X\)-modules
  \[
    \prod_{i} (\iota_i)_{*}\,\mathcal{M}_i
      \;\;\underset{b}{\overset{a}{\rightrightarrows}}\;\;
    \prod_{i,j} (\iota_i \circ f_{ij})_{*}\,\bigl(f_{ij}^{*}\mathcal{M}_i\bigr),
  \]
  where, on the \((i,j)\)-component, leg \(a\) pushes the \(i\)-th factor forward along the
  overlap immersion \(f_{ij}\) using the unit of the pullback--pushforward adjunction
  together with the pushforward-composition comparison
  \((\iota_i)_{*} \cong (\iota_i \circ f_{ij})_{*}\circ f_{ij}^{*}\) (the dual of
  \cref{def:modules_pullbackComp}), and leg \(b\) pushes the \(j\)-th factor forward and
  additionally transports it across the inverse transition isomorphism \((g_{ij})^{-1}\)
  and the glue datum's compatibility identity, so that both legs land on the common overlap
  sheaf \((\iota_i \circ f_{ij})_{*}(f_{ij}^{*}\mathcal{M}_i)\). The glued sheaf is the
  equalizer
  \[
    \mathcal{M} \;=\; \operatorname{eq}(a,b) \;\hookrightarrow\; \prod_{i} (\iota_i)_{*}\,\mathcal{M}_i,
  \]
  the sub-sheaf of tuples of pushed-forward chart sections whose two overlap restrictions
  agree -- precisely the compatible families, now realised as a categorical limit rather than
  a bespoke presheaf. The cocycle conditions (C1)/(C2) on the family \((g_{ij})\) are
  \emph{not} consumed in forming this equalizer object: the limit exists for any family of
  transition maps. They are instead what make the equalizer descend correctly, and are
  exactly the hypotheses pinned down downstream when the restriction isomorphisms
  \cref{def:gr_modules_glueRestrictionIso} are produced. This specialises effective descent
  of sheaves of modules to the concrete open cover \(\{\iota_i : U_i \hookrightarrow X\}\)
  supplied by a scheme glue datum, and is the form in which the Grassmannian's
  tautological bundle and quotient are built.

  The glued sheaf comes with three further structures, each recorded as its own block
  downstream: the \emph{restriction isomorphisms} \(\rho_i : \iota_i^{*}\mathcal{M}
  \xrightarrow{\sim} \mathcal{M}_i\) compatible with the \(g_{ij}\) over the overlaps
  \cref{def:gr_modules_glueRestrictionIso}; the \emph{characterisation up to unique
  isomorphism} by the restriction property \cref{lem:gr_modules_glue_unique}; and the
  \emph{descent of morphisms} -- a family \(\phi_i : \mathcal{M}_i \to \mathcal{N}_i\)
  commuting with the transition isomorphisms glues to a unique \(\phi : \mathcal{M} \to
  \mathcal{N}\) with \(\rho_i^{\mathcal{N}} \circ \iota_i^{*}\phi = \phi_i \circ
  \rho_i^{\mathcal{M}}\) \cref{def:gr_modules_glueHom}.
\end{definition}

\begin{definition}[Restriction isomorphisms of the glued sheaf]
  \label{def:gr_modules_glueRestrictionIso}
  % NOTE: the Lean decl `glueRestrictionIso` is pinned canonically by
  % `def:glueRestrictionIso` in Picard_GlueDescent.tex (its 1:1 home after the
  % iter-067 file split). This block is the downstream consumer's cross-reference
  % node; it carries no \lean{} pin to avoid a duplicate.
  \uses{def:glueRestrictionIso, def:scheme_modules_glue, lem:over_restrict_pullback_iso,
    lem:modules_pullback_basechange_transport, lem:gr_glueData_bridges}
  Let \(\mathcal{M}\) be the glued sheaf of \cref{def:scheme_modules_glue}. For each
  chart index \(i\) there is a \emph{restriction isomorphism} of sheaves of
  \(\mathcal{O}_{U_i}\)-modules
  \[
    \rho_i \;:\; \iota_i^{*}\,\mathcal{M} \;\xrightarrow{\ \sim\ }\; \mathcal{M}_i,
  \]
  compatible with the transition isomorphisms: over each overlap \(V_{ij}\) the two
  restrictions \(\rho_i\) and \(\rho_j\), transported to the common overlap, differ
  exactly by \(g_{ij}\).
\end{definition}
\begin{proof}
  \uses{def:scheme_modules_glue, lem:over_restrict_pullback_iso,
    lem:modules_pullback_basechange_transport, lem:gr_glueData_bridges}
  By the equalizer-of-pushforwards description \cref{def:scheme_modules_glue}, a section of
  \(\mathcal{M}\) over an open of the form \(\iota_i^{-1}V\), restricted to the chart
  \(U_i\), is recovered from and recovers its \(i\)-th component
  \(s_i \in \Gamma(\mathcal{M}_i, \cdot)\) -- the \(i\)-th factor of the product
  \(\prod_i (\iota_i)_{*}\mathcal{M}_i\) into which the equalizer embeds; the
  chart-restriction bridge \cref{lem:over_restrict_pullback_iso} identifies the pullback
  \(\iota_i^{*}\mathcal{M}\) with this restriction, and the \(i\)-th projection
  \((s_\bullet) \mapsto s_i\) is the isomorphism \(\rho_i\) (with inverse extending a chart
  section by the overlap relation).
  The overlap compatibility is the defining relation of the compatible families: over
  \(V_{ij}\) the components \(s_i, s_j\) are related by \(g_{ij}\), which is precisely the
  statement that \(\rho_i\) and \(\rho_j\) differ by \(g_{ij}\) after transport to
  \(V_{ij}\); the triple-overlap coherence needed for this to be consistent is (C2),
  whose transports are aligned by \cref{lem:modules_pullback_basechange_transport} and
  \cref{lem:gr_glueData_bridges}.
\end{proof}

\begin{lemma}[Uniqueness of the glued sheaf up to canonical isomorphism]
  \label{lem:gr_modules_glue_unique}
  \lean{AlgebraicGeometry.Scheme.Modules.glue_unique}
  % NOTE: forward declaration (planned work); the \lean{} target is not yet realised.
  \uses{def:scheme_modules_glue, def:gr_modules_glueRestrictionIso,
    lem:sheaf_existsUnique_gluing_mathlib}
  Suppose \(\mathcal{M}'\) is another sheaf of \(\mathcal{O}_X\)-modules equipped with
  isomorphisms \(\rho'_i : \iota_i^{*}\mathcal{M}' \xrightarrow{\sim} \mathcal{M}_i\)
  compatible with the \(g_{ij}\) over the overlaps. Then there is a unique isomorphism
  \(\mathcal{M} \xrightarrow{\sim} \mathcal{M}'\) intertwining the restriction
  isomorphisms, i.e.\ with \(\rho'_i \circ \iota_i^{*}(\cdot) = \rho_i\) for every \(i\).
\end{lemma}
\begin{proof}
  \uses{def:gr_modules_glueRestrictionIso, lem:sheaf_existsUnique_gluing_mathlib}
  A morphism of sheaves of modules on \(X\) is determined by, and may be assembled from,
  its restrictions to the members of the open cover \(\{\iota_i\}\), agreeing on the
  overlaps -- the locality and gluing of sections
  (\cref{lem:sheaf_existsUnique_gluing_mathlib}) applied to the internal hom. The
  intertwining condition forces the restriction over \(U_i\) to be
  \((\rho'_i)^{-1} \circ \rho_i : \iota_i^{*}\mathcal{M} \to \iota_i^{*}\mathcal{M}'\);
  these local isomorphisms agree on each overlap \(V_{ij}\) because both
  \((\rho_i)\) and \((\rho'_i)\) are compatible with the \emph{same} transition family
  \((g_{ij})\) (\cref{def:gr_modules_glueRestrictionIso}). By locality they glue to a
  unique global isomorphism \(\mathcal{M} \xrightarrow{\sim} \mathcal{M}'\) with the
  required property; uniqueness is the separation half of the same gluing.
\end{proof}

\begin{definition}[Descent of morphisms along the cover]
  \label{def:gr_modules_glueHom}
  \lean{AlgebraicGeometry.Scheme.Modules.glueHom}
  % NOTE: forward declaration (planned work); the \lean{} target is not yet realised.
  \uses{def:scheme_modules_glue, def:gr_modules_glueRestrictionIso,
    lem:sheaf_existsUnique_gluing_mathlib}
  Let \((\mathcal{M}_i, g_{ij})\) and \((\mathcal{N}_i, h_{ij})\) be two glued data
  over the same glue datum \(D\), with glued sheaves \(\mathcal{M}, \mathcal{N}\) and
  restriction isomorphisms \(\rho^{\mathcal{M}}_i, \rho^{\mathcal{N}}_i\). Given a family
  of \(\mathcal{O}_{U_i}\)-linear morphisms \(\phi_i : \mathcal{M}_i \to \mathcal{N}_i\)
  commuting with the transition isomorphisms on every overlap (i.e.\ \(h_{ij}\) intertwines
  the pullbacks of \(\phi_i\) and \(\phi_j\) over \(V_{ij}\)), there is a unique morphism
  of sheaves of \(\mathcal{O}_X\)-modules \(\phi : \mathcal{M} \to \mathcal{N}\) with
  \[
    \rho^{\mathcal{N}}_i \circ \iota_i^{*}\phi \;=\; \phi_i \circ \rho^{\mathcal{M}}_i
    \qquad \text{for every } i.
  \]
\end{definition}
\begin{proof}
  \uses{def:gr_modules_glueRestrictionIso, lem:sheaf_existsUnique_gluing_mathlib}
  On each chart the prescribed identity defines a local morphism
  \((\rho^{\mathcal{N}}_i)^{-1} \circ \phi_i \circ \rho^{\mathcal{M}}_i :
  \iota_i^{*}\mathcal{M} \to \iota_i^{*}\mathcal{N}\). The commutation of the \(\phi_i\)
  with the transition isomorphisms \((g_{ij}, h_{ij})\), together with the overlap
  compatibility of the \(\rho\)-families (\cref{def:gr_modules_glueRestrictionIso}), makes
  these local morphisms agree on every overlap \(V_{ij}\). By the locality and gluing of
  sections (\cref{lem:sheaf_existsUnique_gluing_mathlib}) applied to the internal hom they
  promote to a unique global morphism \(\phi : \mathcal{M} \to \mathcal{N}\) restricting
  to each \(\phi_i\) through the \(\rho_i\) -- the descent of \(\mathrm{Hom}\) along the
  cover.
\end{proof}

\subsection{Pullback of free sheaves along an arbitrary morphism}
\label{sec:grquot_pullback_free}

The functor of points (\cref{def:grassmannian_functor}) acts on morphisms by pullback, so
it needs the pullback of the trivial bundle to be again trivial --
\(\varphi^{*}\mathcal{O}_T^{\,\oplus I} \cong \mathcal{O}_{T'}^{\,\oplus I}\) -- for an
\emph{arbitrary} scheme morphism \(\varphi\). Mathlib's pullback-of-free comparison applies
only when the underlying site functor \(\mathrm{Opens.map}\,\varphi\) is \emph{final}; the
following three blocks remove that restriction for every morphism, and the resulting
pullback machinery preserves rank-\(d\) local freeness.

\begin{lemma}\leanok
[The opens-pullback functor is final]
  \label{lem:gr_opensMap_final}
  \lean{AlgebraicGeometry.Scheme.Modules.opensMap_final}
  For any morphism of schemes \(\varphi : T' \to T\), the inverse-image functor on opens
  \(\mathrm{Opens.map}\,\varphi : \mathrm{Opens}(T) \to \mathrm{Opens}(T')\),
  \(U \mapsto \varphi^{-1}U\), is \emph{final}.
\end{lemma}

\begin{proof}


\leanok
  Finality means that for every open \(V \subseteq T'\) the structured-arrow category of
  opens \(U \subseteq T\) with \(V \le \varphi^{-1}U\) is connected. That category is
  nonempty and has the terminal object \(U = \top\) -- one always has
  \(V \le \varphi^{-1}\top = T'\), and every object admits a unique arrow to \(\top\) --
  and a category with a terminal object is connected.
\end{proof}

\begin{lemma}\leanok
[Pullback of a free sheaf is free]
  \label{lem:gr_pullbackFreeIso}
  \lean{AlgebraicGeometry.Scheme.Modules.pullbackFreeIso}
  \uses{lem:gr_opensMap_final}
  For any scheme morphism \(\varphi : T' \to T\) and index type \(I\), there is a canonical
  isomorphism of sheaves of \(\mathcal{O}_{T'}\)-modules
  \[
    \varphi^{*}\bigl(\mathcal{O}_T^{\,\oplus I}\bigr)
      \;\xrightarrow{\ \sim\ }\;
    \mathcal{O}_{T'}^{\,\oplus I}.
  \]
\end{lemma}

\begin{proof}
  \uses{lem:gr_opensMap_final}
  \leanok
  Mathlib's pullback-of-free comparison produces this isomorphism whenever the site functor
  \(\mathrm{Opens.map}\,\varphi\) is final; that finality is \cref{lem:gr_opensMap_final},
  which holds for every morphism, so the comparison applies unconditionally.
\end{proof}

\begin{lemma}

\leanok
[Free-pullback comparison is natural in the base morphism]
  \label{lem:gr_pullbackFreeIso_eqToHom}
  \lean{AlgebraicGeometry.Scheme.Modules.pullbackFreeIso_eqToHom}
  \uses{lem:gr_pullbackFreeIso}
  Let \(\varphi, \psi : T' \to T\) be morphisms with \(\varphi = \psi\) and \(I\) an index
  type. The two free-pullback comparisons \cref{lem:gr_pullbackFreeIso} along \(\varphi\) and
  \(\psi\) agree after transporting the differing pullback sources by the canonical
  \(\mathrm{eqToHom}\): the composite of that transport with
  \(\mathrm{pullbackFreeIso}(\psi)\) equals \(\mathrm{pullbackFreeIso}(\varphi)\).
\end{lemma}
\begin{proof}
  \uses{lem:gr_pullbackFreeIso}
  \leanok
  Substituting the equality \(\varphi = \psi\) collapses the \(\mathrm{eqToHom}\) to the
  identity and the two comparisons coincide.
\end{proof}

\begin{lemma}

\leanok
[Iso-level free-pullback cancellation]
  \label{lem:gr_pullbackFreeIso_trans_symm_eqToIso}
  \lean{AlgebraicGeometry.Scheme.Modules.pullbackFreeIso_trans_symm_eqToIso}
  \uses{lem:gr_pullbackFreeIso}
  For equal base morphisms \(\varphi = \psi : T' \to T\) and index type \(I\), the composite
  \(\mathrm{pullbackFreeIso}(\varphi) \circ \mathrm{pullbackFreeIso}(\psi)^{-1}\) is the
  \(\mathrm{eqToIso}\) transport between the (differing) pullback sources of the free sheaf
  \(\mathcal{O}_T^{\,\oplus I}\). Stated and proved generically in the base morphisms, so its
  use never forces evaluation of a concrete immersion; this is the load-bearing replacement
  for the cast chain in the self-identity \cref{lem:gr_bundleCocycle_id}.
\end{lemma}
\begin{proof}
  \uses{lem:gr_pullbackFreeIso}
  \leanok
  Substituting \(\varphi = \psi\) turns the composite into
  \(\mathrm{pullbackFreeIso}(\varphi) \circ \mathrm{pullbackFreeIso}(\varphi)^{-1} =
  \mathrm{id}\), which is the \(\mathrm{eqToIso}\) of the now-trivial source equality.
\end{proof}

\begin{lemma}\leanok
[Pullback preserves rank-\(d\) local freeness]
  \label{lem:gr_pullback_isLocallyFreeOfRank}
  \lean{AlgebraicGeometry.Scheme.Modules.pullback_isLocallyFreeOfRank}
  \uses{lem:gr_pullbackFreeIso, def:is_locally_free_of_rank, def:modules_pullbackComp}
  Let \(\varphi : T' \to T\) be a morphism of schemes and \(\mathcal{M}\) a sheaf of
  \(\mathcal{O}_T\)-modules locally free of rank \(d\)
  (\cref{def:is_locally_free_of_rank}). Then the pullback \(\varphi^{*}\mathcal{M}\) is
  locally free of rank \(d\) on \(T'\).
\end{lemma}

\begin{proof}
  \uses{lem:gr_pullbackFreeIso, def:is_locally_free_of_rank, def:modules_pullbackComp}
  \leanok
  Choose an open cover \(\{U_i\}\) of \(T\) trivialising \(\mathcal{M}\),
  \(\mathcal{M}|_{U_i} \cong \mathcal{O}_{U_i}^{\,d}\). The preimages
  \(\{\varphi^{-1}U_i\}\) cover \(T'\). On each member the factorisation
  \(\varphi \circ \iota_{\varphi^{-1}U_i} = (\varphi|_{U_i}) \circ \iota_{U_i}\) of the
  restricted morphism, together with the pseudofunctor comparison of iterated pullbacks
  (\cref{def:modules_pullbackComp}), identifies
  \((\varphi^{*}\mathcal{M})|_{\varphi^{-1}U_i}\) with
  \((\varphi|_{U_i})^{*}\mathcal{O}_{U_i}^{\,d}\); the latter is
  \(\mathcal{O}_{\varphi^{-1}U_i}^{\,d}\) by \cref{lem:gr_pullbackFreeIso}. Hence
  \(\varphi^{*}\mathcal{M}\) is locally free of rank \(d\).
\end{proof}

\section{The tautological quotient on the charts}
\label{sec:grquot_chart}

\emph{Infrastructure: scalar endomorphisms.} Realising a matrix of regular
functions as a morphism of free sheaves of modules has no off-the-shelf primitive
in Mathlib; it is assembled from two small project-local helpers that turn a global
function into a scalar endomorphism of the structure sheaf, after which the matrix
is reconstituted entry-by-entry over a finite biproduct.

\begin{definition}[Global section of the unit module]
  \label{def:gr_globalUnitSection}
  \lean{AlgebraicGeometry.Grassmannian.globalUnitSection}
  \uses{lem:moduleUnit_mathlib}
  Let \(X\) be a scheme and \(a \in \Gamma(X, \mathcal{O}_X)\) a global section of
  the structure sheaf. Viewing \(\mathcal{O}_X\) as the unit module
  \(\mathbf{1} \in \ShMod(\mathcal{O}_X)\) (\cref{lem:moduleUnit_mathlib}), the
  section \(a\) determines a global section of \(\mathbf{1}\): on each open
  \(V \subseteq X\) take the restriction \(a|_V \in \Gamma(V, \mathcal{O}_X)\), the
  family \((a|_V)_V\) being compatible with the restriction maps because
  restriction is functorial. This is the section of the unit sheaf attached to a
  global function.
\end{definition}

\begin{definition}[Scalar endomorphism of \(\mathcal{O}_X\)]
  \label{def:gr_scalarEnd}
  \lean{AlgebraicGeometry.Grassmannian.scalarEnd}
  \uses{def:gr_globalUnitSection, lem:moduleUnit_mathlib}
  For \(a \in \Gamma(X, \mathcal{O}_X)\) the \emph{scalar endomorphism}
  \(\mathrm{scalarEnd}(a) : \mathcal{O}_X \to \mathcal{O}_X\) is the
  \(\mathcal{O}_X\)-linear endomorphism of the unit module given by multiplication
  by \(a\). It is obtained from the global section of \cref{def:gr_globalUnitSection}
  under the canonical identification of endomorphisms of the unit module
  \(\mathbf{1}\) with its global sections, \(\End(\mathbf{1}) \cong
  \Gamma(X, \mathbf{1})\); concretely, on each open \(V\) it acts as
  multiplication by \(a|_V\). These scalar endomorphisms are the matrix-entry
  building blocks of the chart quotient.
\end{definition}

\begin{lemma}[Scalar endomorphism of the unit is the identity]
  \label{lem:gr_scalarEnd_one}
  \lean{AlgebraicGeometry.Grassmannian.scalarEnd_one}
  \uses{def:gr_scalarEnd}
  The scalar endomorphism attached to the unit function \(1 \in
  \Gamma(X, \mathcal{O}_X)\) is the identity of the unit module:
  \(\mathrm{scalarEnd}(1) = \mathrm{id}_{\mathbf{1}}\).
\end{lemma}

\begin{proof}
  \uses{def:gr_scalarEnd}
  Multiplication by \(1\) is the identity on every section, so under the
  identification \(\End(\mathbf{1}) \cong \Gamma(X, \mathbf{1})\) the section \(1\)
  corresponds to \(\mathrm{id}_{\mathbf{1}}\).
\end{proof}

\begin{lemma}[Scalar endomorphism of the zero function is zero]
  \label{lem:gr_scalarEnd_zero}
  \lean{AlgebraicGeometry.Grassmannian.scalarEnd_zero}
  \uses{def:gr_scalarEnd}
  The scalar endomorphism attached to the zero function \(0 \in
  \Gamma(X, \mathcal{O}_X)\) is the zero endomorphism:
  \(\mathrm{scalarEnd}(0) = 0\).
\end{lemma}

\begin{proof}
  \uses{def:gr_scalarEnd}
  Multiplication by \(0\) sends every section to \(0\), so the corresponding
  endomorphism of the unit module is the zero map.
\end{proof}

\begin{lemma}[Composition of scalar endomorphisms is multiplication]
  \label{lem:gr_scalarEnd_comp}
  \lean{AlgebraicGeometry.Grassmannian.scalarEnd_comp}
  \uses{def:gr_scalarEnd}
  For \(a, b \in \Gamma(X, \mathcal{O}_X)\) the scalar endomorphisms compose by
  multiplication: \(\mathrm{scalarEnd}(a) \circ \mathrm{scalarEnd}(b) =
  \mathrm{scalarEnd}(a\,b)\).
\end{lemma}
\begin{proof}
  \uses{def:gr_scalarEnd}
  Under the identification \(\End(\mathbf{1}) \cong \Gamma(X, \mathbf{1})\) the
  scalar endomorphisms act as multiplication by \(a\) and by \(b\); their composite
  is multiplication by \(a\,b\), because the section ring is commutative and
  restriction maps are ring homomorphisms.
\end{proof}

\begin{lemma}[Scalar endomorphism is additive]
  \label{lem:gr_scalarEnd_add}
  \lean{AlgebraicGeometry.Grassmannian.scalarEnd_add}
  \uses{def:gr_scalarEnd}
  For \(a, b \in \Gamma(X, \mathcal{O}_X)\) one has \(\mathrm{scalarEnd}(a + b) =
  \mathrm{scalarEnd}(a) + \mathrm{scalarEnd}(b)\).
\end{lemma}
\begin{proof}
  \uses{def:gr_scalarEnd}
  Multiplication by \(a + b\) equals multiplication by \(a\) plus multiplication by
  \(b\) on every section, since the restriction maps of \(\mathcal{O}_X\) are additive;
  under \(\End(\mathbf{1}) \cong \Gamma(X, \mathbf{1})\) this is the stated additivity.
\end{proof}

\begin{lemma}

[Scalar endomorphism of a finite sum]
  \label{lem:gr_scalarEnd_sum}
  \lean{AlgebraicGeometry.Grassmannian.scalarEnd_sum}
  \uses{def:gr_scalarEnd, lem:gr_scalarEnd_zero, lem:gr_scalarEnd_add}
  For a finite index set \(s\) and a family \(f : \iota \to \Gamma(X, \mathcal{O}_X)\),
  \(\mathrm{scalarEnd}\bigl(\sum_{i \in s} f_i\bigr) = \sum_{i \in s} \mathrm{scalarEnd}(f_i)\).
\end{lemma}
\begin{proof}
  \uses{lem:gr_scalarEnd_zero, lem:gr_scalarEnd_add}
  Induction on the finite set \(s\): the empty case is \(\mathrm{scalarEnd}(0) = 0\)
  (\cref{lem:gr_scalarEnd_zero}) and the insertion step is additivity
  (\cref{lem:gr_scalarEnd_add}).
\end{proof}

\emph{Infrastructure: matrix automorphisms of the free sheaf.} To realise the
\(\mathrm{GL}_d\) bundle transitions one must promote an invertible \(d \times d\)
matrix of global functions to an \emph{automorphism} of the free rank-\(d\) sheaf
\(\mathcal{O}_S^{\,d}\), exactly as the chart quotient \cref{def:gr_chart_quotient}
realises the universal matrix. The next blocks assemble this matrix-to-morphism
functor and record its two structural identities -- that it carries matrix
multiplication to composition and the identity matrix to the identity -- which are
the algebraic heart of the bundle cocycle.

\begin{definition}[Matrix endomorphism of the free sheaf]
  \label{def:gr_matrixEnd}
  \lean{AlgebraicGeometry.Grassmannian.matrixEnd}
  \uses{def:gr_scalarEnd}
  Let \(S\) be a scheme, \(d \in \mathbb{N}\), and
  \(M \in \operatorname{Mat}_{d \times d}(\Gamma(S, \mathcal{O}_S))\) a matrix of
  global functions. The \emph{matrix endomorphism} \(\mathrm{matrixEnd}(M)\) is the
  endomorphism of the free rank-\(d\) sheaf \(\mathcal{O}_S^{\,d}\) whose
  \((p,q)\)-component, under the presentation of \(\mathcal{O}_S^{\,d}\) as the
  finite biproduct of \(d\) copies of the unit module \(\mathbf{1}\), is the scalar
  endomorphism \(\mathrm{scalarEnd}(M_{p,q})\) (\cref{def:gr_scalarEnd}). It is
  assembled from these components by the universal property of the rank-\(d\)
  biproduct, exactly as the chart quotient \(u^I\) is assembled
  (\cref{def:gr_chart_quotient}). This is the substitute for the (non-existent)
  Mathlib ``matrix \(\mapsto\) endomorphism of a free sheaf'' primitive.
\end{definition}

\begin{lemma}[Matrix endomorphism carries multiplication to composition]
  \label{lem:gr_matrixEnd_comp}
  \lean{AlgebraicGeometry.Grassmannian.matrixEnd_comp}
  \uses{def:gr_matrixEnd, lem:gr_scalarEnd_comp, lem:gr_scalarEnd_add}
  For \(M, N \in \operatorname{Mat}_{d \times d}(\Gamma(S, \mathcal{O}_S))\) the matrix
  endomorphisms compose according to the matrix product, with the order reversed by
  the contravariance of the component indexing:
  \[
    \mathrm{matrixEnd}(M) \,\circ\, \mathrm{matrixEnd}(N)
      \;=\; \mathrm{matrixEnd}(N \cdot M).
  \]
\end{lemma}
\begin{proof}
  \uses{def:gr_matrixEnd, lem:gr_scalarEnd_comp, lem:gr_scalarEnd_add}
  The composition of two biproduct-assembled morphisms is again biproduct-assembled,
  with \((i,q)\)-component the sum over the middle index \(p\) of the composites of
  components, \(\sum_p \mathrm{component}_{i,p} \circ \mathrm{component}_{p,q}\) (the
  categorical matrix product over the biproduct). Each composite of scalar
  endomorphisms is the scalar endomorphism of the product,
  \(\mathrm{scalarEnd}(a)\circ\mathrm{scalarEnd}(b) = \mathrm{scalarEnd}(ab)\)
  (\cref{lem:gr_scalarEnd_comp}), and a finite sum of scalar endomorphisms is the
  scalar endomorphism of the sum (\cref{lem:gr_scalarEnd_add}); hence the \((i,q)\)-component
  of the composite is \(\mathrm{scalarEnd}\bigl(\sum_p M_{q,p} N_{p,i}\bigr) =
  \mathrm{scalarEnd}\bigl((N\cdot M)_{q,i}\bigr)\), which is precisely the
  \((i,q)\)-component of \(\mathrm{matrixEnd}(N\cdot M)\). The reversal of the factor
  order is the contravariance of the column/component indexing.
\end{proof}

\begin{lemma}[Matrix endomorphism of the identity is the identity]
  \label{lem:gr_matrixEnd_one}
  \lean{AlgebraicGeometry.Grassmannian.matrixEnd_one}
  \uses{def:gr_matrixEnd, lem:gr_scalarEnd_one, lem:gr_scalarEnd_zero}
  The matrix endomorphism of the identity matrix is the identity morphism of
  \(\mathcal{O}_S^{\,d}\): \(\mathrm{matrixEnd}(1_{d\times d}) =
  \mathrm{id}_{\mathcal{O}_S^{\,d}}\).
\end{lemma}
\begin{proof}
  \uses{def:gr_matrixEnd, lem:gr_scalarEnd_one, lem:gr_scalarEnd_zero}
  The diagonal components are \(\mathrm{scalarEnd}(1) = \mathrm{id}\)
  (\cref{lem:gr_scalarEnd_one}) and the off-diagonal components are
  \(\mathrm{scalarEnd}(0) = 0\) (\cref{lem:gr_scalarEnd_zero}); the morphism with this
  component matrix is the identity of the biproduct.
\end{proof}

\begin{definition}[Free-sheaf automorphism of an invertible matrix]
  \label{def:gr_matrixToFreeIso}
  \lean{AlgebraicGeometry.Grassmannian.matrixToFreeIso}
  \uses{def:gr_matrixEnd, lem:gr_matrixEnd_comp, lem:gr_matrixEnd_one}
  Let \(M, N \in \operatorname{Mat}_{d \times d}(\Gamma(S, \mathcal{O}_S))\) be mutually
  inverse, \(M N = 1\) and \(N M = 1\). The \emph{matrix automorphism}
  \(\mathrm{matrixToFreeIso}(M, N)\) is the isomorphism of free sheaves
  \(\mathcal{O}_S^{\,d} \xrightarrow{\sim} \mathcal{O}_S^{\,d}\) with forward map
  \(\mathrm{matrixEnd}(M)\) and inverse map \(\mathrm{matrixEnd}(N)\); the two
  composites are identities by \cref{lem:gr_matrixEnd_comp} (turning the products
  \(NM\) and \(MN\) into the relevant compositions) and \cref{lem:gr_matrixEnd_one}.
  Its forward map is \(\mathrm{matrixEnd}(M)\) by construction
  (\cref{lem:gr_matrixToFreeIso_hom}).
\end{definition}

\begin{lemma}[Forward map of the matrix automorphism]
  \label{lem:gr_matrixToFreeIso_hom}
  \lean{AlgebraicGeometry.Grassmannian.matrixToFreeIso_hom}
  \uses{def:gr_matrixToFreeIso, def:gr_matrixEnd}
  The forward map of \(\mathrm{matrixToFreeIso}(M, N)\) is \(\mathrm{matrixEnd}(M)\).
\end{lemma}
\begin{proof}
  \uses{def:gr_matrixToFreeIso, def:gr_matrixEnd}
  Immediate from the definition \cref{def:gr_matrixToFreeIso}, whose forward field is
  literally \(\mathrm{matrixEnd}(M)\).
\end{proof}

\begin{lemma}[Matrix automorphisms compose multiplicatively]
  \label{lem:gr_matrixToFreeIso_mul}
  \lean{AlgebraicGeometry.Grassmannian.matrixToFreeIso_mul}
  \uses{def:gr_matrixToFreeIso, lem:gr_matrixToFreeIso_hom, lem:gr_matrixEnd_comp}
  Let \(A, A', B, B'\) be \(d \times d\) matrices of global functions on \(S\) with
  \(A A' = A' A = 1\), \(B B' = B' B = 1\), so that \(\mathrm{matrixToFreeIso}(A, A')\)
  and \(\mathrm{matrixToFreeIso}(B, B')\) are automorphisms of \(\mathcal{O}_S^{\,d}\).
  Their forward maps compose to the forward map of the product automorphism, with the
  order reversed by the component contravariance:
  \[
    \mathrm{matrixToFreeIso}(A, A')_{\mathrm{hom}}
      \,\circ\, \mathrm{matrixToFreeIso}(B, B')_{\mathrm{hom}}
      \;=\; \mathrm{matrixEnd}(B \cdot A).
  \]
  Equivalently, composition of matrix automorphisms realises matrix multiplication of
  the underlying matrices. This is the linkage that turns the matrix-level cocycle
  identity into a composition identity of sheaf-of-module isomorphisms.
\end{lemma}
\begin{proof}
  \uses{def:gr_matrixToFreeIso, lem:gr_matrixToFreeIso_hom, lem:gr_matrixEnd_comp}
  By \cref{lem:gr_matrixToFreeIso_hom} the two forward maps are
  \(\mathrm{matrixEnd}(A)\) and \(\mathrm{matrixEnd}(B)\); their composite is
  \(\mathrm{matrixEnd}(A)\circ\mathrm{matrixEnd}(B) = \mathrm{matrixEnd}(B\cdot A)\) by
  \cref{lem:gr_matrixEnd_comp}.
\end{proof}

\begin{definition}

[Chart quotient \(u^I\)]
  \label{def:gr_chart_quotient}
  \lean{AlgebraicGeometry.Grassmannian.chartQuotientMap}
  \uses{def:gr_affine_chart, def:gr_universal_matrix, def:gr_scalarEnd}
  % SOURCE: [Nitsure], "Construction of Hilbert and Quot Schemes", \S 1,
  % sub-section "Universal quotient" (read from
  % references/nitsure-hilbert-quot-src/nitsure-hilbert-quot.tex, L898-L910).
  % SOURCE QUOTE:
  % "We next define a rank $d$ locally free sheaf $\U$ on $\grass(r,d)$
  %  together with a surjective homomorphism $\oplus^r\,{\mathcal O}_{\grass(r,d)} \to \U$.
  %  On each $U^I$ we define a surjective homomorphism
  %  $u^I :\oplus^r\,{\mathcal O}_{U^I} \to  \oplus^d\,{\mathcal O}_{U^I}$
  %  by the matrix $X^I$."
  \textit{Source: [Nitsure], \S 1.}
  Let \(I \subseteq \{1,\dots,r\}\) with \(\#(I) = d\). On the affine chart
  \(U^I = \Spec R^I\) (\cref{def:gr_affine_chart}) define the
  \(\mathcal{O}_{U^I}\)-linear homomorphism of trivial bundles
  \[
    u^I \;:\; \mathcal{O}_{U^I}^{\,r}
      \;\longrightarrow\; \mathcal{O}_{U^I}^{\,d}
  \]
  given by left multiplication by the universal \(d \times r\) matrix \(X^I\)
  (\cref{def:gr_universal_matrix}): on global sections it is the \(R^I\)-linear
  map \((R^I)^{r} \to (R^I)^{d}\), \(v \mapsto X^I v\). Since the \(I\)-th minor
  \(X^I_I\) of \(X^I\) is the identity matrix \(1_{d \times d}\), the composite of
  \(u^I\) with the inclusion of the \(I\)-indexed coordinates
  \(\mathcal{O}_{U^I}^{\,d} \hookrightarrow \mathcal{O}_{U^I}^{\,r}\) is the
  identity; hence \(u^I\) is a split surjection of locally free sheaves, with
  target the free rank-\(d\) sheaf \(\mathcal{O}_{U^I}^{\,d}\).

  \emph{Realisation.} The free sheaves \(\mathcal{O}_{U^I}^{\,r}\) and
  \(\mathcal{O}_{U^I}^{\,d}\) are the finite biproducts (coproducts) of copies of
  the unit module \(\mathbf{1} = \mathcal{O}_{U^I}\), indexed by \(\{1,\dots,r\}\)
  and \(\{1,\dots,d\}\). Each entry \(X^I_{p,q} \in R^I\) of the universal matrix
  (\cref{def:gr_universal_matrix}) is a global function on \(U^I = \Spec R^I\),
  obtained from \(R^I\) through the canonical identification
  \(\Gamma(\Spec R^I, \mathcal{O}) \cong R^I\); the corresponding scalar
  endomorphism \(\mathrm{scalarEnd}(X^I_{p,q}) : \mathbf{1} \to \mathbf{1}\)
  (\cref{def:gr_scalarEnd}) is the \((p,q)\)-component of \(u^I\). Assembling these
  components over the index biproducts -- the universal property of the biproduct
  packaging a matrix of component morphisms into a single morphism of the free
  sheaves -- yields \(u^I\). This entry-by-entry assembly via the finite biproduct
  is the substitute for a (non-existent) Mathlib ``matrix \(\mapsto\) morphism of
  free sheaves'' primitive.
\end{definition}

\begin{lemma}[The $I$-columns of the chart quotient are the identity]
  \label{lem:gr_chartQuotientMap_iFree}
  \lean{AlgebraicGeometry.Grassmannian.chartQuotientMap_ιFree}
  \uses{def:gr_chart_quotient, def:gr_universal_matrix, def:gr_scalarEnd,
    lem:gr_scalarEnd_one, lem:gr_scalarEnd_zero,
    lem:gr_universalMatrix_submatrix_self}
  For each \(k \in \{1,\dots,d\}\), the chart quotient \(u^I\)
  (\cref{def:gr_chart_quotient}) carries the \(\iota_I(k)\)-th basis section of the
  free sheaf \(\mathcal{O}_{U^I}^{\,r}\) -- where \(\iota_I(k)\) is the \(k\)-th
  element of \(I\) in increasing order -- to the \(k\)-th basis section of
  \(\mathcal{O}_{U^I}^{\,d}\); that is, the \(I\)-indexed columns of \(u^I\) form the
  identity:
  \[
    e_{\iota_I(k)} \,\circ\, u^I \;=\; e_k.
  \]
\end{lemma}

\begin{proof}
  \uses{def:gr_chart_quotient, def:gr_universal_matrix, def:gr_scalarEnd,
    lem:gr_scalarEnd_one, lem:gr_scalarEnd_zero,
    lem:gr_universalMatrix_submatrix_self}
  By the biproduct presentation of \(u^I\) (\cref{def:gr_chart_quotient}), the
  composite \(e_{\iota_I(k)} \circ u^I\) reads off the \(\iota_I(k)\)-th column of the
  component matrix, whose \(p\)-th entry is the scalar endomorphism
  \(\mathrm{scalarEnd}(X^I_{p,\iota_I(k)})\). The \(I\)-indexed columns of the
  universal matrix \(X^I\), read through the order isomorphism \(\{1,\dots,d\} \cong
  I\), form the \(d \times d\) identity (\cref{lem:gr_universalMatrix_submatrix_self}),
  so \(X^I_{p,\iota_I(k)} = \delta_{p,k}\). Hence the diagonal entry is
  \(\mathrm{scalarEnd}(1) = \mathrm{id}\) (\cref{lem:gr_scalarEnd_one}) and the
  off-diagonal entries are \(\mathrm{scalarEnd}(0) = 0\)
  (\cref{lem:gr_scalarEnd_zero}); the resulting column is the \(k\)-th coordinate
  inclusion, giving \(e_{\iota_I(k)} \circ u^I = e_k\).
\end{proof}

\begin{lemma}[The chart quotient is an epimorphism]
  \label{lem:gr_chartQuotientMap_epi}
  \lean{AlgebraicGeometry.Grassmannian.chartQuotientMap_epi}
  \uses{def:gr_chart_quotient}
  The chart quotient \(u^I : \mathcal{O}_{U^I}^{\,r} \to \mathcal{O}_{U^I}^{\,d}\)
  (\cref{def:gr_chart_quotient}) is an epimorphism of sheaves of modules.
\end{lemma}

\begin{proof}
  \uses{def:gr_chart_quotient, lem:gr_chartQuotientMap_iFree}
  We exhibit a section, so that \(u^I\) is a \emph{split} epimorphism and hence an
  epimorphism. Let \(\iota_I : \{1,\dots,d\} \hookrightarrow \{1,\dots,r\}\) be the
  inclusion of the \(I\)-indexed coordinates, \(k \mapsto \iota_I(k)\), the
  \(k\)-th element of \(I\) in increasing order. It induces a morphism of free
  sheaves \(s_I : \mathcal{O}_{U^I}^{\,d} \to \mathcal{O}_{U^I}^{\,r}\), the
  coordinate inclusion sending the \(k\)-th basis section to the
  \(\iota_I(k)\)-th. We claim
  \[
    u^I \circ s_I \;=\; \mathrm{id}_{\mathcal{O}_{U^I}^{\,d}}.
  \]
  Both sides are morphisms out of the free rank-\(d\) sheaf, which is the coproduct
  of \(d\) copies of the unit module; by the universal property of that coproduct
  it suffices to check the equality on each coordinate, i.e.\ to compare component
  matrices. The \((p,k)\)-component of \(u^I \circ s_I\) is the entry
  \(X^I_{p,\iota_I(k)}\) of the universal matrix at the \(I\)-indexed column, which
  is precisely the \(I\)-th minor \(X^I_I\) of \(X^I\). Since that minor is the
  \(d \times d\) identity (the columns of \(X^I\) indexed by \(I\), read through the
  order isomorphism \(\{1,\dots,d\} \cong I\), form the identity block --
  \cref{def:gr_universal_matrix}), one has \(X^I_{p,\iota_I(k)} = \delta_{p,k}\),
  so \(u^I \circ s_I\) has the identity component matrix and equals
  \(\mathrm{id}\). Hence \(s_I\) splits \(u^I\) and \(u^I\) is a split
  epimorphism, in particular an epimorphism.
\end{proof}

\section{The \(\mathrm{GL}_d\) bundle transition cocycle}
\label{sec:grquot_bundle_cocycle}

The universal quotient sheaf \(\mathcal{U}\) (\cref{def:gr_universal_quotient_sheaf}) is
obtained by feeding the gluing primitive \cref{def:scheme_modules_glue} the per-chart free
rank-\(d\) sheaves \(\mathcal{O}_{U^I}^{\,d}\), the bundle transitions
\(g_{I,J} = (X^I_J)^{-1}\), and the two module cocycle hypotheses (C1)/(C2) that primitive
requires. The matrix cocycle for the Cramer inverses \((X^I_J)^{-1}\) is already in hand --
it is the content of the scheme-level cocycle \cref{lem:gr_cocycle} and the minor identities
\cref{lem:gr_universalMinorInv_identities} -- but the gluing primitive consumes its
\emph{module-level} incarnation: an honest isomorphism of sheaves of modules together with
its self-identity and triple-overlap multiplicativity. The three blocks of this section
package exactly that, decoupling the construction of \(\mathcal{U}\) from the remaining
matrix-to-module transport bookkeeping. This is the only net-new infrastructure the
universal quotient requires beyond the gluing primitive itself.

\begin{lemma}

[Self-minor Cramer inverse is the identity]
  \label{lem:gr_universalMinorInv_self}
  \lean{AlgebraicGeometry.Grassmannian.universalMinorInv_self}
  \uses{def:gr_universalMinorInv, lem:gr_universalMatrix_submatrix_self}
  The Cramer inverse of the self-minor is the identity matrix:
  \((X^I_I)^{-1} = 1_{d \times d}\).
\end{lemma}
\begin{proof}
  \uses{lem:gr_universalMatrix_submatrix_self}
  The self-minor \(X^I_I\) is the identity block
  (\cref{lem:gr_universalMatrix_submatrix_self}), and \(1^{-1} = 1\).
\end{proof}

\begin{definition}[The bundle transition \(g_{I,J}\)]
  \label{def:gr_bundleTransition}
  \lean{AlgebraicGeometry.Grassmannian.bundleTransition}
  \uses{def:gr_universalMinorInv, def:gr_scalarEnd, lem:gr_pullbackFreeIso,
    lem:gr_minorDet_unit, def:gr_chart_quotient, def:gr_transition}
  Let \(I, J \subseteq \{1,\dots,r\}\) be size-\(d\) subsets and let
  \(f_{IJ} : U^I_J \hookrightarrow U^I\) and \(t_{IJ} : U^I_J \xrightarrow{\sim} U^J_I\) be
  the overlap immersion and transition of the chart cover (\cref{def:gr_transition}). The
  \emph{bundle transition} \(g_{I,J}\) is the isomorphism of sheaves of modules on the
  overlap \(U^I_J\)
  \[
    g_{I,J} \;:\;
      f_{IJ}^{*}\bigl(\mathcal{O}_{U^I}^{\,d}\bigr)
      \;\xrightarrow{\ \sim\ }\;
      (t_{IJ} \circ f_{JI})^{*}\bigl(\mathcal{O}_{U^J}^{\,d}\bigr)
  \]
  induced by the invertible \(d \times d\) matrix
  \(g_{I,J} = (X^I_J)^{-1} \in \mathrm{GL}_d(R^I_J)\) (\cref{def:gr_universalMinorInv}),
  the Cramer inverse of the localised \(J\)-minor, which is invertible because
  \(\det(X^I_J)\) is a unit of \(R^I_J\) (\cref{lem:gr_minorDet_unit}).

  \emph{Realisation.} The transition is built exactly as the chart quotient
  \cref{def:gr_chart_quotient}: each entry \((X^I_J)^{-1}_{p,q} \in R^I_J\) is a global
  function on the overlap, its scalar endomorphism
  \(\mathrm{scalarEnd}\bigl((X^I_J)^{-1}_{p,q}\bigr)\) (\cref{def:gr_scalarEnd}) is the
  \((p,q)\)-component, and the universal property of the rank-\(d\) biproduct assembles
  these into an automorphism of the free sheaf \(\mathcal{O}_{U^I_J}^{\,d}\); the
  source and target pullbacks of free sheaves are identified with this
  \(\mathcal{O}_{U^I_J}^{\,d}\) through the free-pullback comparisons
  \cref{lem:gr_pullbackFreeIso} along \(f_{IJ}\) and \(t_{IJ} \circ f_{JI}\), conjugating
  the matrix automorphism into the displayed isomorphism. Invertibility of the matrix makes
  \(g_{I,J}\) an isomorphism.
\end{definition}

\begin{lemma}[Self-identity of the bundle transition (C1)]
  \label{lem:gr_bundleCocycle_id}
  \lean{AlgebraicGeometry.Grassmannian.bundleTransition_self}
  \uses{def:gr_bundleTransition, lem:gr_universalMatrix_submatrix_self,
    lem:gr_universalMinorInv_identities, lem:gr_scalarEnd_one, lem:gr_scalarEnd_zero,
    lem:gr_pullbackFreeIso}
  The self-transition is the identity: \(g_{I,I} = \mathrm{id}\) on
  \(f_{II}^{*}\mathcal{O}_{U^I}^{\,d}\). This is the module-level form of the descent
  condition (C1) of \cref{def:scheme_modules_glue} for the bundle data \((g_{I,J})\).
\end{lemma}
\begin{proof}
  \uses{def:gr_bundleTransition, lem:gr_universalMatrix_submatrix_self,
    lem:gr_universalMinorInv_identities, lem:gr_scalarEnd_one, lem:gr_scalarEnd_zero}
  The \(I\)-th minor of the universal matrix is the identity block, \(X^I_I = 1_{d\times d}\)
  (\cref{lem:gr_universalMatrix_submatrix_self}), so its Cramer inverse is again the
  identity, \((X^I_I)^{-1} = 1_{d\times d}\) (\cref{lem:gr_universalMinorInv_identities}).
  Hence in the biproduct presentation of \cref{def:gr_bundleTransition} the diagonal
  components are \(\mathrm{scalarEnd}(1) = \mathrm{id}\) (\cref{lem:gr_scalarEnd_one}) and the
  off-diagonal components are \(\mathrm{scalarEnd}(0) = 0\) (\cref{lem:gr_scalarEnd_zero}), so
  the assembled automorphism is the identity of \(\mathcal{O}_{U^I_I}^{\,d}\); conjugation by
  the free-pullback comparisons preserves the identity, giving \(g_{I,I} = \mathrm{id}\).
\end{proof}

\begin{definition}

[Bundle transition data over the glue datum]
  \label{def:gr_bundleTransitionData}
  \lean{AlgebraicGeometry.Grassmannian.bundleTransitionData}
  \uses{def:gr_bundleTransition, def:scheme_modules_glue}
  The family of bundle transitions \((g_{I,J})\) (\cref{def:gr_bundleTransition}), repackaged
  over the index set of the Grassmannian glue datum \(\mathrm{theGlueData}(d,r)\) in the exact
  shape \(\bigl(\forall I\,J,\; f_{IJ}^{*}\mathcal{O}_{U^I}^{\,d} \cong (t_{IJ} \circ
  f_{JI})^{*}\mathcal{O}_{U^J}^{\,d}\bigr)\) that the gluing primitive
  \cref{def:scheme_modules_glue} consumes.
\end{definition}

\begin{lemma}[Cramer-inverse cocycle on the triple overlap]
  \label{lem:gr_bundleCocycle_matrix}
  \lean{AlgebraicGeometry.Grassmannian.bundleTransition_cocycle_matrix}
  \uses{def:gr_universalMinorInv, lem:gr_cocycle_imageMatrix_eq,
    lem:gr_universalMinorInv_identities, lem:gr_mul_submatrix_col,
    lem:gr_inv_mul_inv_mul_cancel, lem:gr_cocycle}
  Over the common triple-overlap ring (the localisation carrying the three charts'
  data to \(V_{IJK}\)) the Cramer inverses of the localised minors satisfy the
  multiplicative cocycle identity
  \[
    (X^J_K)^{-1}\,(X^I_J)^{-1} \;=\; (X^I_K)^{-1}.
  \]
  This is the pure matrix-algebra core of (C2), independent of any sheaf data.
\end{lemma}
\begin{proof}
  \uses{def:gr_universalMinorInv, lem:gr_cocycle_imageMatrix_eq,
    lem:gr_universalMinorInv_identities, lem:gr_mul_submatrix_col,
    lem:gr_inv_mul_inv_mul_cancel, lem:gr_cocycle}
  Write all matrices over the common triple-overlap ring. The image matrix of the
  transition \(\theta_{I,J}\) is \(X^J = (X^I_J)^{-1} X^I\), and its \(K\)-minor is
  computed by restricting to the \(K\)-indexed columns; since taking a column
  submatrix commutes with left multiplication (\cref{lem:gr_mul_submatrix_col}),
  \[
    X^J_K \;=\; \bigl((X^I_J)^{-1} X^I\bigr)_K \;=\; (X^I_J)^{-1}\, X^I_K .
  \]
  Both \(X^I_J\) and \(X^I_K\) have unit determinant
  (\cref{lem:gr_universalMinorInv_identities}), so this product of invertibles inverts
  by \((AB)^{-1} = B^{-1} A^{-1}\):
  \[
    (X^J_K)^{-1} \;=\; (X^I_K)^{-1}\,(X^I_J).
  \]
  Multiplying on the right by \((X^I_J)^{-1}\) and cancelling \(X^I_J (X^I_J)^{-1} = 1\)
  via the two-sided Cramer identity (\cref{lem:gr_universalMinorInv_identities}, the
  cancellation packaged in \cref{lem:gr_inv_mul_inv_mul_cancel}) gives
  \((X^J_K)^{-1} (X^I_J)^{-1} = (X^I_K)^{-1}\). The agreement of the three image
  matrices over the triple overlap that licenses these manipulations is the content of
  \cref{lem:gr_cocycle_imageMatrix_eq}, the matrix form of the scheme-level cocycle
  \cref{lem:gr_cocycle}.
\end{proof}

\begin{lemma}[Scalar endomorphism is natural under pullback]
  \label{lem:gr_scalarEnd_pullback}
  \lean{AlgebraicGeometry.Grassmannian.scalarEnd_pullback}
  \uses{def:gr_scalarEnd, lem:gr_pullbackFreeIso, def:modules_pullbackComp}
  Let \(p : T \to S\) be a morphism of schemes and \(a \in \Gamma(S, \mathcal{O}_S)\) a
  global function, with comorphism on global sections
  \(p^{\sharp} : \Gamma(S, \mathcal{O}_S) \to \Gamma(T, \mathcal{O}_T)\).
  Write \(q : p^{*}\mathbf{1} \xrightarrow{\sim} \mathbf{1}\) for the unit-pullback comparison
  (\cref{lem:gr_pullbackFreeIso} at a one-element index). Then the pullback of the scalar endomorphism
  \(\mathrm{scalarEnd}(a)\) is conjugate, through \(q\), to the scalar endomorphism of the
  base-changed function:
  \[
    p^{*}\bigl(\mathrm{scalarEnd}(a)\bigr)
      \;=\; q^{-1} \,\circ\, \mathrm{scalarEnd}\bigl(p^{\sharp} a\bigr) \,\circ\, q
  \]
  as endomorphisms of \(p^{*}\mathbf{1}\); equivalently
  \(q \circ p^{*}(\mathrm{scalarEnd}(a)) = \mathrm{scalarEnd}(p^{\sharp} a) \circ q\). This is the
  single-entry naturality atom underlying the matrix-level statement
  \cref{lem:gr_matrixEnd_pullback}.
\end{lemma}
\begin{proof}
  \uses{def:gr_scalarEnd, lem:gr_pullbackFreeIso}
  The scalar endomorphism \(\mathrm{scalarEnd}(a)\) (\cref{def:gr_scalarEnd}) is multiplication
  by the global section \(a\) on the unit module \(\mathbf 1 = \mathcal{O}_S\). Pulling back
  along \(p\) sends multiplication by \(a\) to multiplication by its image, and the comorphism
  \(p^{\sharp}\) is precisely the action of pullback on global functions, so
  \(p^{*}(\mathrm{scalarEnd}(a))\) is multiplication by \(p^{\sharp}a\) on \(p^{*}\mathbf 1\). The
  comparison \(q\) is an isomorphism of unit modules intertwining this multiplication on
  \(p^{*}\mathbf 1\) with multiplication by \(p^{\sharp}a\) on \(\mathbf 1\); conjugating back
  along \(q\) therefore yields the claimed identity. The verification is a single computation on
  the unit module, where multiplication by a function commutes with the natural comparison
  \(q\).
\end{proof}

\begin{lemma}[Matrix endomorphism is natural under pullback]
  \label{lem:gr_matrixEnd_pullback}
  \lean{AlgebraicGeometry.Grassmannian.matrixEnd_pullback}
  \uses{def:gr_matrixEnd, def:gr_scalarEnd, lem:gr_scalarEnd_pullback, lem:gr_pullbackFreeIso,
    def:modules_pullbackComp}
  Let \(p : T \to S\) be a morphism of schemes, \(d \in \mathbb{N}\), and
  \(M \in \operatorname{Mat}_{d\times d}(\Gamma(S, \mathcal{O}_S))\). Let
  \(p^{\sharp}M \in \operatorname{Mat}_{d\times d}(\Gamma(T, \mathcal{O}_T))\) be the entrywise
  image of \(M\) under the comorphism \(p^{\sharp}\), and let
  \(Q : p^{*}(\mathcal{O}_S^{\,d}) \xrightarrow{\sim} \mathcal{O}_T^{\,d}\) be the free-pullback
  comparison \cref{lem:gr_pullbackFreeIso} at the rank-\(d\) index set. Then the
  pullback of the matrix endomorphism \(\mathrm{matrixEnd}(M)\) is conjugate, through \(Q\), to
  the matrix endomorphism of the base-changed matrix:
  \[
    p^{*}\bigl(\mathrm{matrixEnd}(M)\bigr)
      \;=\; Q^{-1} \,\circ\, \mathrm{matrixEnd}(p^{\sharp}M) \,\circ\, Q .
  \]
\end{lemma}
\begin{proof}
  \uses{def:gr_matrixEnd, lem:gr_scalarEnd_pullback, lem:gr_pullbackFreeIso,
    def:modules_pullbackComp}
  By \cref{def:gr_matrixEnd}, \(\mathrm{matrixEnd}(M)\) is the biproduct-assembled endomorphism
  of \(\mathcal{O}_S^{\,d} = \bigoplus_{k} \mathbf 1\) whose \((k,\ell)\)-component is
  \(\mathrm{scalarEnd}(M_{k,\ell})\). The pullback functor \(p^{*}\) is a left adjoint, hence
  additive and biproduct-preserving; through the coproduct comparison underlying the
  free-pullback isomorphism \cref{lem:gr_pullbackFreeIso} it carries the biproduct presentation
  of \(\mathcal{O}_S^{\,d}\) to that of \(\mathcal{O}_T^{\,d}\) and the assembled morphism to the
  morphism reassembled from the pulled-back components, the reassociation of the iterated
  pullbacks being \cref{def:modules_pullbackComp}. On each single \((k,\ell)\)-component the
  pulled-back scalar endomorphism is, by \cref{lem:gr_scalarEnd_pullback}, the
  comparison-conjugate of \(\mathrm{scalarEnd}(p^{\sharp}M_{k,\ell}) =
  \mathrm{scalarEnd}((p^{\sharp}M)_{k,\ell})\). Reassembling over the biproduct, the conjugating
  one-element comparisons combine into the single free comparison \(Q\), giving
  \(p^{*}(\mathrm{matrixEnd}(M)) = Q^{-1}\circ\mathrm{matrixEnd}(p^{\sharp}M)\circ Q\).
\end{proof}

\begin{lemma}[Affine global-sections comorphism is the inducing ring homomorphism]
  \label{lem:gr_baseChange_bridge_gammaSpec}
  \lean{AlgebraicGeometry.Grassmannian.baseChange_bridge_gammaSpec}
  \uses{def:gr_glued_scheme}
  Let \(\varphi : A \to B\) be a homomorphism of commutative rings and
  \(\Spec\varphi : \Spec B \to \Spec A\) the induced morphism of affine schemes. Under the
  natural identification of the global sections of an affine scheme with its coordinate ring
  -- the counit \(\Gamma(\Spec A, \mathcal O) \xrightarrow{\sim} A\) of the
  \(\Gamma \dashv \Spec\) adjunction, and likewise for \(B\) -- the global-sections base-change
  comorphism
  \[
    (\Spec\varphi)^{\sharp} : \Gamma(\Spec A, \mathcal O) \longrightarrow \Gamma(\Spec B, \mathcal O)
  \]
  is identified with \(\varphi\) itself; that is, the square relating
  \((\Spec\varphi)^{\sharp}\) to \(\varphi\) through the two counit isomorphisms commutes.
\end{lemma}
\begin{proof}
  This is the naturality of the counit isomorphism \(\Gamma \circ \Spec \cong \mathrm{id}\) of
  the global-sections / spectrum adjunction: for any ring homomorphism \(\varphi\) the comorphism
  of \(\Spec\varphi\) on global sections is, after transport through the natural counit
  isomorphisms at \(A\) and \(B\), the homomorphism \(\varphi\). No Grassmannian-specific data
  enters; the statement is the affine-scheme naturality square applied to a single ring map.
\end{proof}

\begin{lemma}[First-projection bridge to \texorpdfstring{$\mathrm{awayInclLeft}$}{awayInclLeft}]
  \label{lem:gr_baseChange_bridge_left}
  \lean{AlgebraicGeometry.Grassmannian.baseChange_bridge_left}
  \uses{def:gr_the_glue_data, def:gr_chart_incl, def:gr_away_pullback_iso,
    lem:gr_awayPullbackIso_inv_fst, def:gr_away_incl_left,
    lem:gr_baseChange_bridge_gammaSpec}
  Fix indices \(I, J, K\) with triple overlap \(V_{IJK} = U^I_J \times_{U^I} U^I_K\) of the glue
  datum (\cref{def:gr_the_glue_data}), and let
  \(p^{IJ}_{IJK} : V_{IJK} \to U^I_J\) be the first projection (the pullback projection of the
  chart inclusion \(\mathrm{chartIncl}\), \cref{def:gr_chart_incl}). Through the away-pullback
  identification \(V_{IJK} \cong \Spec R^I[1/(P^I_J P^I_K)]\) (\cref{def:gr_away_pullback_iso})
  the projection \(p^{IJ}_{IJK}\) is the \(\Spec\) morphism of the left away-inclusion
  \(\mathrm{awayInclLeft}(P^I_J, P^I_K)\) (\cref{lem:gr_awayPullbackIso_inv_fst},
  \cref{def:gr_away_incl_left}). Consequently the global-sections base-change map of
  \(p^{IJ}_{IJK}\), transported through the affine global-sections identification, equals the
  ring homomorphism \(\mathrm{awayInclLeft}(P^I_J, P^I_K)\); in particular the entrywise image of
  the Cramer inverse \((X^I_J)^{-1}\) under the geometric base-change agrees with its image under
  \(\mathrm{awayInclLeft}\).
\end{lemma}
\begin{proof}
  \uses{lem:gr_awayPullbackIso_inv_fst, def:gr_away_incl_left,
    lem:gr_baseChange_bridge_gammaSpec}
  By \cref{lem:gr_awayPullbackIso_inv_fst} the first pullback projection, precomposed with the
  away-pullback identification \cref{def:gr_away_pullback_iso}, equals
  \(\Spec(\mathrm{awayInclLeft}(P^I_J, P^I_K))\). The source and overlap charts are affine, so by
  \cref{lem:gr_baseChange_bridge_gammaSpec} the global-sections comorphism of this \(\Spec\)
  morphism is, through the affine global-sections identification, the ring homomorphism
  \(\mathrm{awayInclLeft}(P^I_J, P^I_K)\) (\cref{def:gr_away_incl_left}). Applying the comorphism
  entrywise to \((X^I_J)^{-1}\) gives the asserted agreement.
\end{proof}

\begin{lemma}[Second-projection bridge to \texorpdfstring{$\mathrm{awayInclRight}$}{awayInclRight}]
  \label{lem:gr_baseChange_bridge_right}
  \lean{AlgebraicGeometry.Grassmannian.baseChange_bridge_right}
  \uses{def:gr_the_glue_data, def:gr_chart_incl, def:gr_away_pullback_iso,
    lem:gr_awayPullbackIso_inv_snd, def:gr_away_incl_right,
    lem:gr_baseChange_bridge_gammaSpec}
  With the notation of \cref{lem:gr_baseChange_bridge_left}, let
  \(p^{IK}_{IJK} : V_{IJK} \to U^I_K\) be the second projection (the pullback projection of the
  chart inclusion \(\mathrm{chartIncl}\), \cref{def:gr_chart_incl}). Through the away-pullback
  identification \(V_{IJK} \cong \Spec R^I[1/(P^I_J P^I_K)]\) (\cref{def:gr_away_pullback_iso})
  the projection \(p^{IK}_{IJK}\) is the \(\Spec\) morphism of the right away-inclusion
  \(\mathrm{awayInclRight}(P^I_J, P^I_K)\) (\cref{lem:gr_awayPullbackIso_inv_snd},
  \cref{def:gr_away_incl_right}). Consequently the global-sections base-change map of
  \(p^{IK}_{IJK}\), transported through the affine global-sections identification, equals the
  ring homomorphism \(\mathrm{awayInclRight}(P^I_J, P^I_K)\); in particular the entrywise image of
  the Cramer inverse \((X^I_K)^{-1}\) under the geometric base-change agrees with its image under
  \(\mathrm{awayInclRight}\).
\end{lemma}
\begin{proof}
  \uses{lem:gr_awayPullbackIso_inv_snd, def:gr_away_incl_right,
    lem:gr_baseChange_bridge_gammaSpec}
  Identical to \cref{lem:gr_baseChange_bridge_left} with the second leg in place of the first:
  by \cref{lem:gr_awayPullbackIso_inv_snd} the second pullback projection, precomposed with the
  away-pullback identification \cref{def:gr_away_pullback_iso}, equals
  \(\Spec(\mathrm{awayInclRight}(P^I_J, P^I_K))\). The charts being affine,
  \cref{lem:gr_baseChange_bridge_gammaSpec} identifies its global-sections comorphism with
  \(\mathrm{awayInclRight}(P^I_J, P^I_K)\) (\cref{def:gr_away_incl_right}), and the entrywise
  image of \((X^I_K)^{-1}\) follows.
\end{proof}

\begin{lemma}[Triple-transition bridge to \texorpdfstring{$\Theta_{IJ}$}{Theta_IJ}]
  \label{lem:gr_baseChange_bridge_transition}
  \lean{AlgebraicGeometry.Grassmannian.baseChange_bridge_transition}
  \uses{def:gr_the_glue_data, def:gr_chart_transition', def:gr_cocycle_theta_ij,
    def:gr_away_pullback_iso, def:gr_away_mul_comm_equiv,
    lem:gr_baseChange_bridge_gammaSpec}
  With the notation of \cref{lem:gr_baseChange_bridge_left}, let
  \(t'_{IJK} : V_{IJK} \to U^J_K \times_{U^J} U^J_I\) be the triple transition of the glue datum
  (\cref{def:gr_chart_transition'}). Through the source and target away-pullback identifications
  (\cref{def:gr_away_pullback_iso}) and the order-swap isomorphism
  (\cref{def:gr_away_mul_comm_equiv}), \(t'_{IJK}\) is the \(\Spec\) morphism of the localised
  triple-overlap transition \(\Theta_{IJ} : S_J \to S_I\) (\cref{def:gr_cocycle_theta_ij},
  composed with the order-swap ring isomorphism). Consequently the global-sections base-change
  map of \(t'_{IJK}\), transported through the affine global-sections identification, equals the
  ring homomorphism \(\Theta_{IJ}\) (up to the order-swap); in particular the entrywise image of
  the Cramer inverse \((X^J_K)^{-1}\) under the geometric base-change agrees with its image under
  \(\Theta_{IJ}\).
\end{lemma}
\begin{proof}
  \uses{def:gr_chart_transition', def:gr_cocycle_theta_ij, def:gr_away_mul_comm_equiv,
    lem:gr_baseChange_bridge_gammaSpec}
  By \cref{def:gr_chart_transition'} the triple transition \(t'_{IJK}\) is the composite of the
  source away-pullback identification, \(\Spec\Theta_{IJ}\) (\cref{def:gr_cocycle_theta_ij}), the
  \(\Spec\) of the order-swap isomorphism (\cref{def:gr_away_mul_comm_equiv}), and the inverse
  target away-pullback identification. The charts and overlaps being affine, applying
  \cref{lem:gr_baseChange_bridge_gammaSpec} to each \(\Spec\) factor identifies the global-sections
  comorphism of \(t'_{IJK}\), through the affine global-sections identifications and the order-swap,
  with the ring homomorphism \(\Theta_{IJ}\). The entrywise image of \((X^J_K)^{-1}\) under this
  comorphism therefore agrees with its image under \(\Theta_{IJ}\).
\end{proof}

\begin{lemma}[Base-change bridge to the localised cocycle ring homomorphisms]
  \label{lem:gr_baseChange_bridge}
  \lean{AlgebraicGeometry.Grassmannian.baseChange_bridge}
  \uses{def:gr_transition, def:gr_glued_scheme, def:gr_the_glue_data,
    lem:gr_bundleCocycle_matrix, def:gr_cocycle_theta_ij, def:gr_away_incl_left,
    def:gr_away_incl_right, lem:gr_baseChange_bridge_left,
    lem:gr_baseChange_bridge_right, lem:gr_baseChange_bridge_transition}
  Fix indices \(I, J, K\) and the triple overlap \(V_{IJK}\) of the chart cover
  \(\{U^I\}\) of \(\mathrm{Gr}(r,d)\) (\cref{def:gr_glued_scheme}, \cref{def:gr_the_glue_data}).
  The scheme-pullback base-change map on global sections
  \(\Gamma(U^I_J, \mathcal{O}) \to \Gamma(V_{IJK}, \mathcal{O})\) induced by the first projection
  \(p^{IJ}_{IJK}\), and likewise the maps induced by the second projection and by the triple
  transition \(t'_{IJK}\) (\cref{def:gr_transition}), are identified -- through the natural identification of global sections of an affine scheme
  with its coordinate ring and the fact that the projections and triple transition are the
  \(\Spec\) morphisms of the coordinate-ring homomorphisms by which the Grassmannian glue datum
  (\cref{def:gr_the_glue_data}) is assembled -- with the ring homomorphisms
  \(\mathrm{awayInclLeft}\) (\cref{def:gr_away_incl_left}),
  \(\mathrm{awayInclRight}\) (\cref{def:gr_away_incl_right}) and \(\Theta_{IJ}\)
  (\cref{def:gr_cocycle_theta_ij}) over which the matrix cocycle
  \cref{lem:gr_bundleCocycle_matrix} is stated. Consequently the entrywise image of a Cramer
  inverse \((X^I_J)^{-1}\) under the geometric base-change map agrees with its image under the
  corresponding coordinate-ring homomorphism, so \cref{lem:gr_bundleCocycle_matrix} applies
  to the base-changed matrices.
\end{lemma}
\begin{proof}
  \uses{lem:gr_baseChange_bridge_left, lem:gr_baseChange_bridge_right,
    lem:gr_baseChange_bridge_transition}
  The three identifications are exactly the three projection bridges. The first projection is
  handled by \cref{lem:gr_baseChange_bridge_left}, identifying its global-sections base-change
  map with \(\mathrm{awayInclLeft}\); the second projection by
  \cref{lem:gr_baseChange_bridge_right}, identifying it with \(\mathrm{awayInclRight}\); and the
  triple transition \(t'_{IJK}\) by \cref{lem:gr_baseChange_bridge_transition}, identifying it
  with \(\Theta_{IJ}\). Each bridge already records the agreement of the entrywise image of the
  relevant Cramer inverse under the geometric base-change with its image under the named ring
  homomorphism over which \cref{lem:gr_bundleCocycle_matrix} is stated, so assembling the three
  gives the claim.
\end{proof}

\begin{lemma}[Transport and endpoint alignment of the bundle transitions]
  \label{lem:gr_bundleCocycle_transport}
  \lean{AlgebraicGeometry.Grassmannian.bundleTransition_cocycle_transport}
  \uses{def:gr_bundleTransition, lem:gr_bundleCocycle_matrix, lem:gr_matrixToFreeIso_mul,
    lem:gr_matrixEnd_pullback, lem:gr_baseChange_bridge,
    lem:modules_pullback_basechange_transport, lem:gr_glueData_bridges,
    def:modules_pullbackComp, lem:gr_pullbackFreeIso}
  Fix indices \(I, J, K\) with triple overlap \(V_{IJK} = U^I_J \times_{U^I} U^I_K\),
  and let \(\widehat{g}_{IJ}^{\,K}, \widehat{g}_{JK}^{\,I}, \widehat{g}_{IK}^{\,J}\) be
  the base-change transports (\cref{lem:modules_pullback_basechange_transport}) of the
  three bundle transitions to \(V_{IJK}\). Under the free-pullback comparisons
  (\cref{lem:gr_pullbackFreeIso}) and the endpoint bridges
  (\cref{lem:gr_glueData_bridges}), each transported transition is identified with the
  matrix automorphism of the corresponding base-changed Cramer inverse on the common
  free sheaf \(\mathcal{O}_{V_{IJK}}^{\,d}\); granting the matrix identity
  \cref{lem:gr_bundleCocycle_matrix}, the composite satisfies
  \[
    \widehat{g}_{JK}^{\,I} \circ \widehat{g}_{IJ}^{\,K}
      \;=\; \widehat{g}_{IK}^{\,J}
  \]
  as isomorphisms of sheaves of modules on \(V_{IJK}\). This is the substantive
  transport step: the alignment of the three transports' domains and codomains.
\end{lemma}
\begin{proof}
  \uses{def:gr_bundleTransition, lem:gr_bundleCocycle_matrix, lem:gr_matrixToFreeIso_mul,
    lem:gr_matrixEnd_pullback, lem:gr_baseChange_bridge,
    lem:modules_pullback_basechange_transport, lem:gr_glueData_bridges,
    def:modules_pullbackComp, lem:gr_pullbackFreeIso}
  This is pure assembly of the two infrastructure pieces over the endpoint bridges. Each
  bundle transition is, by \cref{def:gr_bundleTransition}, \(\mathrm{matrixToFreeIso}\) of a
  localised Cramer inverse conjugated by the free-pullback comparisons
  \cref{lem:gr_pullbackFreeIso}; its base-change transport
  (\cref{lem:modules_pullback_basechange_transport}) reassociates the iterated pullbacks via
  \cref{def:modules_pullbackComp}, and the three endpoint bridges \cref{lem:gr_glueData_bridges}
  pin the source, middle and target sheaves to a single \(\mathcal{O}_{V_{IJK}}^{\,d}\). By
  \cref{lem:gr_matrixEnd_pullback} each transport then equals \(\mathrm{matrixEnd}\) of the
  base-changed Cramer inverse, and by \cref{lem:gr_baseChange_bridge} those base-changed matrices
  are exactly \((X^I_J)^{-1}, (X^J_K)^{-1}, (X^I_K)^{-1}\) over the common triple-overlap ring of
  \cref{lem:gr_bundleCocycle_matrix}. The linkage \cref{lem:gr_matrixToFreeIso_mul} composes the
  first two to \(\mathrm{matrixEnd}\bigl((X^J_K)^{-1}(X^I_J)^{-1}\bigr)\), and
  by the matrix cocycle identity \cref{lem:gr_bundleCocycle_matrix} the argument equals
  \((X^I_K)^{-1}\), so the composite equals \(\widehat{g}_{IK}^{\,J}\); the bridges make both
  sides isomorphisms between the same pair of sheaves on \(V_{IJK}\).
\end{proof}

\begin{lemma}[Triple-overlap multiplicativity of the bundle transition (C2)]
  \label{lem:gr_bundleCocycle_mul}
  \lean{AlgebraicGeometry.Grassmannian.bundleTransition_cocycle}
  \uses{def:gr_bundleTransition, lem:gr_bundleCocycle_matrix, lem:gr_matrixToFreeIso_mul,
    lem:gr_bundleCocycle_transport, lem:modules_pullback_basechange_transport,
    lem:gr_glueData_bridges, def:modules_pullbackComp}
  Over each triple overlap \(V_{IJK} = U^I_J \times_{U^I} U^I_K\) the base-change transports
  (\cref{lem:modules_pullback_basechange_transport}) of the three bundle transitions satisfy
  \[
    \widehat{g}_{JK}^{\,I} \circ \widehat{g}_{IJ}^{\,K}
      \;=\; \widehat{g}_{IK}^{\,J},
  \]
  the module-level form of the descent condition (C2) of \cref{def:scheme_modules_glue} for
  the bundle data \((g_{I,J})\). This is the \(\mathrm{GL}_d\)-valued multiplicative cocycle
  of \(\mathcal{U}\).
\end{lemma}
\begin{proof}
  \uses{lem:gr_bundleCocycle_matrix, lem:gr_matrixToFreeIso_mul,
    lem:gr_bundleCocycle_transport}
  The statement decomposes into three independent steps. First, the matrix-algebra core
  \cref{lem:gr_bundleCocycle_matrix} establishes the Cramer-inverse cocycle
  \((X^J_K)^{-1}(X^I_J)^{-1} = (X^I_K)^{-1}\) over the triple overlap. Second, the linkage
  \cref{lem:gr_matrixToFreeIso_mul} records that composition of matrix automorphisms realises
  matrix multiplication. Third, the transport step \cref{lem:gr_bundleCocycle_transport}
  carries the three transitions to the common overlap \(V_{IJK}\), aligns their source,
  middle, and target sheaves through the base-change transport and the endpoint bridges, and
  -- feeding it the matrix identity through the linkage -- yields exactly the stated equality
  \(\widehat{g}_{JK}^{\,I} \circ \widehat{g}_{IJ}^{\,K} = \widehat{g}_{IK}^{\,J}\).
\end{proof}

\begin{definition}

[The universal quotient sheaf \(\mathcal{U}\)]
  \label{def:gr_universal_quotient_sheaf}
  \lean{AlgebraicGeometry.Grassmannian.universalQuotient}
  \uses{def:gr_chart_quotient, def:gr_universalMinorInv, def:gr_transition,
    lem:gr_cocycle, def:gr_glued_scheme, def:scheme_modules_glue,
    def:gr_modules_glueRestrictionIso, def:gr_bundleTransition,
    lem:gr_bundleCocycle_id, lem:gr_bundleCocycle_mul,
    lem:gr_glueData_bridges, def:is_locally_free_of_rank}
  % SOURCE: [Nitsure], "Construction of Hilbert and Quot Schemes", \S 1,
  % sub-section "Universal quotient" (read from
  % references/nitsure-hilbert-quot-src/nitsure-hilbert-quot.tex, L903-L910).
  % SOURCE QUOTE:
  % "Compatible with the co-cycle $(\theta_{I,J})$
  %  for gluing the affine pieces $U^I$, we give gluing data
  %  $(g_{I,J})$ for gluing together the trivial bundles $\oplus^d\,{\mathcal O}_{U^I}$
  %  by putting
  %  $$g_{I,J} = (X^I_J)^{-1} \in GL_d(U^I_J)$$"
  \textit{Source: [Nitsure], \S 1.}
  Over each overlap \(U^I_J = \Spec R^I_J\) define the bundle-transition matrix
  \[
    g_{I,J} \;=\; (X^I_J)^{-1} \;\in\; \mathrm{GL}_d(R^I_J),
  \]
  the Cramer inverse of the localised \(J\)-minor (\cref{def:gr_universalMinorInv}),
  which is invertible because \(\det(X^I_J)\) is a unit of \(R^I_J\). The matrices
  \((g_{I,J})\) are compatible with the chart-gluing cocycle \((\theta_{I,J})\)
  (\cref{def:gr_transition}, \cref{lem:gr_cocycle}): over a triple overlap they
  satisfy \(g_{I,K} = g_{J,K}\, g_{I,J}\) and \(g_{I,I} = 1\), the multiplicative
  cocycle condition for a vector bundle. Gluing the free rank-\(d\) sheaves
  \(\mathcal{O}_{U^I}^{\,d}\) along the isomorphisms \(g_{I,J}\) over the cover
  \(\{U^I\}\) of \(\mathrm{Gr}(r,d)\) (\cref{def:gr_glued_scheme}) therefore
  produces a locally free \(\mathcal{O}_{\mathrm{Gr}(r,d)}\)-module
  \(\mathcal{U}\) of rank \(d\). Its determinant line bundle \(\det \mathcal{U}\)
  has transition functions \(\det(g_{I,J}) = 1/P^I_J\), the data underlying the
  Pl\"ucker embedding.
\end{definition}

The tautological quotient \(u\) is assembled from the chart quotients \(u^I\) by the
section-level descent \cref{def:gr_modules_glueHom}; its one remaining content is the
\emph{overlap-compatibility} of the family \((u^I)\) with the bundle gluing
\((g_{I,J})\). The next four blocks isolate that compatibility: a rectangular analogue of
the square matrix-endomorphism API (\cref{def:gr_matrixEnd},
\cref{lem:gr_matrixEnd_pullback}, \cref{lem:gr_matrixEnd_comp}) for the
\(r \to d\) chart quotient (\cref{def:gr_matrixEndRect},
\cref{lem:gr_matrixEndRect_pullback}, \cref{lem:gr_matrixEndRect_comp}), the
adjunction transposition that turns the descent condition into a pullback-level identity
(\cref{lem:gr_tautologicalQuotientComponent_transpose}), and the matrix-core assembly
that closes it (\cref{lem:gr_tautologicalQuotient_overlap}).

\subsection{Matrix-calculus toolbox}

The chart-cover lemma (\cref{lem:chartLocus_isOpenCover}) and the overlap compatibility of
\cref{def:grPointOfRankQuotient} both rest on a small calculus of \emph{rectangular} matrix
homomorphisms of free sheaves, generalising the square API
(\cref{def:gr_matrixEnd}, \cref{lem:gr_matrixEnd_comp}, \cref{lem:gr_matrixEnd_one}). The
blocks below collect that toolbox: the bridge to the square case, the identity and
composition laws, injectivity of the matrix presentation, the column-restriction law, the
affine splitting of epimorphisms of free sheaves and the resulting matrix right inverse, the
field-level minor extraction, and the presented-matrix machinery that packages a quotient as
a matrix along a morphism and proves Nitsure's change-of-basis identity. They are
project-bespoke infrastructure.

\begin{lemma}[Square and rectangular matrix homomorphisms coincide]
  \label{lem:gr_matrixEnd_eq_matrixEndRect}
  \lean{AlgebraicGeometry.Grassmannian.matrixEnd_eq_matrixEndRect}
  \uses{def:gr_matrixEnd, def:gr_matrixEndRect}
  For a square matrix \(M \in \operatorname{Mat}_{d\times d}(\Gamma(S,\mathcal{O}_S))\), the
  square and rectangular matrix homomorphisms agree:
  \(\mathrm{matrixEnd}(M) = \mathrm{matrixEndRect}(M)\).
\end{lemma}
\begin{proof}
  \uses{def:gr_matrixEnd, def:gr_matrixEndRect}
  Both are the morphism assembled from the same biproduct components
  \(\mathrm{scalarEnd}(M_{k,\ell})\) over the index set \(\{1,\dots,d\}\); the two
  constructions are literally the same, so the equality holds by definition.
\end{proof}

\begin{lemma}[Identity law for \(\mathrm{matrixEndRect}\)]
  \label{lem:gr_matrixEndRect_one}
  \lean{AlgebraicGeometry.Grassmannian.matrixEndRect_one}
  \uses{def:gr_matrixEndRect, lem:gr_matrixEnd_one, lem:gr_matrixEnd_eq_matrixEndRect}
  The rectangular homomorphism of the identity matrix is the identity morphism:
  \(\mathrm{matrixEndRect}(1) = \mathrm{id}_{\mathcal{O}_S^{\,d}}\).
\end{lemma}
\begin{proof}
  \uses{lem:gr_matrixEnd_one, lem:gr_matrixEnd_eq_matrixEndRect}
  By \cref{lem:gr_matrixEnd_eq_matrixEndRect} the left side equals \(\mathrm{matrixEnd}(1)\),
  which is the identity by \cref{lem:gr_matrixEnd_one}.
\end{proof}

\begin{lemma}[Fully rectangular composition law]
  \label{lem:gr_matrixEndRect_comp_rect}
  \lean{AlgebraicGeometry.Grassmannian.matrixEndRect_comp_rect}
  \uses{def:gr_matrixEndRect, lem:gr_scalarEnd_comp, lem:gr_scalarEnd_sum,
    lem:gr_scalarEnd_add, lem:gr_matrixEndRect_comp}
  For \(A \in \operatorname{Mat}_{e\times n}(\Gamma(S,\mathcal{O}_S))\) and
  \(B \in \operatorname{Mat}_{d\times e}(\Gamma(S,\mathcal{O}_S))\), composing two
  rectangular homomorphisms realises the matrix product, the order reversed by the
  contravariant component indexing exactly as in \cref{lem:gr_matrixEndRect_comp}: the
  composite
  \[
    \mathcal{O}_S^{\,n} \xrightarrow{\mathrm{matrixEndRect}(A)} \mathcal{O}_S^{\,e}
      \xrightarrow{\mathrm{matrixEndRect}(B)} \mathcal{O}_S^{\,d}
  \]
  equals \(\mathrm{matrixEndRect}(B \cdot A)\) (the \(d \times n\) product).
\end{lemma}
\begin{proof}
  \uses{lem:gr_scalarEnd_comp, lem:gr_scalarEnd_sum, lem:gr_scalarEnd_add}
  The same biproduct-extensionality computation as \cref{lem:gr_matrixEndRect_comp}, with the
  middle biproduct indexed by \(\{1,\dots,e\}\). The \((k,\ell)\)-component of the composite
  is \(\sum_m \mathrm{scalarEnd}(B_{k,m})\circ\mathrm{scalarEnd}(A_{m,\ell})\); each composite
  is the scalar endomorphism of the product (\cref{lem:gr_scalarEnd_comp}) and a finite sum
  of scalar endomorphisms is the scalar endomorphism of the sum (\cref{lem:gr_scalarEnd_sum},
  \cref{lem:gr_scalarEnd_add}), so the component is
  \(\mathrm{scalarEnd}((B\cdot A)_{k,\ell})\), the \((k,\ell)\)-component of
  \(\mathrm{matrixEndRect}(B\cdot A)\).
\end{proof}

\begin{lemma}[Uniqueness of the presenting matrix]
  \label{lem:gr_matrixEndRect_injective}
  \lean{AlgebraicGeometry.Grassmannian.matrixEndRect_injective}
  \uses{def:gr_matrixEndRect, def:gr_scalarEnd}
  The matrix presentation is injective: if
  \(\mathrm{matrixEndRect}(M) = \mathrm{matrixEndRect}(N)\) for
  \(M, N \in \operatorname{Mat}_{d\times r}(\Gamma(S,\mathcal{O}_S))\), then \(M = N\).
\end{lemma}
\begin{proof}
  \uses{def:gr_scalarEnd}
  Extract the \((p,j)\)-entry by pre-composing with the free-basis inclusion of the
  \(j\)-th summand and post-composing with the projection onto the \(p\)-th summand: this
  isolates the scalar endomorphism \(\mathrm{scalarEnd}(M_{p,j})\) from
  \(\mathrm{matrixEndRect}(M)\). Equal homomorphisms therefore give
  \(\mathrm{scalarEnd}(M_{p,j}) = \mathrm{scalarEnd}(N_{p,j})\); and \(a \mapsto
  \mathrm{scalarEnd}(a)\) is injective (its inverse reads off the underlying section), whence
  \(M_{p,j} = N_{p,j}\) for all \(p, j\).
\end{proof}

\begin{lemma}[Column restriction against the index inclusion]
  \label{lem:gr_freeMap_matrixEndRect}
  \lean{AlgebraicGeometry.Grassmannian.freeMap_matrixEndRect}
  \uses{def:gr_matrixEndRect}
  For an index map \(g : \{1,\dots,e\} \to \{1,\dots,r\}\) and
  \(M \in \operatorname{Mat}_{d\times r}(\Gamma(S,\mathcal{O}_S))\), pre-composing the
  rectangular homomorphism with the free index inclusion
  \(\mathrm{freeMap}(g) : \mathcal{O}_S^{\,e} \to \mathcal{O}_S^{\,r}\) restricts to the
  \(g\)-columns:
  \[
    \mathrm{freeMap}(g) \,\circ\, \mathrm{matrixEndRect}(M)
      \;=\; \mathrm{matrixEndRect}\bigl(M_{\bullet,\, g}\bigr),
  \]
  where \(M_{\bullet,\, g}\) is the \(d \times e\) submatrix with \(\ell\)-th column the
  \(g(\ell)\)-th column of \(M\).
\end{lemma}
\begin{proof}
  \uses{def:gr_matrixEndRect}
  Both sides agree on each free generator \(j\): the inclusion sends it to the \(g(j)\)-th
  generator, and reading off the \((p,j)\)-entry on either side gives \(M_{p, g(j)}\). Cofan
  extensionality over the \(e\) generators concludes.
\end{proof}

\begin{lemma}[Epimorphisms of free sheaves split over an affine]
  \label{lem:gr_exists_section_of_epi_free_spec}
  \lean{AlgebraicGeometry.Grassmannian.exists_section_of_epi_free_spec}
  On an affine base \(\Spec R\), every epimorphism
  \(\psi : \mathcal{O}_{\Spec R}^{\,r} \twoheadrightarrow \mathcal{O}_{\Spec R}^{\,d}\) of
  free sheaves splits: there is a section
  \(\Phi : \mathcal{O}_{\Spec R}^{\,d} \to \mathcal{O}_{\Spec R}^{\,r}\) with
  \(\psi \circ \Phi = \mathrm{id}\).
\end{lemma}
\begin{proof}
  Under the fully faithful additive functor \(\mathrm{Mod}_R \to (\Spec R)\text{-Mod}\)
  taking a module to its associated quasi-coherent sheaf, the free sheaves correspond to the
  free modules \(R^{\,r}, R^{\,d}\) and \(\psi\) to a surjection of modules. A free module is
  projective, so this surjection admits a module-level section; transporting that section
  back through full faithfulness produces the required \(\Phi\).
\end{proof}

\begin{lemma}[Matrix right inverse over \(\Spec R\)]
  \label{lem:gr_exists_rightInverse_of_epi_matrixEndRect_spec}
  \lean{AlgebraicGeometry.Grassmannian.exists_rightInverse_of_epi_matrixEndRect_spec}
  \uses{lem:gr_exists_section_of_epi_free_spec, lem:gr_matrixEndRect_comp_rect,
    lem:gr_matrixEndRect_injective, lem:gr_matrixEndRect_one}
  For \(M \in \operatorname{Mat}_{d\times r}(\Gamma(\Spec R, \mathcal{O}))\) with
  \(\mathrm{matrixEndRect}(M)\) an epimorphism, the matrix admits a right inverse: there is
  \(G \in \operatorname{Mat}_{r\times d}(\Gamma(\Spec R, \mathcal{O}))\) with \(M \cdot G = 1\).
\end{lemma}
\begin{proof}
  \uses{lem:gr_exists_section_of_epi_free_spec, lem:gr_matrixEndRect_comp_rect,
    lem:gr_matrixEndRect_injective, lem:gr_matrixEndRect_one}
  Split the epimorphism by \cref{lem:gr_exists_section_of_epi_free_spec} and read its section
  off as a matrix \(G\), so that \(\mathrm{matrixEndRect}(G)\) is the section. The composition
  law (\cref{lem:gr_matrixEndRect_comp_rect}) gives \(\mathrm{matrixEndRect}(M\cdot G)\) as the
  composite of section and epimorphism, which is the identity; with
  \(\mathrm{matrixEndRect}(1) = \mathrm{id}\) (\cref{lem:gr_matrixEndRect_one}) and injectivity
  of the presentation (\cref{lem:gr_matrixEndRect_injective}) this forces \(M \cdot G = 1\).
\end{proof}

\begin{lemma}[Matrix right inverse over an affine scheme]
  \label{lem:gr_exists_rightInverse_of_epi_matrixEndRect}
  \lean{AlgebraicGeometry.Grassmannian.exists_rightInverse_of_epi_matrixEndRect}
  \uses{lem:gr_exists_rightInverse_of_epi_matrixEndRect_spec, lem:gr_matrixEndRect_pullback}
  For an affine scheme \(S\) and
  \(M \in \operatorname{Mat}_{d\times r}(\Gamma(S, \mathcal{O}_S))\) with
  \(\mathrm{matrixEndRect}(M)\) an epimorphism, there is
  \(G \in \operatorname{Mat}_{r\times d}(\Gamma(S, \mathcal{O}_S))\) with \(M \cdot G = 1\).
\end{lemma}
\begin{proof}
  \uses{lem:gr_exists_rightInverse_of_epi_matrixEndRect_spec, lem:gr_matrixEndRect_pullback}
  Transport along the canonical isomorphism \(S \cong \Spec\Gamma(S, \mathcal{O}_S)\). The
  base-change naturality of \(\mathrm{matrixEndRect}\) (\cref{lem:gr_matrixEndRect_pullback})
  identifies the pulled-back homomorphism, through the free-pullback comparisons, with
  \(\mathrm{matrixEndRect}\) of the transported matrix, reducing the claim to the \(\Spec\)
  case (\cref{lem:gr_exists_rightInverse_of_epi_matrixEndRect_spec}); pulling the resulting
  right inverse back along the (invertible) global-sections map gives \(G\) over \(S\).
\end{proof}

\begin{lemma}[Field-level minor extraction]
  \label{lem:gr_exists_isUnit_submatrix}
  \lean{AlgebraicGeometry.Grassmannian.exists_isUnit_submatrix}
  Let \(K\) be a field, \(M \in \operatorname{Mat}_{d\times r}(K)\) and
  \(G \in \operatorname{Mat}_{r\times d}(K)\) with \(M \cdot G = 1\). Then some size-\(d\)
  column-index set \(I \subseteq \{1,\dots,r\}\) gives an invertible minor: the \(d \times d\)
  submatrix \(M_I\) of the \(I\)-columns (in their natural order) is a unit.
\end{lemma}
\begin{proof}
  The right inverse makes the column map \(K^{\,r} \to K^{\,d}\) surjective, so the columns of
  \(M\) span \(K^{\,d}\). A spanning family contains a basis; choosing \(d\) columns forming a
  basis, the corresponding \(d \times d\) submatrix has linearly independent columns and is
  therefore invertible.
\end{proof}

\begin{lemma}[Inverse comparison intertwines the index inclusion]
  \label{lem:gr_pullbackFreeIso_inv_freeMap}
  \lean{AlgebraicGeometry.Grassmannian.pullbackFreeIso_inv_freeMap}
  \uses{lem:gr_pullbackFreeIso}
  For a morphism \(p : W \to V\) and an index map \(g : \{1,\dots,n\} \to \{1,\dots,m\}\), the
  inverse free-pullback comparisons \(Q_n, Q_m\) (\cref{lem:gr_pullbackFreeIso}) intertwine the
  index inclusions:
  \[
    Q_n^{-1} \,\circ\, p^{*}\bigl(\mathrm{freeMap}(g)\bigr)
      \;=\; \mathrm{freeMap}(g) \,\circ\, Q_m^{-1}.
  \]
\end{lemma}
\begin{proof}
  \uses{lem:gr_pullbackFreeIso}
  This is the inverse form of the forward naturality of \(\mathrm{freeMap}\) under pullback:
  cancel the comparison isomorphisms \(Q_n, Q_m\) on both sides and apply the forward
  intertwining identity \(p^{*}(\mathrm{freeMap}(g))\circ Q_m = Q_n \circ \mathrm{freeMap}(g)\).
\end{proof}

\begin{definition}[Presented matrix of a quotient along a morphism]
  \label{def:gr_presentedMatrix}
  \lean{AlgebraicGeometry.Grassmannian.presentedMatrix}
  \uses{def:gr_rankQuotient, lem:gr_pullbackFreeIso, def:gr_matrixEndRect}
  Let \(x = \langle\mathcal{F}, q\rangle\) be a rank-\(d\) quotient of \(\mathcal{O}_T^{\,r}\)
  (\cref{def:gr_rankQuotient}), \(j : P \to T\) a morphism, and \(I\) a size-\(d\) subset such
  that the pulled-back chart composite \(j^{*}(s_I \circ q)\) is an isomorphism. The
  \emph{presented matrix} \(\mathrm{presentedMatrix}(x, j, I) \in
  \operatorname{Mat}_{d\times r}(\Gamma(P, \mathcal{O}_P))\) is the matrix that reads off,
  entrywise via the section-extraction \(\mathrm{unitEndSection}\), the conjugated free-sheaf
  homomorphism
  \[
    Q_r^{-1} \,\circ\, j^{*}q \,\circ\, \bigl(j^{*}(s_I \circ q)\bigr)^{-1}
      \,\circ\, Q_d
      \;:\; \mathcal{O}_P^{\,r} \longrightarrow \mathcal{O}_P^{\,d},
  \]
  where \(Q_r, Q_d\) are the free-pullback comparisons (\cref{lem:gr_pullbackFreeIso}). It
  generalises the chart matrix (the case where \(j\) is the inclusion of the chart locus
  \(T_I\)) and is the vehicle of the overlap comparison.
\end{definition}

\begin{lemma}[The presented matrix presents]
  \label{lem:gr_matrixEndRect_presentedMatrix}
  \lean{AlgebraicGeometry.Grassmannian.matrixEndRect_presentedMatrix}
  \uses{def:gr_presentedMatrix, def:gr_matrixEndRect}
  With the data of \cref{def:gr_presentedMatrix},
  \[
    \mathrm{matrixEndRect}\bigl(\mathrm{presentedMatrix}(x, j, I)\bigr)
      \;=\; Q_r^{-1} \,\circ\, j^{*}q \,\circ\, \bigl(j^{*}(s_I \circ q)\bigr)^{-1}
      \,\circ\, Q_d.
  \]
\end{lemma}
\begin{proof}
  \uses{def:gr_presentedMatrix}
  Immediate from the definition: \(\mathrm{matrixEndRect}\) of the
  \(\mathrm{unitEndSection}\)-matrix of a free-sheaf homomorphism recovers that homomorphism.
\end{proof}

\begin{lemma}[Minor of the presented matrix presents the \(J\)-chart composite]
  \label{lem:gr_matrixEndRect_presentedMatrix_minor}
  \lean{AlgebraicGeometry.Grassmannian.matrixEndRect_presentedMatrix_minor}
  \uses{lem:gr_freeMap_matrixEndRect, lem:gr_pullbackFreeIso_inv_freeMap,
    lem:gr_matrixEndRect_presentedMatrix}
  For a second size-\(d\) subset \(J\), the \(J\)-minor of the \(I\)-presented matrix presents
  the pulled-back \(J\)-chart composite against the inverted \(I\)-chart composite:
  \[
    \mathrm{matrixEndRect}\bigl(\mathrm{presentedMatrix}(x, j, I)_{\bullet,\, J}\bigr)
      \;=\; Q_d^{-1} \,\circ\, j^{*}(s_J \circ q) \,\circ\,
        \bigl(j^{*}(s_I \circ q)\bigr)^{-1} \,\circ\, Q_d.
  \]
\end{lemma}
\begin{proof}
  \uses{lem:gr_freeMap_matrixEndRect, lem:gr_pullbackFreeIso_inv_freeMap,
    lem:gr_matrixEndRect_presentedMatrix}
  Column restriction (\cref{lem:gr_freeMap_matrixEndRect}) turns the \(J\)-minor into
  pre-composition with the index inclusion \(s_J\); the inverse-comparison intertwiner
  (\cref{lem:gr_pullbackFreeIso_inv_freeMap}) and the presentation identity
  (\cref{lem:gr_matrixEndRect_presentedMatrix}) move that inclusion across the pullback, where
  \(s_J \circ q\) is exactly the \(J\)-chart composite.
\end{proof}

\begin{lemma}[Change of basis between chart presentations]
  \label{lem:gr_presentedMatrix_changeOfBasis}
  \lean{AlgebraicGeometry.Grassmannian.presentedMatrix_changeOfBasis}
  \uses{def:gr_presentedMatrix, lem:gr_matrixEndRect_presentedMatrix,
    lem:gr_matrixEndRect_presentedMatrix_minor, lem:gr_matrixEndRect_comp_rect,
    lem:gr_matrixEndRect_injective}
  Along a morphism \(j\) on which both the \(I\)- and \(J\)-chart composites pull back to
  isomorphisms, the \(I\)-presented matrix is the product of its own \(J\)-minor with the
  \(J\)-presented matrix -- Nitsure's overlap identity \(M^I = M^I_J \cdot M^J\):
  \[
    \mathrm{presentedMatrix}(x, j, I)
      \;=\; \mathrm{presentedMatrix}(x, j, I)_{\bullet,\, J}
        \,\cdot\, \mathrm{presentedMatrix}(x, j, J).
  \]
\end{lemma}
\begin{proof}
  \uses{lem:gr_matrixEndRect_presentedMatrix, lem:gr_matrixEndRect_presentedMatrix_minor,
    lem:gr_matrixEndRect_comp_rect, lem:gr_matrixEndRect_injective}
  Apply \(\mathrm{matrixEndRect}\) and use injectivity (\cref{lem:gr_matrixEndRect_injective}).
  By the composition law (\cref{lem:gr_matrixEndRect_comp_rect}) the right side becomes the
  composite of the two presented homomorphisms; substituting
  \cref{lem:gr_matrixEndRect_presentedMatrix} (for both factors) and
  \cref{lem:gr_matrixEndRect_presentedMatrix_minor} (for the minor), the pulled-back
  \(J\)-chart composite \(j^{*}(s_J \circ q)\) cancels with its inverse, leaving
  \(Q_r^{-1} \circ j^{*}q \circ (j^{*}(s_I \circ q))^{-1} \circ Q_d =
  \mathrm{matrixEndRect}(\mathrm{presentedMatrix}(x, j, I))\).
\end{proof}

\begin{lemma}[Invertible presentation forces a unit matrix]
  \label{lem:gr_isUnit_of_isIso_matrixEndRect}
  \lean{AlgebraicGeometry.Grassmannian.isUnit_of_isIso_matrixEndRect}
  \uses{def:gr_matrixToFreeIso, lem:gr_matrixEndRect_comp_rect,
    lem:gr_matrixEndRect_injective, lem:gr_matrixEndRect_one}
  Converse of \cref{def:gr_matrixToFreeIso}: if \(\mathrm{matrixEndRect}(N)\) is an
  isomorphism for a square \(N \in \operatorname{Mat}_{d\times d}(\Gamma(S, \mathcal{O}_S))\),
  then \(N\) is a unit (an invertible matrix).
\end{lemma}
\begin{proof}
  \uses{lem:gr_matrixEndRect_comp_rect, lem:gr_matrixEndRect_injective,
    lem:gr_matrixEndRect_one}
  Read off the inverse morphism by its \(\mathrm{unitEndSection}\) matrix \(N'\), so that
  \(\mathrm{matrixEndRect}(N') = \mathrm{matrixEndRect}(N)^{-1}\). The composition law
  (\cref{lem:gr_matrixEndRect_comp_rect}), the identity law
  (\cref{lem:gr_matrixEndRect_one}) and injectivity (\cref{lem:gr_matrixEndRect_injective})
  turn the two-sided inverse relation of morphisms into \(N \cdot N' = 1 = N' \cdot N\),
  exhibiting \(N\) as a unit.
\end{proof}

\begin{definition}

[Rectangular matrix homomorphism of free sheaves]
  \label{def:gr_matrixEndRect}
  \lean{AlgebraicGeometry.Grassmannian.matrixEndRect}
  \uses{def:gr_scalarEnd}
  Let \(S\) be a scheme, \(d, r \in \mathbb{N}\), and
  \(M \in \operatorname{Mat}_{d \times r}(\Gamma(S, \mathcal{O}_S))\) a rectangular matrix
  of global functions. The \emph{rectangular matrix homomorphism}
  \(\mathrm{matrixEndRect}(M)\) is the morphism of free sheaves
  \[
    \mathrm{matrixEndRect}(M) \;:\; \mathcal{O}_S^{\,r}
      \;\longrightarrow\; \mathcal{O}_S^{\,d}
  \]
  whose \((k,\ell)\)-component -- \(k \in \{1,\dots,d\}\) an output index,
  \(\ell \in \{1,\dots,r\}\) an input index -- under the presentation of
  \(\mathcal{O}_S^{\,r}\) and \(\mathcal{O}_S^{\,d}\) as finite biproducts of copies of the
  unit module \(\mathbf 1\), is the scalar endomorphism \(\mathrm{scalarEnd}(M_{k,\ell})\)
  (\cref{def:gr_scalarEnd}). It is assembled from these components by the universal property
  of the two biproducts, exactly as the square \(\mathrm{matrixEnd}\)
  (\cref{def:gr_matrixEnd}) and the chart quotient (\cref{def:gr_chart_quotient}). When
  \(r = d\) it specialises to \(\mathrm{matrixEnd}\); and the chart quotient \(u^I\)
  (\cref{def:gr_chart_quotient}) is by construction \(\mathrm{matrixEndRect}(X^I)\) of the
  universal \(d \times r\) matrix \(X^I\) (\cref{def:gr_universal_matrix}).
\end{definition}

\begin{lemma}

[Rectangular matrix homomorphism is natural under pullback]
  \label{lem:gr_matrixEndRect_pullback}
  \lean{AlgebraicGeometry.Grassmannian.matrixEndRect_pullback}
  \uses{def:gr_matrixEndRect, def:gr_scalarEnd, lem:gr_scalarEnd_pullback,
    lem:gr_pullbackFreeIso, def:modules_pullbackComp, lem:gr_matrixEnd_pullback}
  Let \(p : T \to S\) be a morphism of schemes and
  \(M \in \operatorname{Mat}_{d\times r}(\Gamma(S,\mathcal{O}_S))\), with entrywise
  base-changed matrix \(p^{\sharp}M \in \operatorname{Mat}_{d\times r}(\Gamma(T,\mathcal{O}_T))\).
  Write \(Q_r : p^{*}\mathcal{O}_S^{\,r} \xrightarrow{\sim} \mathcal{O}_T^{\,r}\) and
  \(Q_d : p^{*}\mathcal{O}_S^{\,d} \xrightarrow{\sim} \mathcal{O}_T^{\,d}\) for the
  free-pullback comparisons (\cref{lem:gr_pullbackFreeIso}) at the index sets
  \(\{1,\dots,r\}\) and \(\{1,\dots,d\}\). Then the pullback of
  \(\mathrm{matrixEndRect}(M)\) is conjugate, through \(Q_r\) and \(Q_d\), to the
  rectangular homomorphism of the base-changed matrix:
  \[
    p^{*}\bigl(\mathrm{matrixEndRect}(M)\bigr)
      \;=\; Q_d^{-1} \,\circ\, \mathrm{matrixEndRect}(p^{\sharp}M) \,\circ\, Q_r .
  \]
\end{lemma}
\begin{proof}
  \uses{def:gr_matrixEndRect, lem:gr_scalarEnd_pullback, lem:gr_pullbackFreeIso,
    def:modules_pullbackComp, lem:gr_matrixEnd_pullback}
  Identical to the square case \cref{lem:gr_matrixEnd_pullback}, with the source and target
  biproducts now indexed by \(\{1,\dots,r\}\) and \(\{1,\dots,d\}\) respectively. The
  pullback functor \(p^{*}\), a left adjoint, is additive and biproduct-preserving; through
  the coproduct comparisons underlying the free-pullback isomorphisms
  \cref{lem:gr_pullbackFreeIso} it carries the biproduct presentation of
  \(\mathcal{O}_S^{\,r}\) (resp. \(\mathcal{O}_S^{\,d}\)) to that of \(\mathcal{O}_T^{\,r}\)
  (resp. \(\mathcal{O}_T^{\,d}\)), the reassociation of the iterated pullbacks being
  \cref{def:modules_pullbackComp}. On each single \((k,\ell)\)-component the pulled-back
  scalar endomorphism is, by \cref{lem:gr_scalarEnd_pullback}, the comparison-conjugate of
  \(\mathrm{scalarEnd}((p^{\sharp}M)_{k,\ell})\). Reassembling over the two biproducts, the
  conjugating one-element comparisons combine into \(Q_r\) on the source and \(Q_d\) on the
  target, giving \(p^{*}(\mathrm{matrixEndRect}(M)) = Q_d^{-1}\circ
  \mathrm{matrixEndRect}(p^{\sharp}M)\circ Q_r\).
\end{proof}

\begin{lemma}

[Square-after-rectangular composition law]
  \label{lem:gr_matrixEndRect_comp}
  \lean{AlgebraicGeometry.Grassmannian.matrixEndRect_comp}
  \uses{def:gr_matrixEndRect, def:gr_matrixEnd, lem:gr_scalarEnd_comp, lem:gr_scalarEnd_add,
    lem:gr_matrixEnd_comp}
  For \(M \in \operatorname{Mat}_{d\times r}(\Gamma(S,\mathcal{O}_S))\) and
  \(N \in \operatorname{Mat}_{d\times d}(\Gamma(S,\mathcal{O}_S))\), post-composing the
  rectangular homomorphism with a square one realises the matrix product, with the order
  reversed by the contravariance of the component indexing exactly as in
  \cref{lem:gr_matrixEnd_comp}:
  \[
    \mathrm{matrixEndRect}(M) \,\circ\, \mathrm{matrixEnd}(N)
      \;=\; \mathrm{matrixEndRect}(N \cdot M),
  \]
  the composite \(\mathcal{O}_S^{\,r} \xrightarrow{\mathrm{matrixEndRect}(M)}
  \mathcal{O}_S^{\,d} \xrightarrow{\mathrm{matrixEnd}(N)} \mathcal{O}_S^{\,d}\) (here
  \(N \cdot M\) is the \(d \times r\) product).
\end{lemma}
\begin{proof}
  \uses{def:gr_matrixEndRect, def:gr_matrixEnd, lem:gr_scalarEnd_comp, lem:gr_scalarEnd_add,
    lem:gr_matrixEnd_comp}
  The same biproduct-extensionality computation as \cref{lem:gr_matrixEnd_comp}, the only
  change being that the leftmost (input) biproduct is indexed by \(\{1,\dots,r\}\) rather
  than \(\{1,\dots,d\}\). The composite of two biproduct-assembled morphisms is again
  biproduct-assembled, with \((k,\ell)\)-component the sum over the middle index \(m\) of the
  composites of components, \(\sum_m \mathrm{scalarEnd}(N_{k,m}) \circ
  \mathrm{scalarEnd}(M_{m,\ell})\). Each composite of scalar endomorphisms is the scalar
  endomorphism of the product (\cref{lem:gr_scalarEnd_comp}) and a finite sum of scalar
  endomorphisms is the scalar endomorphism of the sum (\cref{lem:gr_scalarEnd_add}); hence
  that component is \(\mathrm{scalarEnd}(\sum_m N_{k,m} M_{m,\ell}) =
  \mathrm{scalarEnd}((N\cdot M)_{k,\ell})\), the \((k,\ell)\)-component of
  \(\mathrm{matrixEndRect}(N\cdot M)\).
\end{proof}

\begin{lemma}

[Adjunction transpose of the chart-overlap condition]
  \label{lem:gr_tautologicalQuotientComponent_transpose}
  \lean{AlgebraicGeometry.Grassmannian.tautologicalQuotientComponent_transpose}
  \uses{def:gr_chart_quotient, def:gr_bundleTransition, def:gr_modules_glueHom,
    lem:modules_pullback_basechange_transport, lem:gr_homEquiv_conjugateEquiv_app,
    lem:conjugateEquiv_pullbackComp_inv_mathlib, lem:gr_pullbackFreeIso,
    def:modules_pullbackComp}
  Fix size-\(d\) subsets \(I, J\) and let \(f_{IJ} : U^I_J \hookrightarrow U^I\),
  \(t_{IJ} : U^I_J \xrightarrow{\sim} U^J_I\) be the overlap immersion and transition
  (\cref{def:gr_transition}). The per-chart components \(u^I\) of the tautological quotient
  are the adjoint transposes, along the chart immersions, of the chart quotients
  \cref{def:gr_chart_quotient} (precomposed with the free-pullback comparison). Their
  descent compatibility -- the hypothesis the gluing primitive \cref{def:gr_modules_glueHom}
  consumes, namely that the bundle transition \(g_{I,J}\) (\cref{def:gr_bundleTransition})
  intertwines the pullbacks of \(u^I\) and \(u^J\) over \(U^I_J\) -- is equivalent, under the
  pullback--pushforward adjunction of the overlap immersion, to the pullback-level identity
  of morphisms \(\mathcal{O}_{U^I_J}^{\,r} \to \mathcal{O}_{U^I_J}^{\,d}\)
  \[
    g_{I,J} \,\circ\, f_{IJ}^{*}\,u^I
      \;=\; (t_{IJ} \circ f_{JI})^{*}\,u^J,
  \]
  understood through the free-pullback comparisons (\cref{lem:gr_pullbackFreeIso}) on
  source and target.
\end{lemma}
\begin{proof}
  \uses{lem:gr_homEquiv_conjugateEquiv_app, lem:conjugateEquiv_pullbackComp_inv_mathlib,
    lem:gr_pullbackFreeIso, def:modules_pullbackComp}
  The components are produced by applying the adjunction hom-equivalence to the chart
  quotients; the descent condition of \cref{def:gr_modules_glueHom} is stated on those
  transposes. Transposing back across the adjunction is the naturality of the conjugate
  (mate) construction against the hom-equivalences, \cref{lem:gr_homEquiv_conjugateEquiv_app},
  together with the compatibility of the conjugate with the iterated-pullback comparison
  \cref{lem:conjugateEquiv_pullbackComp_inv_mathlib} reassociating the two pullbacks
  \(f_{IJ}^{*}\) and \((t_{IJ}\circ f_{JI})^{*}\) via \cref{def:modules_pullbackComp}. The
  free-pullback comparisons \cref{lem:gr_pullbackFreeIso} re-present both sides as morphisms
  of the trivial bundles \(\mathcal{O}_{U^I_J}^{\,r} \to \mathcal{O}_{U^I_J}^{\,d}\), under
  which the transposed condition becomes the displayed pullback-level identity.
\end{proof}

\begin{lemma}

[Overlap compatibility of the tautological quotient]
  \label{lem:gr_tautologicalQuotient_overlap}
  \lean{AlgebraicGeometry.Grassmannian.tautologicalQuotient_overlap}
  \uses{lem:gr_tautologicalQuotientComponent_transpose, lem:gr_matrixEndRect_pullback,
    lem:gr_matrixEndRect_comp, def:gr_bundleTransition, def:gr_chart_quotient,
    def:gr_universalMinorInv, lem:gr_universalMatrix_map_transitionPreMap,
    lem:gr_imageMatrix_map_eq, lem:modules_pullback_basechange_transport,
    lem:gr_baseChange_bridge_left}
  For every pair of size-\(d\) subsets \(I, J\), the chart quotients \(u^I\)
  (\cref{def:gr_chart_quotient}) satisfy the overlap-compatibility condition required by the
  section-level descent \cref{def:gr_modules_glueHom} for the bundle gluing \((g_{I,J})\):
  the pullback-level identity \(g_{I,J} \circ f_{IJ}^{*} u^I = (t_{IJ}\circ f_{JI})^{*} u^J\)
  of \cref{lem:gr_tautologicalQuotientComponent_transpose} holds on each overlap \(U^I_J\).
\end{lemma}
\begin{proof}
  \uses{lem:gr_tautologicalQuotientComponent_transpose, lem:gr_matrixEndRect_pullback,
    lem:gr_matrixEndRect_comp, def:gr_bundleTransition,
    lem:gr_universalMatrix_map_transitionPreMap, lem:gr_imageMatrix_map_eq,
    lem:gr_baseChange_bridge_left}
  By \cref{lem:gr_tautologicalQuotientComponent_transpose} it suffices to prove the
  pullback-level identity \(g_{I,J} \circ f_{IJ}^{*} u^I = (t_{IJ}\circ f_{JI})^{*} u^J\) on
  \(U^I_J\). Both sides are morphisms of trivial bundles \(\mathcal{O}_{U^I_J}^{\,r} \to
  \mathcal{O}_{U^I_J}^{\,d}\); write each as \(\mathrm{matrixEndRect}\) of a \(d \times r\)
  matrix over the overlap ring \(R^I_J\) and compare matrices.

  Since \(u^I = \mathrm{matrixEndRect}(X^I)\) (\cref{def:gr_chart_quotient}), the
  base-change naturality \cref{lem:gr_matrixEndRect_pullback} identifies \(f_{IJ}^{*} u^I\),
  through the free-pullback comparisons, with \(\mathrm{matrixEndRect}\) of the localised
  matrix \(X^I\) over \(R^I_J\) (the global-sections base change is the away-localisation
  \cref{lem:gr_baseChange_bridge_left}). The bundle transition \(g_{I,J}\) is
  \(\mathrm{matrixEnd}\) of the Cramer inverse \((X^I_J)^{-1}\)
  (\cref{def:gr_bundleTransition}, \cref{def:gr_universalMinorInv}); hence by the
  square-after-rectangular composition law \cref{lem:gr_matrixEndRect_comp},
  \[
    g_{I,J} \circ f_{IJ}^{*} u^I
      \;=\; \mathrm{matrixEndRect}\bigl((X^I_J)^{-1} \cdot X^I\bigr).
  \]
  The transition image matrix identity \(X^J = (X^I_J)^{-1} X^I\)
  (\cref{lem:gr_universalMatrix_map_transitionPreMap}, transported across the overlap by
  \cref{lem:gr_imageMatrix_map_eq}) turns the right-hand matrix into \(X^J\). Symmetrically,
  \(\mathrm{matrixEndRect}_J\) of \(X^J\) is, by \cref{lem:gr_matrixEndRect_pullback} along
  \(t_{IJ}\circ f_{JI}\), the comparison form of \((t_{IJ}\circ f_{JI})^{*} u^J\). Therefore
  the two sides agree, which is the required overlap compatibility.
\end{proof}

\begin{definition}

[The tautological quotient \(u : \mathcal{O}^{r} \twoheadrightarrow \mathcal{U}\)]
  \label{def:tautological_quotient}
  \lean{AlgebraicGeometry.Grassmannian.tautologicalQuotient}
  \uses{def:gr_universal_quotient_sheaf, def:gr_chart_quotient,
    def:gr_glued_scheme, def:gr_universalMinorInv, lem:gr_cocycle,
    def:scheme_modules_glue, def:gr_modules_glueRestrictionIso,
    def:gr_modules_glueHom, lem:gr_tautologicalQuotient_overlap}
  % SOURCE: [Nitsure], "Construction of Hilbert and Quot Schemes", \S 1,
  % sub-section "Universal quotient" (read from
  % references/nitsure-hilbert-quot-src/nitsure-hilbert-quot.tex, L908-L910).
  % SOURCE QUOTE:
  % "This is compatible with the homomorphisms $u^I$,
  %  so we get a surjective homomorphism
  %  $u : \oplus^r\,{\mathcal O}_{\grass(r,d)} \to \U$."
  \textit{Source: [Nitsure], \S 1.}
  The chart quotients \(u^I : \mathcal{O}_{U^I}^{\,r} \to \mathcal{O}_{U^I}^{\,d}\)
  (\cref{def:gr_chart_quotient}) are compatible with the bundle gluing
  \((g_{I,J})\) of \(\mathcal{U}\) (\cref{def:gr_universal_quotient_sheaf}): over
  each overlap \(U^I_J\) one has the identity of \(d \times r\) matrices
  \[
    g_{I,J}\, X^I \;=\; (X^I_J)^{-1} X^I \;=\; X^J \qquad \text{on } U^I_J,
  \]
  which is precisely the defining formula of the chart transition
  \(\theta_{I,J}\) (\cref{def:gr_transition}, \cref{def:gr_universalMinorInv}). At the
  module level this overlap compatibility -- the equalising condition consumed by the
  section-level descent \cref{def:gr_modules_glueHom} -- is exactly
  \cref{lem:gr_tautologicalQuotient_overlap}, whose proof routes the matrix identity above
  through the rectangular pullback/composition API
  (\cref{lem:gr_matrixEndRect_pullback}, \cref{lem:gr_matrixEndRect_comp}) and the
  adjunction transpose \cref{lem:gr_tautologicalQuotientComponent_transpose}.
  Hence the local maps \(u^I\), transported through the trivialisations of
  \(\mathcal{U}\), agree on overlaps and glue to a single surjective
  \(\mathcal{O}_{\mathrm{Gr}(r,d)}\)-linear homomorphism
  \[
    u \;:\; \mathcal{O}_{\mathrm{Gr}(r,d)}^{\,r}
      \;\twoheadrightarrow\; \mathcal{U},
  \]
  the \emph{tautological} (or \emph{universal}) rank-\(d\) quotient of the trivial
  rank-\(r\) sheaf. It is surjective because each \(u^I\) is (split) surjective and
  surjectivity is local on \(\mathrm{Gr}(r,d)\), and \(\mathcal{U}\) is locally
  free of rank \(d\) by \cref{def:gr_universal_quotient_sheaf}. The pair
  \(\langle \mathcal{U}, u \rangle\) is the family of rank-\(d\) quotients carried
  by \(\mathrm{Gr}(r,d)\) itself.
\end{definition}

\section{The functor of points}
\label{sec:grquot_functor}

The Grassmannian functor sends a scheme \(T\) to the set of \emph{equivalence classes} of
rank-\(d\) quotients of \(\mathcal{O}_T^{\,r}\). The following block packages a single such
quotient, the equivalence relation on them, and the pullback action through which the
functor acts on morphisms; the two coherence blocks after it isolate the comparison
identities that make that action functorial.

\begin{definition}[Rank-\(d\) quotients and their pullback action]
  \label{def:gr_rankQuotient}
  \lean{AlgebraicGeometry.Grassmannian.RankQuotient,
        AlgebraicGeometry.Grassmannian.RankQuotient.Rel,
        AlgebraicGeometry.Grassmannian.RankQuotient.rel_refl,
        AlgebraicGeometry.Grassmannian.RankQuotient.rel_symm,
        AlgebraicGeometry.Grassmannian.RankQuotient.rel_trans,
        AlgebraicGeometry.Grassmannian.rqSetoid,
        AlgebraicGeometry.Grassmannian.rqPullback,
        AlgebraicGeometry.Grassmannian.rqPullback_rel}
  \uses{lem:gr_pullback_isLocallyFreeOfRank, lem:gr_pullbackFreeIso,
        def:is_locally_free_of_rank}
  A \emph{rank-\(d\) quotient of \(\mathcal{O}_T^{\,r}\)} on a scheme \(T\) is a triple
  \(x = \langle \mathcal{F}, q \rangle\) consisting of a sheaf of \(\mathcal{O}_T\)-modules
  \(\mathcal{F}\) that is locally free of rank \(d\) (\cref{def:is_locally_free_of_rank})
  together with an epimorphism \(q : \mathcal{O}_T^{\,r} \twoheadrightarrow \mathcal{F}\) out
  of the trivial rank-\(r\) bundle. Two such quotients \(x = \langle \mathcal{F}, q\rangle\)
  and \(y = \langle \mathcal{F}', q'\rangle\) are \emph{equivalent}, written \(x \sim y\),
  when there is an isomorphism \(f : \mathcal{F} \xrightarrow{\sim} \mathcal{F}'\) with
  \(q \circ f = q'\) (equivalently \(\ker q = \ker q'\)). This relation is reflexive
  (\(f = \mathrm{id}\)), symmetric (\(f \mapsto f^{-1}\)), and transitive (composition of the
  witnessing isomorphisms), hence an equivalence relation -- the setoid whose quotient is the
  functor's value at \(T\).

  For a morphism \(\psi : T' \to T\) the \emph{pullback action} sends
  \(x = \langle \mathcal{F}, q\rangle\) to
  \[
    \psi^{*}x \;=\;
      \bigl\langle\, \psi^{*}\mathcal{F},\;
        \mathcal{O}_{T'}^{\,r} \xrightarrow{\ \sim\ }
        \psi^{*}\bigl(\mathcal{O}_T^{\,r}\bigr) \xrightarrow{\ \psi^{*}q\ }
        \psi^{*}\mathcal{F} \,\bigr\rangle,
  \]
  re-presenting the source as the trivial bundle via the free-pullback isomorphism
  \cref{lem:gr_pullbackFreeIso}. The result is again a rank-\(d\) quotient: the composite is
  an epimorphism because pullback, a left adjoint, preserves epimorphisms and the
  re-presentation is an isomorphism, and \(\psi^{*}\mathcal{F}\) is locally free of rank \(d\)
  by \cref{lem:gr_pullback_isLocallyFreeOfRank}. The pullback action respects the equivalence
  relation -- given \(f\) witnessing \(x \sim y\), the pulled-back isomorphism \(\psi^{*}f\)
  witnesses \(\psi^{*}x \sim \psi^{*}y\), by functoriality of pullback -- so it descends to a
  map of equivalence classes.
\end{definition}

The two blocks below record the comparison identities, between the free-pullback isomorphism
\cref{lem:gr_pullbackFreeIso} and Mathlib's pseudofunctor structure on pullback, that make
the descended pullback action of \cref{def:gr_rankQuotient} functorial.

\begin{lemma}\leanok
[Identity coherence of the free-pullback comparison]
  \label{lem:gr_pullbackObjUnitToUnit_id}
  \lean{AlgebraicGeometry.Scheme.Modules.pullbackObjUnitToUnit_id}
  \uses{lem:gr_pullbackFreeIso}
  For the identity morphism \(\mathrm{id}_T\), the canonical comparison morphism of unit
  modules \(\mathrm{id}_T^{*}\mathbf{1} \to \mathbf{1}\) -- the map underlying the
  free-pullback isomorphism \cref{lem:gr_pullbackFreeIso} at a one-element index -- coincides
  with the component at \(\mathbf{1}\) of the pullback-identity isomorphism
  \(\mathrm{id}_T^{*} \cong \mathrm{id}\).
\end{lemma}

\begin{proof}
  \uses{lem:gr_pullbackFreeIso}
  \leanok
  The comparison morphism is the mate of \(\mathrm{id}_{\mathbf 1}\) under the
  pullback--pushforward adjunction: the image of the identity of \(\mathbf 1\) through the
  adjunction's hom-equivalence on the unit. For the identity scheme morphism the pushforward
  is the identity functor, so the adjunction unit is itself an identity; Mathlib's
  pullback-identity isomorphism is exactly the conjugate (left-adjoint transpose) of the
  pushforward-identity isomorphism. Tracing the identity through the hom-equivalence and the
  conjugate-transpose identity for the pullback-identity isomorphism shows the comparison
  equals the \(\mathbf 1\)-component of that isomorphism.
\end{proof}

\begin{lemma}\leanok
[Identity coherence at the free sheaf]
  \label{lem:gr_pullbackFreeIso_id}
  \lean{AlgebraicGeometry.Scheme.Modules.pullbackFreeIso_id}
  \uses{lem:gr_pullbackFreeIso, lem:gr_pullbackObjUnitToUnit_id}
  For the identity morphism \(\mathrm{id}_T\) and an index set \(I\), the free-pullback
  isomorphism \(\mathrm{id}_T^{*}(\bigoplus_I \mathbf 1) \cong \bigoplus_I \mathbf 1\) coincides
  with the component at the free sheaf \(\bigoplus_I \mathbf 1\) of the pullback-identity
  isomorphism \(\mathrm{id}_T^{*} \cong \mathrm{id}\).
\end{lemma}

\begin{proof}
  \uses{lem:gr_pullbackObjUnitToUnit_id}
  \leanok
  Both sides are cocones on the coproduct \(\bigoplus_I \mathbf 1\); by coproduct
  extensionality it suffices to check equality after precomposing with each coprojection
  \(\mathbf 1 \to \bigoplus_I \mathbf 1\), which reduces to the one-element case
  \cref{lem:gr_pullbackObjUnitToUnit_id}.
\end{proof}

\begin{lemma}\leanok
[Composite coherence of the free-pullback comparison]
  \label{lem:gr_pullbackObjUnitToUnit_comp}
  \lean{AlgebraicGeometry.Scheme.Modules.pullbackObjUnitToUnit_comp}
  \uses{lem:gr_pullbackObjUnitToUnit_id, def:modules_pullbackComp}
  For composable scheme morphisms \(a : T'' \to T'\) and \(b : T' \to T\), the comparison
  morphism of unit modules \((b \circ a)^{*}\mathbf{1} \to \mathbf{1}\) factors as
  \(a^{*}\) of the comparison for \(b\), the comparison for \(a\), and the component at
  \(\mathbf{1}\) of the iterated-pullback comparison \(a^{*}b^{*} \cong (b\circ a)^{*}\)
  (\cref{def:modules_pullbackComp}).
\end{lemma}

\begin{proof}
  \uses{lem:gr_pullbackObjUnitToUnit_id, def:modules_pullbackComp,
    lem:gr_homEquiv_conjugateEquiv_app}
    \leanok
  This is the composite analogue of \cref{lem:gr_pullbackObjUnitToUnit_id}: the mate
  construction is compatible with composition of adjunctions, so the comparison for
  \(b \circ a\) is the pasting of the comparisons for \(a\) and \(b\) reassociated by the
  pseudofunctor comparison \cref{def:modules_pullbackComp}. The compatibility of the mate
  with the adjunction hom-equivalences is supplied by
  \cref{lem:gr_homEquiv_conjugateEquiv_app}, the reusable conjugate-naturality engine. The
  pentagon coherence of the iterated-pullback comparisons makes the factorisation
  well-defined.
\end{proof}

\begin{lemma}\leanok
[Mate/conjugate compatibility of the adjunction hom-equivalence]
  \label{lem:gr_homEquiv_conjugateEquiv_app}
  \lean{AlgebraicGeometry.Scheme.Modules.homEquiv_conjugateEquiv_app}
  Let \(\mathrm{adj}_1, \mathrm{adj}_2\) be two adjunctions and \(\alpha\) a natural
  transformation between their left adjoints. For a morphism \(f\) the hom-equivalence of
  \(\mathrm{adj}_2\) applied to \(\alpha_c \circ f\) equals the hom-equivalence of
  \(\mathrm{adj}_1\) applied to \(f\), post-composed with the component at \(d\) of the
  conjugate (mate) of \(\alpha\). This is the general naturality of the conjugate
  construction against the two hom-equivalences; it is the reusable engine behind
  \cref{lem:gr_pullbackObjUnitToUnit_comp}.
\end{lemma}

\begin{proof}


\leanok
  Both sides unfold to the same pasting of unit/counit triangles for the two adjunctions;
  the equality is the standard compatibility of the mate of \(\alpha\) with the
  hom-set bijections.
\end{proof}

\begin{lemma}\leanok
[Composite coherence at the free sheaf]
  \label{lem:gr_pullbackFreeIso_comp}
  \lean{AlgebraicGeometry.Scheme.Modules.pullbackFreeIso_comp}
  \uses{lem:gr_pullbackFreeIso, lem:gr_pullbackObjUnitToUnit_comp, def:modules_pullbackComp}
  The composite analogue of \cref{lem:gr_pullbackFreeIso_id}: for composable \(a, b\) and an
  index set \(I\), the free-pullback isomorphism for \(b \circ a\) at \(\bigoplus_I \mathbf 1\)
  factors through \(a^{*}\) of the free-pullback isomorphism for \(b\), the one for \(a\), and
  the iterated-pullback comparison \cref{def:modules_pullbackComp}.
\end{lemma}

\begin{proof}
  \uses{lem:gr_pullbackObjUnitToUnit_comp}
  \leanok
  By coproduct extensionality it suffices to check the identity after precomposing with each
  coprojection \(\mathbf 1 \to \bigoplus_I \mathbf 1\), which reduces to the one-element
  composite coherence \cref{lem:gr_pullbackObjUnitToUnit_comp} (with the base-ring sheaf
  tracked through each pullback).
\end{proof}

Together \cref{lem:gr_pullbackObjUnitToUnit_id} and \cref{lem:gr_pullbackObjUnitToUnit_comp}
upgrade to the corresponding free-sheaf coherences -- the free-pullback isomorphism at the
identity equals the component of the pullback-identity isomorphism, and likewise for the
composite -- by coproduct extensionality over the decomposition
\(\mathcal{O}^{\,\oplus I} = \bigoplus_I \mathbf 1\). These coherences discharge both
functoriality laws (identity and composition) of the Grassmannian
functor \cref{def:grassmannian_functor}.

\begin{definition}

[The Grassmannian / Quot-of-trivial-sheaf functor \(\operatorname{Grass}(r,d)\)]
  \label{def:grassmannian_functor}
  \lean{AlgebraicGeometry.Grassmannian.functor}
  \uses{def:is_locally_free_of_rank, def:quot_functor, def:gr_rankQuotient,
    lem:gr_pullbackObjUnitToUnit_id, lem:gr_pullbackObjUnitToUnit_comp}
  % SOURCE: [Nitsure], "Construction of Hilbert and Quot Schemes", \S 1,
  % Exercise (2) "Grassmannian as a Quot scheme" (read from
  % references/nitsure-hilbert-quot-src/nitsure-hilbert-quot.tex, L557-L564).
  % SOURCE QUOTE:
  % "represents the
  %  contravariant functor
  %  $$\Grass(r,d) = \Quot_{\oplus^r\,{\mathcal O}_{\Z}/\Z/\Z}^{d, {\mathcal O}_{\Z}}$$
  %  from schemes to sets, which associates to any $T$ the
  %  set of all equivalence classes
  %  $\langle \F, q\rangle$ of quotients $q: \oplus^r\,{\mathcal O}_T \to \F$
  %  where $\F$ is a locally free sheaf on $T$ of rank $d$."
  \textit{Source: [Nitsure], \S 1, Exercise (2).}
  The \emph{Grassmannian functor}
  \(\operatorname{Grass}(r,d) = \Quot^{\,d,\,\mathcal{O}_{\mathbb{Z}}}_{\oplus^{r}
  \mathcal{O}_{\mathbb{Z}}/\mathbb{Z}/\mathbb{Z}}\) is the contravariant functor
  from schemes to sets defined as follows. For a scheme \(T\), a \emph{family of
  rank-\(d\) quotients of \(\mathcal{O}_T^{\,r}\)} is a pair
  \(\langle \mathcal{F}, q \rangle\) consisting of a locally free
  \(\mathcal{O}_T\)-module \(\mathcal{F}\) of rank \(d\) together with a surjective
  \(\mathcal{O}_T\)-linear homomorphism \(q : \mathcal{O}_T^{\,r}
  \twoheadrightarrow \mathcal{F}\). Two such families
  \(\langle \mathcal{F}, q \rangle\) and \(\langle \mathcal{F}', q' \rangle\) are
  \emph{equivalent} if there is an isomorphism \(f : \mathcal{F}
  \xrightarrow{\sim} \mathcal{F}'\) with \(f \circ q = q'\); equivalently
  \(\ker(q) = \ker(q')\) as subsheaves of \(\mathcal{O}_T^{\,r}\). Set
  \[
    \operatorname{Grass}(r,d)(T)
      \;=\; \bigl\{\, \text{equivalence classes } \langle \mathcal{F}, q \rangle
        \text{ of rank-}d\text{ quotients of } \mathcal{O}_T^{\,r} \,\bigr\}.
  \]
  For a morphism \(\varphi : T' \to T\), the map
  \(\operatorname{Grass}(r,d)(\varphi)\) sends \(\langle \mathcal{F}, q \rangle\)
  to the pulled-back quotient \(\langle \varphi^{*}\mathcal{F},\, \varphi^{*}q
  \rangle\); this is well-defined on equivalence classes because pullback is right
  exact and preserves local freeness of rank \(d\). (As \(\mathcal{F}\) is locally
  free of rank \(d\), surjectivity is preserved and the support is automatically
  proper -- all of \(T'\) -- so this is the special case
  \(E = \mathcal{O}^{r}\), \(X = S = \operatorname{base}\) of the general Quot
  functor of \cref{def:quot_functor}.)
\end{definition}

\section{The universal property}
\label{sec:grquot_universal}

The representability theorem \cref{thm:grassmannian_universal_property} has two halves:
the forward map \(\varphi \mapsto \langle \varphi^{*}\mathcal{U}, \varphi^{*}u\rangle\)
and its inverse, which manufactures a morphism \(T \to \mathrm{Gr}(r,d)\) out of a rank-\(d\)
quotient \(\langle \mathcal{F}, q\rangle\) on \(T\). The forward map needs the tautological
quotient to actually be a quotient -- that \(u\) is an epimorphism
(\cref{lem:tautologicalQuotient_epi}). The inverse is built Zariski-locally and glued: for
each size-\(d\) subset \(I\) one isolates the open locus \(T_I \subseteq T\) where the
\(I\)-indexed columns of \(q\) form a basis of \(\mathcal{F}\), classifies \(T_I \to U^I\)
there, and glues over the cover \(\{T_I\}\). The blocks of this section supply the four
ingredients of that inverse: the iso-locus of a sheaf morphism (\cref{def:isoLocus}) and
the fact that a morphism becomes invertible after restriction to its iso-locus
(\cref{lem:isIso_pullback_isoLocus_map}); the chart locus \(T_I\) (\cref{def:chartLocus})
and the proof that the chart loci cover \(T\) (\cref{lem:chartLocus_isOpenCover}); and the
glued classifying morphism itself (\cref{def:grPointOfRankQuotient}).

\begin{lemma}[The tautological quotient is an epimorphism]
  \label{lem:tautologicalQuotient_epi}
  \lean{AlgebraicGeometry.Grassmannian.tautologicalQuotient_epi}
  \uses{lem:gr_chartQuotientMap_epi, def:tautological_quotient}
  The tautological quotient \(u : \mathcal{O}_{\mathrm{Gr}(r,d)}^{\,r}
  \twoheadrightarrow \mathcal{U}\) (\cref{def:tautological_quotient}) is an epimorphism of
  sheaves of \(\mathcal{O}_{\mathrm{Gr}(r,d)}\)-modules.
\end{lemma}
\begin{proof}
  \uses{lem:gr_chartQuotientMap_epi, def:tautological_quotient}
  Being an epimorphism of sheaves of modules is local on the base: a morphism is epi
  precisely when its restriction to every member of an open cover is epi, since the
  cokernel is computed chartwise and a sheaf vanishing on a cover vanishes. The charts
  \(\{U^I\}\) cover \(\mathrm{Gr}(r,d)\), and over \(U^I\) the tautological quotient,
  transported through the trivialisation \(\mathcal{U}|_{U^I} \cong \mathcal{O}_{U^I}^{\,d}\)
  of \cref{def:gr_universal_quotient_sheaf}, is the chart quotient \(u^I :
  \mathcal{O}_{U^I}^{\,r} \to \mathcal{O}_{U^I}^{\,d}\) (\cref{def:tautological_quotient}).
  Each \(u^I\) is a split epimorphism: its \(I\)-minor is the identity \(1_{d\times d}\),
  so the coordinate inclusion \(\mathcal{O}_{U^I}^{\,d} \hookrightarrow
  \mathcal{O}_{U^I}^{\,r}\) on the \(I\)-columns is a section of \(u^I\), whence \(u^I\) is
  epi (\cref{lem:gr_chartQuotientMap_epi}). As \(u^I\) is epi on every chart, \(u\) is epi.
\end{proof}

\begin{definition}[The universal quotient sheaf restricts to the free sheaf on a chart]
  \label{def:gr_universalQuotient_restrictionIso}
  \lean{AlgebraicGeometry.Grassmannian.universalQuotient_restrictionIso}
  \uses{def:gr_universal_quotient_sheaf, def:gr_glued_scheme, def:glueRestrictionIso,
    def:gr_bundleTransitionData, lem:gr_bundleCocycle_id, lem:gr_bundleCocycle_mul}
  For each chart index \(I\), the restriction of the universal quotient sheaf
  \(\mathcal{U}\) (\cref{def:gr_universal_quotient_sheaf}) to the \(I\)-chart along the
  glue-data immersion \(\iota_I : U^I \to \mathrm{Gr}(r,d)\) is canonically isomorphic to
  the free rank-\(d\) sheaf,
  \[
    \phi_I \;:\; \iota_I^{*}\mathcal{U}
      \;\xrightarrow{\ \sim\ }\; \mathcal{O}_{U^I}^{\,d}.
  \]
  It is the instantiation, at the Grassmannian glue data (\cref{def:gr_glued_scheme}) and
  the bundle-transition cocycle \((g_{I,J})\) (\cref{def:gr_bundleTransitionData},
  \cref{lem:gr_bundleCocycle_id}, \cref{lem:gr_bundleCocycle_mul}), of the descent
  restriction isomorphism (\cref{def:glueRestrictionIso}); its underlying morphism is the
  adjoint transpose of the \(I\)-th descent-equaliser projection, an isomorphism because
  the descent comparison morphism is invertible.
\end{definition}

\begin{lemma}[The universal quotient sheaf is locally free of rank \(d\)]
  \label{lem:gr_universalQuotient_isLocallyFreeOfRank}
  \lean{AlgebraicGeometry.Grassmannian.universalQuotient_isLocallyFreeOfRank}
  \uses{def:gr_universal_quotient_sheaf, def:gr_universalQuotient_restrictionIso,
    def:gr_glued_scheme, def:is_locally_free_of_rank, def:modules_pullbackComp,
    lem:gr_pullbackFreeIso}
  The universal quotient sheaf \(\mathcal{U}\) (\cref{def:gr_universal_quotient_sheaf}) is
  locally free of rank \(d\) (\cref{def:is_locally_free_of_rank}).
\end{lemma}
\begin{proof}
  \uses{def:gr_universalQuotient_restrictionIso, def:modules_pullbackComp,
    lem:gr_pullbackFreeIso}
  The chart images \(\{\iota_I(U^I)\}\) form an open cover of \(\mathrm{Gr}(r,d)\), since
  the glue-data immersions are jointly surjective (\cref{def:gr_glued_scheme}). Over each
  member, factor the immersion as \(\iota_I = e_I \circ j_I\) through the isomorphism
  \(e_I : U^I \xrightarrow{\sim} \iota_I(U^I)\) onto its open image followed by the open
  inclusion \(j_I\); transporting the chart restriction isomorphism
  (\cref{def:gr_universalQuotient_restrictionIso}) across this factorisation -- inverting
  the parametrisation \(e_I\) through the pullback pseudofunctor comparison
  (\cref{def:modules_pullbackComp}) and the free-pullback comparison
  (\cref{lem:gr_pullbackFreeIso}) -- exhibits the restriction
  \(j_I^{*}\mathcal{U} \cong \mathcal{O}_{\iota_I(U^I)}^{\,d}\) as free of rank \(d\). As
  the \(\iota_I(U^I)\) cover \(\mathrm{Gr}(r,d)\), \(\mathcal{U}\) is locally free of
  rank \(d\).
\end{proof}

\begin{definition}[The tautological point as a rank-\(d\) quotient]
  \label{def:gr_tautologicalRankQuotient}
  \lean{AlgebraicGeometry.Grassmannian.tautologicalRankQuotient}
  \uses{def:gr_rankQuotient, def:gr_universal_quotient_sheaf, def:tautological_quotient,
    lem:tautologicalQuotient_epi, lem:gr_universalQuotient_isLocallyFreeOfRank}
  The tautological pair \(\langle \mathcal{U}, u\rangle\) on \(\mathrm{Gr}(r,d)\) itself,
  packaged as a rank-\(d\) quotient of \(\mathcal{O}_{\mathrm{Gr}(r,d)}^{\,r}\)
  (\cref{def:gr_rankQuotient}): its sheaf is the universal quotient sheaf \(\mathcal{U}\)
  (\cref{def:gr_universal_quotient_sheaf}), its quotient map the tautological quotient
  \(u\) (\cref{def:tautological_quotient}), which is an epimorphism
  (\cref{lem:tautologicalQuotient_epi}) onto a sheaf locally free of rank \(d\)
  (\cref{lem:gr_universalQuotient_isLocallyFreeOfRank}). Pulling it back along a morphism
  \(\psi : T \to \mathrm{Gr}(r,d)\) realises the forward direction of the universal
  property.
\end{definition}

\begin{definition}[Iso-locus of a morphism of sheaves of modules]
  \label{def:isoLocus}
  \lean{AlgebraicGeometry.Grassmannian.isoLocus}
  Let \(X\) be a scheme and \(\varphi : \mathcal{M} \to \mathcal{N}\) a morphism of sheaves
  of \(\mathcal{O}_X\)-modules. The \emph{iso-locus} \(\mathrm{isoLocus}(\varphi) \subseteq X\)
  is the open subscheme
  \[
    \mathrm{isoLocus}(\varphi)
      \;=\; \bigsqcup\,\bigl\{\, U \subseteq X \text{ open} \;:\;
        \varphi|_U := (\iota_U^{*}\varphi) \text{ is an isomorphism} \,\bigr\},
  \]
  the supremum (union) of all opens \(U\) over which the pullback \(\iota_U^{*}\varphi\) of
  \(\varphi\) to the open subscheme \(U\) is an isomorphism of sheaves of
  \(\mathcal{O}_U\)-modules. Equivalently, a point \(t \in X\) lies in
  \(\mathrm{isoLocus}(\varphi)\) iff \(\varphi\) restricts to an isomorphism over some open
  neighbourhood of \(t\). It is the largest open on which \(\varphi\) is invertible; Mathlib
  has no iso-locus primitive for morphisms of sheaves of modules, so it is set up directly
  as this supremum.
\end{definition}

\begin{lemma}[Stalk-wise iso criterion]
  \label{lem:isIso_of_stalkFunctor_map_iso_mathlib}
  \lean{TopCat.Presheaf.isIso_of_stalkFunctor_map_iso}
  \mathlibok
  \textit{Provided by Mathlib.}
  Let \(\mathcal{F}, \mathcal{G}\) be sheaves on a topological space \(X\) valued in a
  concrete category whose forgetful functor reflects isomorphisms and preserves limits and
  filtered colimits, and let \(f : \mathcal{F} \to \mathcal{G}\) be a morphism. If the
  induced map on stalks \(f_t : \mathcal{F}_t \to \mathcal{G}_t\) is an isomorphism for every
  point \(t \in X\), then \(f\) is an isomorphism.
\end{lemma}

\begin{lemma}[A morphism is invertible over its iso-locus]
  \label{lem:isIso_pullback_isoLocus_map}
  \lean{AlgebraicGeometry.Grassmannian.isIso_pullback_isoLocus_map}
  \uses{def:isoLocus, lem:isIso_of_stalkFunctor_map_iso_mathlib}
  For a morphism \(\varphi : \mathcal{M} \to \mathcal{N}\) of sheaves of
  \(\mathcal{O}_X\)-modules, the pullback of \(\varphi\) to the iso-locus
  (\cref{def:isoLocus}),
  \[
    (\iota_{\mathrm{isoLocus}(\varphi)})^{*}\varphi
      \;:\;
    (\iota_{\mathrm{isoLocus}(\varphi)})^{*}\mathcal{M}
      \;\longrightarrow\;
    (\iota_{\mathrm{isoLocus}(\varphi)})^{*}\mathcal{N},
  \]
  is an isomorphism.
\end{lemma}
\begin{proof}
  \uses{def:isoLocus, lem:isIso_of_stalkFunctor_map_iso_mathlib}
  Being an isomorphism is a stalk-local property: by the stalk-wise criterion
  (\cref{lem:isIso_of_stalkFunctor_map_iso_mathlib}) it suffices to show that the pulled-back
  morphism induces an isomorphism on the stalk at every point \(t \in
  \mathrm{isoLocus}(\varphi)\). By the definition of the iso-locus as a union of opens on
  which \(\varphi\) restricts to an isomorphism (\cref{def:isoLocus}), each such \(t\) has an
  open neighbourhood \(U \subseteq \mathrm{isoLocus}(\varphi)\) with \(\iota_U^{*}\varphi\) an
  isomorphism; restriction of an isomorphism is an isomorphism, and pullback commutes with
  taking stalks, so the stalk of the morphism at \(t\) coincides with the stalk of
  \(\iota_U^{*}\varphi\) at \(t\), which is invertible. Hence every stalk map is an
  isomorphism, and the criterion concludes that the global pullback over
  \(\mathrm{isoLocus}(\varphi)\) is an isomorphism.
\end{proof}

\begin{definition}[Chart composite of a rank-\(d\) quotient]
  \label{def:gr_chartComposite}
  \lean{AlgebraicGeometry.Grassmannian.chartComposite}
  \uses{def:gr_rankQuotient}
  Let \(x = \langle \mathcal{F}, q\rangle\) be a rank-\(d\) quotient of
  \(\mathcal{O}_T^{\,r}\) (\cref{def:gr_rankQuotient}) and \(I \subseteq \{1,\dots,r\}\) a
  size-\(d\) subset, enumerated in increasing order by the order isomorphism
  \(\{1,\dots,d\} \xrightarrow{\sim} I\). The \emph{chart composite} is the morphism of
  sheaves of modules
  \[
    s_I \circ q \;:\; \mathcal{O}_T^{\,d}
      \overset{\ s_I\ }{\hookrightarrow} \mathcal{O}_T^{\,r}
      \xrightarrow{\ q\ } \mathcal{F},
  \]
  the quotient \(q\) precomposed with the coordinate inclusion \(s_I\) of the
  \(I\)-indexed factors of the trivial bundle. Over the locus where it is invertible the
  \(I\)-columns of \(q\) form a basis of \(\mathcal{F}\).
\end{definition}

\begin{definition}[Chart locus of a rank-\(d\) quotient]
  \label{def:chartLocus}
  \lean{AlgebraicGeometry.Grassmannian.chartLocus}
  \uses{def:isoLocus, def:gr_chartComposite, def:gr_rankQuotient}
  Let \(x = \langle \mathcal{F}, q\rangle\) be a rank-\(d\) quotient of
  \(\mathcal{O}_T^{\,r}\) on a scheme \(T\) (\cref{def:gr_rankQuotient}) and \(I \subseteq
  \{1,\dots,r\}\) a size-\(d\) subset. The \emph{chart composite} is the morphism of sheaves
  of modules
  \[
    s_I \circ q
      \;:\;
    \mathcal{O}_T^{\,d}
      \overset{\ s_I\ }{\hookrightarrow} \mathcal{O}_T^{\,r}
      \xrightarrow{\ q\ } \mathcal{F},
  \]
  obtained by precomposing the quotient \(q\) with the inclusion \(s_I\) of the
  \(I\)-indexed coordinates of the trivial bundle. The \emph{chart locus}
  \(T_I = \mathrm{chartLocus}(x, I) \subseteq T\) is the iso-locus
  (\cref{def:isoLocus}) of this chart composite:
  \[
    T_I \;=\; \mathrm{isoLocus}(s_I \circ q),
  \]
  the largest open of \(T\) over which the \(I\)-columns of \(q\) form a basis of
  \(\mathcal{F}\) -- i.e.\ over which \(s_I \circ q\) is an isomorphism. Over \(T_I\) the
  quotient \(q\) is therefore presented, after trivialising \(\mathcal{F}\) by
  \((s_I \circ q)^{-1}\), by a \(d \times r\) matrix whose \(I\)-minor is the identity,
  exactly the shape of the universal matrix \(X^I\) realising the chart quotient \(u^I\)
  (\cref{def:gr_chart_quotient}).
\end{definition}

\begin{lemma}[The chart loci cover \(T\)]
  \label{lem:chartLocus_isOpenCover}
  \lean{AlgebraicGeometry.Grassmannian.chartLocus_isOpenCover}
  \uses{def:chartLocus, def:gr_rankQuotient, lem:gr_exists_isUnit_submatrix}
  % SOURCE: [Nitsure], "Construction of Hilbert and Quot Schemes", \S 1,
  % sub-section "Construction by gluing together affine patches" (read from
  % references/nitsure-hilbert-quot-src/nitsure-hilbert-quot.tex, L823-L828).
  % SOURCE QUOTE:
  % "For any $J\subset \{1,\ldots, r\}$ with $\#(J) =d$, let
  %  $P^I_J = \det(X^I_J) \in \Z[X^I]$
  %  where $X^I_J$ is the $J$ th minor of $X^I$.
  %  Let $U^I_J = \spec \Z[X^I, 1/P^I_J]$
  %  the open subscheme of $U^I$ where $P^I_J$ is invertible. This means
  %  the $d\times d$-matrix $X^I_J$ admits an inverse $(X^I_J)^{-1}$
  %  on $U^I_J$."
  \textit{Source: [Nitsure], \S 1.}
  For a rank-\(d\) quotient \(x = \langle \mathcal{F}, q\rangle\) of \(\mathcal{O}_T^{\,r}\)
  on \(T\) (\cref{def:gr_rankQuotient}), the family of chart loci
  \(\{\,T_I = \mathrm{chartLocus}(x, I)\,\}_{\#(I) = d}\) (\cref{def:chartLocus}) is an open
  cover of \(T\): \(\bigcup_{\#(I)=d} T_I = T\).
\end{lemma}
\begin{proof}
  \uses{def:chartLocus, def:gr_rankQuotient, def:isoLocus, def:gr_matrixEndRect,
    lem:gr_exists_rightInverse_of_epi_matrixEndRect, lem:gr_exists_isUnit_submatrix}
  Each \(T_I\) is open, being an iso-locus (\cref{def:isoLocus}). It remains to show the
  \(T_I\) cover \(T\), i.e.\ that every point \(t \in T\) lies in some \(T_I\). Fix \(t\) and
  choose an \emph{affine} open \(W \ni t\) over which \(\mathcal{F}\) is trivialised,
  \(\mathcal{F}|_W \cong \mathcal{O}_W^{\,d}\) (such \(W\) exists because \(\mathcal{F}\) is
  locally free of rank \(d\)); over \(W\) the quotient \(q\) is presented by a \(d \times r\)
  matrix \(M\) of sections of \(\mathcal{O}_W\), i.e.\ \(q|_W\) is conjugate to
  \(\mathrm{matrixEndRect}(M)\) (\cref{def:gr_matrixEndRect}).

  \emph{Global splitting over \(W\).} Because \(q\) is an epimorphism
  (\cref{def:gr_rankQuotient}), so is the presenting homomorphism
  \(\mathrm{matrixEndRect}(M)\) of free sheaves on the affine \(W\). Over an affine base such
  an epimorphism of free sheaves splits -- free modules are projective, so the epimorphism
  admits a section -- and reading the section off as a matrix yields a right inverse
  \(G \in \operatorname{Mat}_{r\times d}(\Gamma(W,\mathcal{O}_W))\) with \(M \cdot G = 1\)
  (\cref{lem:gr_exists_rightInverse_of_epi_matrixEndRect}). Thus \(q\) is globally split over
  \(W\), with no recourse to stalks or Nakayama propagation.

  \emph{Choosing the index set at \(t\).} Reduce the identity \(M \cdot G = 1\) modulo the
  maximal ideal of the local ring at \(t\): writing \(\kappa(t)\) for the residue field, the
  images \(\bar M, \bar G\) over \(\kappa(t)\) still satisfy \(\bar M \cdot \bar G = 1\), so
  \(\bar M\) is a right-invertible \(d \times r\) matrix over the field \(\kappa(t)\). A
  right-invertible matrix over a field has an invertible \(d\)-column minor: its columns span
  \(\kappa(t)^{d}\), a spanning family contains a basis, and the corresponding \(d\) columns
  form an invertible \(d \times d\) submatrix (\cref{lem:gr_exists_isUnit_submatrix}). Pick
  such a size-\(d\) index set \(I \subseteq \{1,\dots,r\}\).

  \emph{Cutting out the basic open.} Let \(P_I = \det(M_I)\) be the determinant of the
  \(I\)-indexed \(d \times d\) minor of \(M\). By the previous step its residue at \(t\) is
  nonzero, so \(P_I\) is a unit at \(t\); hence \(t\) lies in the basic open
  \(W_{P_I} = \{P_I \ne 0\} \subseteq W\) where \(P_I\) is invertible. On \(W_{P_I}\) the
  minor matrix \(M_I\) is invertible, so the chart composite \(s_I \circ q\) -- conjugate
  there to \(\mathrm{matrixEndRect}(M_I)\) -- restricts to an isomorphism, i.e.\
  \(W_{P_I} \subseteq \mathrm{isoLocus}(s_I \circ q) = T_I\) (\cref{def:chartLocus}).
  Therefore \(t \in T_I\), and the chart loci cover \(T\).
\end{proof}

\begin{definition}[Chart matrix of a rank-\(d\) quotient]
  \label{def:gr_chartMatrix}
  \lean{AlgebraicGeometry.Grassmannian.chartMatrix}
  \uses{def:gr_chartComposite, def:chartLocus, lem:isIso_pullback_isoLocus_map,
    lem:gr_pullbackFreeIso, def:gr_presentedMatrix}
  Over the chart locus \(T_I = \mathrm{chartLocus}(x, I)\) (\cref{def:chartLocus}) the chart
  composite \(s_I \circ q\) pulls back to an isomorphism
  (\cref{lem:isIso_pullback_isoLocus_map}). Trivialising \(\mathcal{F}|_{T_I}\) by its
  inverse presents the restricted quotient as the conjugated free-sheaf homomorphism
  \[
    Q_r^{-1} \circ \iota_{T_I}^{*}q \circ (\iota_{T_I}^{*}(s_I \circ q))^{-1} \circ Q_d
      \;:\; \mathcal{O}_{T_I}^{\,r} \longrightarrow \mathcal{O}_{T_I}^{\,d},
  \]
  with \(Q_r, Q_d\) the free-pullback comparisons (\cref{lem:gr_pullbackFreeIso}). The
  \emph{chart matrix} \(M^I = \mathrm{chartMatrix}(x, I) \in
  \operatorname{Mat}_{d\times r}(\Gamma(T_I, \mathcal{O}_{T_I}))\) reads off this
  homomorphism entrywise through the section-extraction \(\mathrm{unitEndSection}\). It is
  the special case of the presented matrix (\cref{def:gr_presentedMatrix}) along the
  chart-locus inclusion; its \(I\)-minor is the identity \(1_{d\times d}\).
\end{definition}

\begin{definition}[Chart morphism of a rank-\(d\) quotient]
  \label{def:gr_chartMorphism}
  \lean{AlgebraicGeometry.Grassmannian.chartMorphism}
  \uses{def:gr_chartMatrix, def:gr_affine_chart, def:gr_universal_matrix}
  The \emph{chart morphism} \(\varphi_I : T_I \to U^I\) into the affine chart
  \(U^I = \Spec \mathbb{Z}[X^I]\) (\cref{def:gr_affine_chart}) is the morphism classified,
  through the \(\Gamma\)--\(\Spec\) adjunction, by the \(\mathbb{Z}\)-algebra homomorphism
  \(R^I = \mathbb{Z}[X^I] \to \Gamma(T_I, \mathcal{O}_{T_I})\) sending each variable
  \(x^I_{p,q}\) of the universal matrix \(X^I\) (\cref{def:gr_universal_matrix}) to the
  corresponding entry of the chart matrix \(M^I\) (\cref{def:gr_chartMatrix}). By
  construction \(\varphi_I^{*}X^I = M^I\), and hence \(\varphi_I^{*}u^I = q|_{T_I}\) in the
  \(I\)-trivialisation.
\end{definition}

\subsection{Overlap compatibility of the chart morphisms}

On the intersection \(P = T_I \times_T T_J\) of two chart loci both chart composites pull
back to isomorphisms, so the pulled-back quotient is presented by both chart matrices. The
blocks below assemble the transports comparing the two presentations and proving the two
chart morphisms agree into the glued scheme: invertibility along a composite
(\cref{lem:gr_isIso_pullback_map_comp}), the presented matrix along a composite
(\cref{lem:gr_presentedMatrix_comp}) and under an equality of base morphisms
(\cref{lem:gr_presentedMatrix_congr}), the chart matrix as a presented matrix
(\cref{lem:gr_chartMatrix_eq_presentedMatrix}), its own minor
(\cref{lem:gr_presentedMatrix_submatrix_self}), the realisation
\(\mathrm{aeval}(M^I)(X^I) = M^I\) (\cref{lem:gr_universalMatrix_map_presentedMatrix}), the
image-matrix transport (\cref{lem:gr_imageMatrix_map_ringHom}), chart-morphism naturality
(\cref{lem:gr_comp_chartMorphism}), the glue-condition identity
(\cref{lem:gr_chart_point_eq}), and finally the overlap compatibility itself
(\cref{lem:gr_chartMorphism_glue_compat}).

\begin{lemma}[Pullback-invertibility persists along precomposition]
  \label{lem:gr_isIso_pullback_map_comp}
  \lean{AlgebraicGeometry.Grassmannian.isIso_pullback_map_comp}
  \uses{def:modules_pullbackComp}
  For composable scheme morphisms \(p : W \to V\), \(a : V \to X\) and a morphism
  \(c : \mathcal{M} \to \mathcal{N}\) of sheaves of \(\mathcal{O}_X\)-modules: if the
  pullback \(a^{*}c\) is an isomorphism, so is \((p \circ a)^{*}c\).
\end{lemma}
\begin{proof}
  \uses{def:modules_pullbackComp}
  Functoriality of pullback makes \(p^{*}(a^{*}c)\) an isomorphism; the pseudofunctor
  comparison \((p \circ a)^{*} \cong p^{*}a^{*}\) (\cref{def:modules_pullbackComp})
  identifies \((p \circ a)^{*}c\) with \(p^{*}(a^{*}c)\), which is invertible.
\end{proof}

\begin{lemma}[The presented matrix depends only on the base morphism]
  \label{lem:gr_presentedMatrix_congr}
  \lean{AlgebraicGeometry.Grassmannian.presentedMatrix_congr}
  \uses{def:gr_presentedMatrix}
  If \(j = j'\) as morphisms \(P \to T\), then the presented matrices along \(j\) and \(j'\)
  agree, \(\mathrm{presentedMatrix}(x, j, I) = \mathrm{presentedMatrix}(x, j', I)\); the
  choice of invertibility witness for the pulled-back chart composite is irrelevant.
\end{lemma}
\begin{proof}
  \uses{def:gr_presentedMatrix}
  Substitute the equality \(j = j'\); the invertibility hypotheses are propositions, hence
  proof-irrelevant, and the two sides become definitionally equal.
\end{proof}

\begin{lemma}[The chart matrix is the presented matrix along the locus inclusion]
  \label{lem:gr_chartMatrix_eq_presentedMatrix}
  \lean{AlgebraicGeometry.Grassmannian.chartMatrix_eq_presentedMatrix}
  \uses{def:gr_chartMatrix, def:gr_presentedMatrix}
  The chart matrix \(M^I = \mathrm{chartMatrix}(x, I)\) (\cref{def:gr_chartMatrix}) equals
  the presented matrix (\cref{def:gr_presentedMatrix}) of \(x\) along the chart-locus
  inclusion \(\iota_{T_I} : T_I \hookrightarrow T\):
  \(\mathrm{chartMatrix}(x, I) = \mathrm{presentedMatrix}(x, \iota_{T_I}, I)\).
\end{lemma}
\begin{proof}
  \uses{def:gr_chartMatrix, def:gr_presentedMatrix}
  Both sides are the same \(\mathrm{unitEndSection}\)-matrix of the same conjugated
  free-sheaf homomorphism; the equality is definitional.
\end{proof}

\begin{lemma}[Transport of the presented matrix along a composite]
  \label{lem:gr_presentedMatrix_comp}
  \lean{AlgebraicGeometry.Grassmannian.presentedMatrix_comp}
  \uses{def:gr_presentedMatrix, lem:gr_matrixEndRect_presentedMatrix,
    lem:gr_matrixEndRect_injective}
  Along composable \(p : W \to V\), \(a : V \to T\) on which the \(I\)-chart composite
  pulls back to isomorphisms, the presented matrix after pulling back along \(p \circ a\)
  is the entrywise restriction along \(p\) of the presented matrix after pulling back
  along \(a\):
  \[
    \mathrm{presentedMatrix}(x, p \circ a, I)
      \;=\; \bigl(\mathrm{presentedMatrix}(x, a, I)\bigr)\big|_{p},
  \]
  the matrix obtained by applying the restriction ring map \(\Gamma(V, \mathcal{O}_V) \to
  \Gamma(W, \mathcal{O}_W)\) of \(p\) to every entry.
\end{lemma}
\begin{proof}
  \uses{lem:gr_matrixEndRect_presentedMatrix, lem:gr_matrixEndRect_injective}
  By injectivity of the rectangular matrix presentation
  (\cref{lem:gr_matrixEndRect_injective}) it suffices to compare after applying
  \(\mathrm{matrixEndRect}\). The presentation identity
  (\cref{lem:gr_matrixEndRect_presentedMatrix}) expresses both sides as conjugated
  pulled-back homomorphisms; the conjugation distributes over the composite \(p \circ a\)
  through the pullback pseudofunctor comparisons, and the entrywise restriction along
  \(p\) is exactly the pullback along \(p\) of the \(a\)-level homomorphism. Comparing the
  two yields the claimed equality.
\end{proof}

\begin{lemma}[The own-minor of the presented matrix is the identity]
  \label{lem:gr_presentedMatrix_submatrix_self}
  \lean{AlgebraicGeometry.Grassmannian.presentedMatrix_submatrix_self}
  \uses{def:gr_presentedMatrix, lem:gr_matrixEndRect_presentedMatrix_minor,
    lem:gr_matrixEndRect_one, lem:gr_matrixEndRect_injective}
  Along any \(j\) on which the \(I\)-chart composite pulls back to an isomorphism, the
  \(I\)-minor of the presented matrix is the identity:
  \(\mathrm{presentedMatrix}(x, j, I)_{\bullet,\, I} = 1_{d\times d}\).
\end{lemma}
\begin{proof}
  \uses{lem:gr_matrixEndRect_presentedMatrix_minor, lem:gr_matrixEndRect_one,
    lem:gr_matrixEndRect_injective}
  By injectivity (\cref{lem:gr_matrixEndRect_injective}) compare after
  \(\mathrm{matrixEndRect}\). The minor presentation identity
  (\cref{lem:gr_matrixEndRect_presentedMatrix_minor}) at \(J = I\) presents the
  \(I\)-chart composite against its own inverse, which is the identity morphism, i.e.\
  \(\mathrm{matrixEndRect}(1_{d\times d})\) (\cref{lem:gr_matrixEndRect_one}).
\end{proof}

\begin{lemma}[The classifying map realises the universal matrix as the presented matrix]
  \label{lem:gr_universalMatrix_map_presentedMatrix}
  \lean{AlgebraicGeometry.Grassmannian.universalMatrix_map_presentedMatrix}
  \uses{def:gr_universal_matrix, def:gr_presentedMatrix,
    lem:gr_presentedMatrix_submatrix_self}
  Applying entrywise the \(\mathbb{Z}\)-algebra map \(\mathrm{aeval}(M)\) -- which sends
  the free variables of \(X^I\) to the complementary entries of
  \(M = \mathrm{presentedMatrix}(x, j, I)\) -- to the universal matrix \(X^I\)
  (\cref{def:gr_universal_matrix}) returns \(M\):
  \(\mathrm{aeval}(M)(X^I) = M\).
\end{lemma}
\begin{proof}
  \uses{lem:gr_presentedMatrix_submatrix_self}
  Entrywise. On a complementary (free) entry \(x^I_{p,q}\) with \(q \notin I\),
  \(\mathrm{aeval}(M)\) returns the corresponding entry of \(M\) by the variable-evaluation
  rule. On an \(I\)-column entry the universal matrix has the identity minor, so the value
  there is the corresponding entry of \(1_{d\times d}\), which matches the \(I\)-minor of
  \(M\) -- the identity by \cref{lem:gr_presentedMatrix_submatrix_self}.
\end{proof}

\begin{lemma}[Transport of the image matrix along a ring homomorphism]
  \label{lem:gr_imageMatrix_map_ringHom}
  \lean{AlgebraicGeometry.Grassmannian.imageMatrix_map_ringHom}
  \uses{def:gr_universalMinorInv, def:gr_universal_matrix, lem:gr_minorDet_unit,
    lem:gr_map_nonsing_inv}
  For size-\(d\) subsets \(I, J\) and any ring homomorphism
  \(\mathrm{incl} : R^I_J \to D\) out of the localised overlap ring whose composite with
  the structure map \(R^I \to R^I_J\) equals \(g\), pushing the image matrix
  \((X^I_J)^{-1} X^I\) forward along \(\mathrm{incl}\) yields \(\bigl((g_{*}X^I)_J\bigr)^{-1}
  (g_{*}X^I)\), where \(g_{*}X^I\) is \(X^I\) base-changed along \(g\):
  \[
    \mathrm{incl}_{*}\bigl((X^I_J)^{-1} X^I\bigr)
      \;=\; \bigl((g_{*}X^I)_J\bigr)^{-1}\, (g_{*}X^I).
  \]
\end{lemma}
\begin{proof}
  \uses{lem:gr_minorDet_unit, lem:gr_map_nonsing_inv}
  The image matrix is the product of the Cramer inverse \(g_{I,J} = (X^I_J)^{-1}\)
  (\cref{def:gr_universalMinorInv}) with \(X^I\); ring homomorphisms commute with matrix
  multiplication, with the submatrix (minor) operation, and -- because the minor
  determinant \(P^I_J\) is a unit (\cref{lem:gr_minorDet_unit}) -- with the nonsingular
  inverse (\cref{lem:gr_map_nonsing_inv}). Pushing each factor forward along
  \(\mathrm{incl}\) and rewriting the structure-map composite as \(g\) gives the stated
  base-changed expression.
\end{proof}

\begin{lemma}[Naturality of the chart morphism]
  \label{lem:gr_comp_chartMorphism}
  \lean{AlgebraicGeometry.Grassmannian.comp_chartMorphism}
  \uses{def:gr_chartMorphism, lem:gr_chartMatrix_eq_presentedMatrix,
    lem:gr_presentedMatrix_comp}
  For \(p : W \to T_I\) on which the pulled-back \(I\)-chart composite is invertible, the
  composite \(p \circ \varphi_I : W \to U^I\) is the morphism classified by the matrix
  presenting the quotient along \(p \circ \iota_{T_I}\) -- i.e.\ classified by
  \(\mathrm{presentedMatrix}(x, p \circ \iota_{T_I}, I)\).
\end{lemma}
\begin{proof}
  \uses{lem:gr_chartMatrix_eq_presentedMatrix, lem:gr_presentedMatrix_comp}
  Naturality of the \(\Gamma\)--\(\Spec\) classifying map turns \(p \circ \varphi_I\) into
  the morphism classified by the \(p\)-restriction of the chart matrix. The chart matrix is
  the presented matrix along \(\iota_{T_I}\) (\cref{lem:gr_chartMatrix_eq_presentedMatrix}),
  and its \(p\)-restriction is the presented matrix along \(p \circ \iota_{T_I}\)
  (\cref{lem:gr_presentedMatrix_comp}); the two classifying maps therefore coincide.
\end{proof}

\begin{lemma}[Chart-point identity over a commutative ring]
  \label{lem:gr_chart_point_eq}
  \lean{AlgebraicGeometry.Grassmannian.chart_point_eq}
  \uses{def:gr_transition, def:gr_glued_scheme, lem:gr_chartTransition_comp_chartIncl}
  Let \(K\) be a commutative ring and \(f : R^I \to K\) a ring map under which the minor
  \(P^I_J\) is a unit. Then the transported classifying map
  \(\mathrm{lift}(f) \circ \tilde\theta_{I,J}\) presents the same \(K\)-point through chart
  \(J\) as \(f\) does through chart \(I\):
  \[
    \Spec(\mathrm{lift}(f) \circ \tilde\theta_{I,J}) \circ \iota_J
      \;=\; \Spec(f) \circ \iota_I.
  \]
\end{lemma}
\begin{proof}
  \uses{def:gr_transition, def:gr_glued_scheme, lem:gr_chartTransition_comp_chartIncl}
  Factor \(\tilde\theta_{I,J}\) as the chart transition followed by the chart inclusion
  (\cref{lem:gr_chartTransition_comp_chartIncl}); the glue condition of the glue data
  (\cref{def:gr_glued_scheme}) makes the transition square commute, and the localisation
  lift recovers \(f\) after the chart inclusion. Assembling the two reductions identifies
  the two composites into the glued scheme.
\end{proof}

\begin{lemma}[Overlap compatibility of the chart morphisms]
  \label{lem:gr_chartMorphism_glue_compat}
  \lean{AlgebraicGeometry.Grassmannian.chartMorphism_glue_compat}
  \uses{def:gr_chartMorphism, def:gr_glued_scheme, lem:gr_isIso_pullback_map_comp,
    lem:gr_comp_chartMorphism, lem:gr_presentedMatrix_changeOfBasis,
    lem:gr_universalMatrix_map_presentedMatrix, lem:gr_imageMatrix_map_ringHom,
    lem:gr_chart_point_eq, lem:gr_isUnit_of_isIso_matrixEndRect,
    lem:gr_matrixEndRect_presentedMatrix_minor}
  On the intersection \(P = T_I \times_T T_J\) the two composites of the chart morphisms
  with the glue-data immersions agree:
  \[
    \mathrm{pr}_1 \circ (\varphi_I \circ \iota_I)
      \;=\; \mathrm{pr}_2 \circ (\varphi_J \circ \iota_J).
  \]
  This is the gluing hypothesis from which \(\mathrm{grPointOfRankQuotient}(x)\) is
  assembled.
\end{lemma}
\begin{proof}
  \uses{lem:gr_isIso_pullback_map_comp, lem:gr_comp_chartMorphism,
    lem:gr_presentedMatrix_changeOfBasis, lem:gr_universalMatrix_map_presentedMatrix,
    lem:gr_imageMatrix_map_ringHom, lem:gr_chart_point_eq,
    lem:gr_isUnit_of_isIso_matrixEndRect, lem:gr_matrixEndRect_presentedMatrix_minor}
  On \(P\) both chart composites pull back to isomorphisms
  (\cref{lem:gr_isIso_pullback_map_comp}), so the pulled-back quotient is presented by both
  matrices \(M^I_P\) and \(M^J_P\); through \cref{lem:gr_comp_chartMorphism} the two
  composites are classified by these matrices. The \(J\)-minor of \(M^I_P\) is invertible
  (\cref{lem:gr_isUnit_of_isIso_matrixEndRect},
  \cref{lem:gr_matrixEndRect_presentedMatrix_minor}), and the change-of-basis identity
  \(M^I_P = (M^I_P)_J \cdot M^J_P\) (\cref{lem:gr_presentedMatrix_changeOfBasis}) --
  realised on the universal matrices via
  \cref{lem:gr_universalMatrix_map_presentedMatrix} and the image-matrix transport
  \cref{lem:gr_imageMatrix_map_ringHom} -- identifies the \(J\)-classifying ring map with
  the localisation-transport \(\mathrm{lift}(\mathrm{aeval}\,M^I_P) \circ
  \tilde\theta_{I,J}\) of the \(I\)-classifying map. The glue-condition identity
  (\cref{lem:gr_chart_point_eq}) then closes the square into the glued scheme.
\end{proof}

\begin{definition}[The morphism induced by a rank-\(d\) quotient]
  \label{def:grPointOfRankQuotient}
  \lean{AlgebraicGeometry.Grassmannian.grPointOfRankQuotient}
  \uses{lem:chartLocus_isOpenCover, lem:isIso_pullback_isoLocus_map,
    def:gr_affine_chart, def:gr_chart_quotient, def:gr_transition, lem:gr_cocycle,
    def:gr_glued_scheme, def:gr_chartComposite, def:gr_chartMatrix,
    def:gr_chartMorphism, lem:gr_chartMorphism_glue_compat}
  Let \(x = \langle \mathcal{F}, q\rangle\) be a rank-\(d\) quotient of
  \(\mathcal{O}_T^{\,r}\) on a scheme \(T\) (\cref{def:gr_rankQuotient}). The morphism
  \[
    \mathrm{grPointOfRankQuotient}(x) \;:\; T \;\longrightarrow\; \mathrm{Gr}(r,d)
  \]
  is constructed as follows. Over each chart locus \(T_I = \mathrm{chartLocus}(x, I)\)
  (\cref{def:chartLocus}) the chart composite \(s_I \circ q\) is an isomorphism
  (\cref{lem:isIso_pullback_isoLocus_map}), so trivialising \(\mathcal{F}|_{T_I}\) by its
  inverse presents \(q|_{T_I}\) as a \(d \times r\) matrix \(M^I\) of global functions on
  \(T_I\) whose \(I\)-minor is the identity. Its \(d(r-d)\) complementary entries determine a
  unique \(\mathbb{Z}\)-algebra homomorphism \(R^I = \mathbb{Z}[X^I] \to \Gamma(T_I,
  \mathcal{O}_{T_I})\), \(x^I_{p,q} \mapsto M^I_{p,q}\), hence by the \(\Gamma\)--\(\Spec\)
  adjunction a unique \emph{chart morphism} \(\varphi_I : T_I \to U^I\)
  (\cref{def:gr_affine_chart}) with \(\varphi_I^{*}X^I = M^I\), and thus
  \(\varphi_I^{*}u^I = q|_{T_I}\) (\cref{def:gr_chart_quotient}). The chart loci cover \(T\)
  (\cref{lem:chartLocus_isOpenCover}), and on each overlap \(T_I \cap T_J\) the two
  presenting matrices are change-of-basis related by the invertible \(I\)-columns minor,
  \(M^J = (M^I_J)^{-1} M^I = \varphi_I^{*}\theta_{I,J}(X^J)\)
  (\cref{def:gr_transition}); hence the composites \(T_I \xrightarrow{\varphi_I} U^I
  \hookrightarrow \mathrm{Gr}(r,d)\) with the chart immersions agree on overlaps after the
  transition \(\theta_{I,J}\). By the cocycle condition (\cref{lem:gr_cocycle}) these
  chart-local morphisms glue to a single morphism \(T \to \mathrm{Gr}(r,d)\)
  (\cref{def:gr_glued_scheme}); this glued morphism is
  \(\mathrm{grPointOfRankQuotient}(x)\), the local-to-global inverse of the universal
  property: pulling the tautological quotient back along it recovers \(\langle \mathcal{F},
  q\rangle\) up to equivalence.
\end{definition}

The two blocks below carry the chart data of a quotient through a base change
\(\psi : T' \to T\); they are the bridge underlying the two inverse laws of the
representability theorem.

\begin{lemma}[Chart composite of a pulled-back rank-\(d\) quotient]
  \label{lem:gr_chartComposite_rqPullback}
  \lean{AlgebraicGeometry.Grassmannian.chartComposite_rqPullback}
  \uses{def:gr_chartComposite, def:gr_rankQuotient, lem:gr_pullbackFreeIso_inv_freeMap}
  For a morphism \(\psi : T' \to T\) and a rank-\(d\) quotient \(y\) on \(T\), the chart
  composite of the pulled-back quotient \(\psi^{*}y\) (\cref{def:gr_rankQuotient}) is the
  \(\psi\)-pullback of the chart composite of \(y\), conjugated through the free-pullback
  comparison:
  \[
    s_I \circ (\psi^{*}y)
      \;=\; Q_d^{-1} \circ \psi^{*}(s_I \circ y).
  \]
\end{lemma}
\begin{proof}
  \uses{lem:gr_pullbackFreeIso_inv_freeMap}
  Unfolding the pullback action, \(s_I \circ (\psi^{*}y)\) is the coordinate inclusion
  \(s_I\) followed by \(Q_r^{-1}\) and \(\psi^{*}y\). The inverse free-pullback comparison
  intertwines the index inclusion (\cref{lem:gr_pullbackFreeIso_inv_freeMap}): sliding
  \(s_I\) across \(Q_r^{-1}\) turns it into \(Q_d^{-1}\) followed by \(\psi^{*}(s_I)\),
  whence \(s_I \circ (\psi^{*}y) = Q_d^{-1} \circ \psi^{*}(s_I \circ y)\).
\end{proof}

\begin{lemma}[The chart locus of a pullback contains the preimage of the chart locus]
  \label{lem:gr_chartLocus_rqPullback}
  \lean{AlgebraicGeometry.Grassmannian.chartLocus_rqPullback}
  \uses{def:chartLocus, lem:gr_chartComposite_rqPullback, lem:gr_isIso_pullback_map_comp,
    def:modules_pullbackComp}
  For \(\psi : T' \to T\) and a rank-\(d\) quotient \(y\) on \(T\), the \(\psi\)-preimage of
  the \(I\)-chart locus of \(y\) is contained in the \(I\)-chart locus of the pulled-back
  quotient:
  \[
    \psi^{-1}\bigl(\mathrm{chartLocus}(y, I)\bigr)
      \;\le\; \mathrm{chartLocus}(\psi^{*}y, I).
  \]
\end{lemma}
\begin{proof}
  \uses{lem:gr_chartComposite_rqPullback, lem:gr_isIso_pullback_map_comp,
    def:modules_pullbackComp}
  Let \(U = \psi^{-1}(\mathrm{chartLocus}(y, I))\). The inclusion \(U \hookrightarrow T'\)
  followed by \(\psi\) factors through the chart-locus inclusion
  \(\iota : \mathrm{chartLocus}(y, I) \hookrightarrow T\), over which the chart composite
  \(s_I \circ y\) pulls back to an isomorphism (\cref{def:chartLocus}). Hence its pullback
  along the factoring composite is an isomorphism
  (\cref{lem:gr_isIso_pullback_map_comp}), and the pullback pseudofunctor comparison
  (\cref{def:modules_pullbackComp}) together with the chart-composite bridge
  (\cref{lem:gr_chartComposite_rqPullback}) identifies it with the pullback to \(U\) of the
  chart composite of \(\psi^{*}y\). So \(s_I \circ (\psi^{*}y)\) is invertible over \(U\),
  i.e.\ \(U\) is one of the opens whose supremum defines
  \(\mathrm{chartLocus}(\psi^{*}y, I)\); taking the supremum gives
  \(U \le \mathrm{chartLocus}(\psi^{*}y, I)\).
\end{proof}

\begin{theorem}

[\(\mathrm{Gr}(r,d)\) represents the Grassmannian functor]
  \label{thm:grassmannian_universal_property}
  \lean{AlgebraicGeometry.Grassmannian.represents}
  \uses{def:grassmannian_functor, def:tautological_quotient,
    def:gr_universal_quotient_sheaf, def:gr_glued_scheme, def:gr_affine_chart,
    def:gr_universal_matrix, def:gr_transition, lem:gr_cocycle}
  % SOURCE: [Nitsure], "Construction of Hilbert and Quot Schemes", \S 1,
  % Exercise (2) "Grassmannian as a Quot scheme" (read from
  % references/nitsure-hilbert-quot-src/nitsure-hilbert-quot.tex, L546-L566).
  % SOURCE QUOTE:
  % "Show that $\grass(r,d)$ together with the quotient
  %  $u:\oplus^r\,{\mathcal O}_{\grass(r,d)} \to \U$ represents the
  %  contravariant functor
  %  $$\Grass(r,d) = \Quot_{\oplus^r\,{\mathcal O}_{\Z}/\Z/\Z}^{d, {\mathcal O}_{\Z}}$$
  %  from schemes to sets, which associates to any $T$ the
  %  set of all equivalence classes
  %  $\langle \F, q\rangle$ of quotients $q: \oplus^r\,{\mathcal O}_T \to \F$
  %  where $\F$ is a locally free sheaf on $T$ of rank $d$.
  %  Therefore, $\quot_{\oplus^r\,{\mathcal O}_{\Z}/\Z/\Z}^{d, {\mathcal O}_{\Z}}$
  %  exists, and equals $\grass(r,d)$."
  \textit{Source: [Nitsure], \S 1, Exercise (2).}
  The Grassmannian scheme \(\mathrm{Gr}(r,d)\) (\cref{def:gr_glued_scheme}),
  together with the tautological quotient \(\langle \mathcal{U}, u \rangle\)
  (\cref{def:tautological_quotient}), represents the functor
  \(\operatorname{Grass}(r,d)\) (\cref{def:grassmannian_functor}). That is, for
  every scheme \(T\) the natural map
  \[
    \operatorname{Hom}(T, \mathrm{Gr}(r,d))
      \;\xrightarrow{\;\sim\;}\;
    \operatorname{Grass}(r,d)(T),
    \qquad
    \bigl(\,\varphi : T \to \mathrm{Gr}(r,d)\,\bigr)
      \;\longmapsto\;
    \langle \varphi^{*}\mathcal{U},\, \varphi^{*}u \rangle,
  \]
  is a bijection, natural in \(T\). Consequently the Quot scheme
  \(\quot^{\,d,\,\mathcal{O}_{\mathbb{Z}}}_{\oplus^{r}\mathcal{O}_{\mathbb{Z}}/
  \mathbb{Z}/\mathbb{Z}}\) of rank-\(d\) quotients of the trivial rank-\(r\) sheaf
  exists and equals \(\mathrm{Gr}(r,d)\); the tautological quotient is the
  universal family.
\end{theorem}

\begin{proof}
  \uses{def:tautological_quotient, def:gr_universal_quotient_sheaf,
    def:gr_chart_quotient, def:gr_affine_chart, def:gr_universal_matrix,
    def:gr_transition, lem:gr_cocycle, def:gr_glued_scheme,
    def:grassmannian_functor, lem:tautologicalQuotient_epi,
    def:grPointOfRankQuotient, lem:chartLocus_isOpenCover, def:chartLocus,
    lem:isIso_pullback_isoLocus_map, def:gr_tautologicalRankQuotient,
    def:gr_chartMorphism, def:gr_universalQuotient_restrictionIso,
    lem:gr_comp_chartMorphism, lem:gr_chartComposite_rqPullback,
    lem:gr_chartLocus_rqPullback, lem:gr_chartMorphism_glue_compat}
  Naturality of \(\varphi \mapsto \langle \varphi^{*}\mathcal{U}, \varphi^{*}u
  \rangle\) is immediate: pullback is functorial and \(\varphi^{*}u\) is again a
  surjection onto a rank-\(d\) locally free sheaf, so the assignment defines a
  morphism of functors \(\operatorname{Hom}(-,\mathrm{Gr}(r,d)) \to
  \operatorname{Grass}(r,d)\). It remains to show it is a bijection for each
  \(T\); we exhibit an inverse, constructed Zariski-locally on \(T\) and glued.

  \emph{From a quotient to a morphism.} Let \(\langle \mathcal{F}, q \rangle \in
  \operatorname{Grass}(r,d)(T)\), so \(q : \mathcal{O}_T^{\,r} \twoheadrightarrow
  \mathcal{F}\) with \(\mathcal{F}\) locally free of rank \(d\). For a size-\(d\)
  subset \(I \subseteq \{1,\dots,r\}\) let \(T_I \subseteq T\) be the locus where
  the composite
  \[
    \mathcal{O}_T^{\,d}
      \overset{\ \iota_I\ }{\hookrightarrow} \mathcal{O}_T^{\,r}
      \xrightarrow{\ q\ } \mathcal{F},
  \]
  with \(\iota_I\) the inclusion of the \(I\)-indexed coordinates, is an
  isomorphism. This locus is open: in a local trivialisation of \(\mathcal{F}\)
  the composite is a \(d \times d\) matrix and \(T_I\) is the non-vanishing locus
  of its determinant. The \(T_I\) cover \(T\): at any point \(t\), the
  \(d\)-dimensional fibre quotient \(q(t) : \kappa(t)^{r} \twoheadrightarrow
  \mathcal{F}\otimes\kappa(t)\) is surjective, so some \(d\) of the \(r\) standard
  basis vectors map to a basis of the \(d\)-dimensional fibre, i.e.\ \(t \in T_I\)
  for that \(I\) (by Nakayama the isomorphism then holds in a neighbourhood).

  On \(T_I\), trivialise \(\mathcal{F}\) using \(\iota_I^{*}\): the quotient
  becomes a surjection \(q_I : \mathcal{O}_{T_I}^{\,r} \to \mathcal{O}_{T_I}^{\,d}\)
  represented by a \(d \times r\) matrix \(M^I\) whose \(I\)-th minor \(M^I_I\) is
  the identity \(1_{d \times d}\) -- exactly the shape of the universal matrix
  \(X^I\) (\cref{def:gr_universal_matrix}). The \(d(r-d)\) free entries of \(M^I\)
  are global functions on \(T_I\) and hence determine a unique
  \(\mathbb{Z}\)-algebra homomorphism \(R^I = \mathbb{Z}[X^I] \to
  \Gamma(T_I, \mathcal{O}_{T_I})\) sending \(x^I_{p,q}\) to the corresponding entry
  of \(M^I\), i.e.\ a unique morphism \(\varphi_I : T_I \to U^I\)
  (\cref{def:gr_affine_chart}) with \(\varphi_I^{*}X^I = M^I\) and thus
  \(\varphi_I^{*}u^I = q_I\) (\cref{def:gr_chart_quotient}).

  \emph{Gluing.} On the overlap \(T_I \cap T_J\) the two trivialisations differ by
  the change-of-basis \((M^I_J)^{-1}\), which is exactly the bundle transition
  \(g_{I,J} = (X^I_J)^{-1}\) of \(\mathcal{U}\)
  (\cref{def:gr_universal_quotient_sheaf}); concretely
  \(M^J = (M^I_J)^{-1} M^I = \varphi_I^{*}\theta_{I,J}(X^J)\)
  (\cref{def:gr_transition}, \cref{def:gr_universalMinorInv}). Hence
  \(\varphi_I\) and \(\varphi_J\) land in the overlap charts and agree there after
  the transition \(\theta_{I,J}\) -- this is the overlap compatibility
  \cref{lem:gr_chartMorphism_glue_compat} -- so by the cocycle condition
  (\cref{lem:gr_cocycle}) the \(\varphi_I\) glue to a single morphism
  \(\varphi = \mathrm{grPointOfRankQuotient}(\langle \mathcal{F}, q\rangle) :
  T \to \mathrm{Gr}(r,d)\) (\cref{def:grPointOfRankQuotient},
  \cref{def:gr_glued_scheme}). This morphism is the candidate inverse; that it
  is a two-sided inverse is the content of the two laws below.

  \emph{Existence (the law \(\mathrm{right\_inv}\)):
  \(\varphi^{*}\langle \mathcal{U}, u\rangle \sim \langle \mathcal{F}, q\rangle\).}
  The equivalence is checked over the chart cover \(\{T_I\}\). Over \(T_I\) the
  morphism \(\varphi\) restricts to \(\varphi_I\) followed by the chart immersion
  \(\iota_I\) (\cref{def:gr_chartMorphism}), and the restriction of the universal
  quotient sheaf along \(\iota_I\) is identified with the free rank-\(d\) sheaf by
  the chart restriction isomorphism
  \(\phi_I : \iota_I^{*}\mathcal{U} \xrightarrow{\sim} \mathcal{O}_{U^I}^{\,d}\)
  (\cref{def:gr_universalQuotient_restrictionIso}). Through this identification
  \(\varphi^{*}u\) becomes \(\varphi_I^{*}u^I\), which by construction of the chart
  morphism (\(\varphi_I^{*}X^I = M^I\)) equals \(q|_{T_I}\) in the
  \(I\)-trivialisation (\cref{def:gr_chartMorphism}, \cref{def:gr_chart_quotient}).
  The witnessing isomorphisms \(\varphi^{*}\mathcal{F} \cong \mathcal{F}\) so
  produced agree on overlaps -- the class is determined by \(\ker q\), which the
  trivialisations match -- and assemble to a single isomorphism over \(T\), giving
  \(\varphi^{*}\langle \mathcal{U}, u\rangle \sim \langle \mathcal{F}, q\rangle\).

  \emph{Uniqueness (the law \(\mathrm{left\_inv}\)):
  \(\mathrm{grPointOfRankQuotient}(\psi^{*}\langle \mathcal{U}, u\rangle) = \psi\)
  for every \(\psi : T \to \mathrm{Gr}(r,d)\).} It suffices to check the equality
  after restricting to each chart locus
  \(\mathrm{chartLocus}(\psi^{*}\langle \mathcal{U}, u\rangle, I)\) of the
  pulled-back tautological quotient. Over such a locus the composite of its
  inclusion with \(\psi\) factors through the \(\psi\)-preimage of the tautological
  chart locus, which is contained in that locus by
  \cref{lem:gr_chartLocus_rqPullback}; the chart-composite bridge
  \cref{lem:gr_chartComposite_rqPullback} identifies the chart data of
  \(\psi^{*}\langle \mathcal{U}, u\rangle\) with the \(\psi\)-pullback of the
  tautological chart data, so the matrix presenting it there is the \(\psi\)-image
  of the universal matrix \(X^I\). Hence the chart morphism classifying
  \(\psi^{*}\langle \mathcal{U}, u\rangle\) over this locus coincides with the
  restriction of \(\psi\) into the chart \(U^I\) (\cref{lem:gr_comp_chartMorphism}).
  As the chart loci cover, the glued morphism
  \(\mathrm{grPointOfRankQuotient}(\psi^{*}\langle \mathcal{U}, u\rangle)\) equals
  \(\psi\) by the uniqueness clause of the glue-data universal property.

  Existence and uniqueness exhibit \(\mathrm{grPointOfRankQuotient}\) as a two-sided
  inverse of \(\varphi \mapsto \langle \varphi^{*}\mathcal{U}, \varphi^{*}u\rangle\),
  so the latter is bijective.

  Therefore \(\mathrm{Gr}(r,d)\) with \(\langle \mathcal{U}, u \rangle\) is a fine
  moduli space: it represents \(\operatorname{Grass}(r,d)\), so
  \(\quot^{\,d,\,\mathcal{O}_{\mathbb{Z}}}_{\oplus^{r}\mathcal{O}_{\mathbb{Z}}/
  \mathbb{Z}/\mathbb{Z}}\) exists and equals \(\mathrm{Gr}(r,d)\).
\end{proof}
